%%%%%%%%%%%%%%%%%%%%%%%%%%%%%%%%%%%%%%%%%%%%%%%%%%%%%%%%%%%%%%%%%%%%%%%%%%%%%%
%
% Structural Semantics
% Constraint Geometry, Scope Calculus,
% and Resonant Models of Representation
%
% Consolidated working draft: Front Matter + Chapters 1-5
%
%%%%%%%%%%%%%%%%%%%%%%%%%%%%%%%%%%%%%%%%%%%%%%%%%%%%%%%%%%%%%%%%%%%%%%%%%%%%%%
\documentclass[11pt,twoside,openright]{book}

%% Engine
\usepackage{fontspec}
\usepackage{unicode-math}

%% Fonts
%% NOTE: Inconsolata is not installed in this build environment; substituted
%% with Latin Modern Mono, which is available and metrically compatible with
%% the rest of the TeX Gyre / Latin Modern family used here.
\setmainfont{TeX Gyre Pagella}
\setsansfont{TeX Gyre Heros}
\setmonofont{Latin Modern Mono}
\setmathfont{Latin Modern Math}

%% Mathematics
%% NOTE: amssymb/amsfonts removed. unicode-math already supplies the full
%% Unicode math symbol repertoire (including everything amssymb defines),
%% and loading amssymb alongside it under XeLaTeX raises fatal
%% "already defined" errors (\eth, \smallsetminus, \digamma, \backepsilon,
%% ...). amsmath is still loaded for its environments/macros, which do not
%% conflict.
\usepackage{amsmath}
\usepackage{amsthm}
\usepackage{mathtools}
\usepackage{mathrsfs}
\usepackage{bm}
%% stmaryrd removed for the same reason as amssymb above; none of its
%% distinctive symbols (e.g. \llbracket/\rrbracket) are currently used.

%% Typography
%% NOTE: expansion=true removed; font expansion is not supported by XeTeX,
%% only by pdfTeX, and raised a fatal error. protrusion works fine here.
\usepackage[protrusion=true,expansion=false]{microtype}

%% Layout
\usepackage{geometry}
\geometry{paper=a4paper,inner=35mm,outer=28mm,top=30mm,bottom=30mm,headsep=8mm,footskip=12mm}

%% Graphics
\usepackage{graphicx}
\usepackage{xcolor}

%% Tables
\usepackage{booktabs}
\usepackage{array}
\usepackage{longtable}
\usepackage{multirow}

%% Lists (used sparingly; most enumerations below have been converted to prose)
\usepackage{enumitem}

%% Hyperlinks
\usepackage[colorlinks=true,linkcolor=blue!50!black,citecolor=blue!50!black,urlcolor=blue!50!black]{hyperref}
\usepackage[nameinlink]{cleveref}

%% Bibliography
%% NOTE: biblatex/biber is not installed in this build environment, so this
%% consolidated draft uses natbib/bibtex instead, which are available. Swap
%% back to biblatex if/when biber is available in the target build chain.
\usepackage[numbers,sort&compress]{natbib}

%% Index
\usepackage{imakeidx}
\makeindex
%% Separate index of symbols, populated in Appendix H (Notation and
%% Mathematical Conventions), which is where every recurring symbol is
%% already collected in one place — indexing it there rather than at each
%% of its many usage sites keeps one authoritative page reference per
%% symbol instead of a page list diluted by every occurrence.
\makeindex[name=symbols,title=Index of Symbols]

%% Page style
\usepackage{fancyhdr}
\pagestyle{fancy}
\fancyhf{}
\fancyhead[LE]{\thepage}
\fancyhead[RO]{\thepage}
\fancyhead[RE]{\leftmark}
\fancyhead[LO]{\rightmark}

%% Equation numbering
\numberwithin{equation}{chapter}

%%%%%%%%%%%%%%%%%%%%%%%%%%%%%%%%%%%%%%%%%%%%%%%%%%%%%%%%%%%%%%%%%%%%%%%%%%%%%%
%% Theorem environments
%%%%%%%%%%%%%%%%%%%%%%%%%%%%%%%%%%%%%%%%%%%%%%%%%%%%%%%%%%%%%%%%%%%%%%%%%%%%%%
\theoremstyle{definition}
\newtheorem{definition}{Definition}[chapter]
\newtheorem{example}[definition]{Example}
\newtheorem{construction}[definition]{Construction}
\theoremstyle{plain}
\newtheorem{theorem}[definition]{Theorem}
\newtheorem{lemma}[definition]{Lemma}
\newtheorem{proposition}[definition]{Proposition}
\newtheorem{corollary}[definition]{Corollary}
\theoremstyle{remark}
\newtheorem{remark}[definition]{Remark}
\newtheorem{observation}[definition]{Observation}
%% Empirical hypotheses are typeset in their own style, distinct from
%% mathematical theorems/propositions, to keep the mathematics/empirical-claim
%% distinction (Chapter 1's Methodological Statement) visible at the level of
%% typography, not just prose.
\theoremstyle{definition}
\newtheorem{hypothesis}[definition]{Hypothesis}
%% Conjectures are explicitly unproved and are kept typographically distinct
%% from theorems/propositions for the same reason.
\newtheorem{conjecture}[definition]{Conjecture}
%% Starred/unnumbered variant, used in the proofs appendix to restate a
%% theorem (already numbered in the main text) before giving its proof.
\newtheorem*{theorem*}{Theorem}

%%%%%%%%%%%%%%%%%%%%%%%%%%%%%%%%%%%%%%%%%%%%%%%%%%%%%%%%%%%%%%%%%%%%%%%%%%%%%%
%% Global mathematical notation
%%
%% CONSOLIDATED: this registry now covers both the Scope Calculus symbols
%% and the Whitehead-reconstruction symbols introduced in Ch.4-5, which
%% were previously undocumented here. \Tau is corrected from a single
%% function to a set of admissible operators, matching how it is used
%% throughout (tau_i \in \Tau).
%%%%%%%%%%%%%%%%%%%%%%%%%%%%%%%%%%%%%%%%%%%%%%%%%%%%%%%%%%%%%%%%%%%%%%%%%%%%%%
% Categories
\newcommand{\CatScope}{\mathbf{Scope}}
\newcommand{\CatProcess}{\mathbf{Process}}
% Objects
\newcommand{\Scope}{\mathbf{S}}
% Morphisms
\newcommand{\ScopeMorph}{\sigma}
% Primitive alphabet
\newcommand{\Alphabet}{\Sigma}
\newcommand{\Pop}{\mathcal{P}}
\newcommand{\Refuse}{\mathcal{R}}
\newcommand{\Bind}{\mathcal{B}_d}
\newcommand{\Collapse}{\mathcal{C}_l}
\newcommand{\ScopeAlphabet}{\Alphabet=\{\Pop,\Refuse,\Bind,\Collapse\}}
% Scope components
\newcommand{\Boundary}{\mathcal{B}}
\newcommand{\Interior}{\mathcal{E}}
\newcommand{\Children}{\mathcal{S}_{\mathrm{child}}}
% Computational domains
\newcommand{\Domain}{\mathcal{D}}
\newcommand{\System}{\mathcal{C}}
% State spaces
\newcommand{\State}{\mathcal{X}}
% Operators
\DeclareMathOperator{\Hom}{Hom}
\DeclareMathOperator{\Obj}{Ob}
\DeclareMathOperator{\Ker}{Ker}
\DeclareMathOperator{\Image}{Im}
\DeclareMathOperator{\PLV}{PLV}
% Admissibility
\newcommand{\Adm}{\mathsf{Adm}}
% Transition set (a computational system's admissible operators; corrected
% from the earlier single-function reading)
\newcommand{\Trans}{\mathcal{T}}
% MEM|8
\newcommand{\Hilbert}{\mathcal{H}}
%% Semantic double brackets, used for the categorical interpretation of
%% types in Appendix C. stmaryrd is unavailable in this build (see the
%% Mathematics section above), so these are defined directly.
\newcommand{\llbracket}{[\mkern-3.5mu[}
\newcommand{\rrbracket}{]\mkern-3.5mu]}
% Whitehead reconstruction (Ch.4-5): antecedent data, prehension set,
% valuation polarity, satisfaction state, and Subjective Form. These are
% left as plain italic letters, consistent with how they were introduced,
% but are registered here so the notation stays traceable across chapters.
\newcommand{\WData}{D}
\newcommand{\WPrehensions}{P}
\newcommand{\WValuation}{v}
\newcommand{\WSatisfaction}{S}
\newcommand{\WSubjForm}{A}
\newcommand{\WOccasion}{E}

\title{Structural Semantics}
\author{Flyxion}
\date{\today}

\begin{document}

%%%%%%%%%%%%%%%%%%%%%%%%%%%%%%%%%%%%%%%%%%%%%%%%%%%%%%%%%%%%%%%%%%%%%%%%%%%%%%
%% Front Matter
%%%%%%%%%%%%%%%%%%%%%%%%%%%%%%%%%%%%%%%%%%%%%%%%%%%%%%%%%%%%%%%%%%%%%%%%%%%%%%
\frontmatter

%% ---------------------------------------------------------------- Title page
\begin{titlepage}
\centering
\vspace*{2cm}
{\Huge\bfseries Structural Semantics\par}
\vspace{1cm}
{\Large Constraint Geometry, Scope Calculus,\\ and Resonant Models of Representation\par}
\vspace{2cm}
{\Large Flyxion\par}
\vfill
A Research Monograph
\vspace{0.5cm}

Version 1.0
\vspace{2cm}

{\large \today}
\end{titlepage}
\cleardoublepage

%% ------------------------------------------------------------- Copyright page
\thispagestyle{empty}
\vspace*{\fill}
\begin{center}
Copyright \copyright\ \the\year\ Flyxion

All rights reserved. No part of this publication may be reproduced, stored
in a retrieval system, or transmitted in any form without permission,
except for brief quotations used in scholarly review.

This work develops original mathematical constructions, formal definitions,
and computational models. Where established results are discussed,
appropriate references are provided.
\end{center}
\vspace*{\fill}
\cleardoublepage

%% --------------------------------------------------------------- Epigraph
\thispagestyle{empty}
\vspace*{0.25\textheight}
\begin{flushright}
\itshape
``The aim of science is not things themselves, as the dogmatists in their
simplicity imagine, but the relations among things.''
\vspace{1em}

--- Henri Poincar\'e
\end{flushright}
\cleardoublepage

\tableofcontents
\cleardoublepage

%% --------------------------------------------------------------- Preface
\chapter*{Preface}
\addcontentsline{toc}{chapter}{Preface}

The rapid development of modern machine learning has revived questions
that have occupied philosophy, mathematics, linguistics, and cognitive
science for more than a century. Systems capable of producing coherent
mathematical arguments, computer programs, scientific summaries, and
extended dialogue have blurred traditional distinctions between syntax,
semantics, reasoning, and understanding. At the same time, the conceptual
vocabulary used to discuss these systems has become increasingly
heterogeneous: terms such as \emph{understanding}, \emph{representation},
\emph{meaning}, \emph{reasoning}, and \emph{intelligence} are often
employed without a common mathematical framework, leading to disagreements
that arise as much from differences in definition as from differences in
evidence.

The purpose of this monograph is to construct such a framework. Rather
than beginning with phenomenological assumptions, behavioral benchmarks,
or biological commitments, we begin with the mathematics of constrained
transformation. The central question investigated throughout this work is
whether semantic capacity can be characterized as a structural property of
admissible transformations over bounded contexts.

The resulting framework, called Structural Semantics, combines ideas from
topology, algebra, category theory, dynamical systems, signal processing,
and machine learning, and engages process philosophy directly through a
formal correspondence with Whitehead's ontology. The mathematical
developments are separated carefully from their philosophical
interpretations, allowing each to be evaluated on its own terms. Where a
result depends on mathematics beyond what a given chapter needs to state
it, the corresponding construction --- from the $\lambda$-calculus and
abstract rewriting theory through categorical, topological,
sheaf-theoretic, and measure-theoretic foundations --- is developed
independently in the appendices rather than assumed. The work is intended
equally for mathematicians, computer scientists, philosophers of mind,
and researchers in mechanistic interpretability. No commitment to any
particular metaphysical position is assumed.

\paragraph{Relation to earlier structuralism.} This monograph does not
replace the structural semantics developed within twentieth-century
linguistics. Rather, it generalizes the structuralist methodology from
systems of linguistic signs to systems of admissible computational
transformations. The resulting framework shares with classical
structuralism an emphasis on relational organization while introducing a
different mathematical foundation based upon topology, algebra, category
theory, and dynamical systems. The precise relationship between the two
uses of the term is stated in \cref{sec:terminological-note} and developed
at length in \cref{sec:structuralist-tradition}.

%% --------------------------------------------------------- Methodological Statement
\chapter*{Methodological Statement}
\addcontentsline{toc}{chapter}{Methodological Statement}

Throughout this monograph, three distinct classes of claims are carefully
distinguished, and every chapter closes with a status summary that sorts
its content among them.

\emph{Mathematical claims} consist of definitions, lemmas, theorems, and
proofs concerning topology, scope algebra, category theory, and related
formal structures; these are established by proof from stated hypotheses
and by nothing else. \emph{Computational claims} concern algorithms,
diagnostic procedures, and models of transformer representations; these
are empirical hypotheses whose validity depends upon experimental
evaluation, and they are presented as such rather than as consequences of
the mathematics. \emph{Philosophical claims} concern interpretation: they
discuss possible implications of the mathematical and computational
developments but are not presented as logical consequences of the formal
theory.

Maintaining this distinction is a central methodological principle of the
present work. A theorem in Volume~II never licenses a claim about
consciousness or moral status; a diagnostic algorithm in Volume~III is
never treated as proven merely because it is well motivated; and a
philosophical reading offered in Volume~IV is never mistaken for a
mathematical consequence of what precedes it.

%% ------------------------------------------------------- How to Read This Monograph
\chapter*{How to Read This Monograph}
\addcontentsline{toc}{chapter}{How to Read This Monograph}

Although the eighteen chapters are arranged sequentially, the book is
organized into four relatively independent volumes, and readers are
encouraged to enter at the volume suited to their interests rather than
proceed straight through. Readers primarily interested in philosophical
questions may begin with Volume~I, which establishes the epistemic
framing, historical lineage, and the boundary separating structural
claims from phenomenological ones. Readers interested in the mathematical
framework itself should concentrate on Volume~II, where bounded scope
geometries, the primitive algebra, the category of scopes, and the
representation-invariance and process-geometry results are developed; the
notation introduced in \cref{sec:notation-and-conventions} is sufficient
preparation, and the historical material is not a prerequisite. Readers
interested in transformer analysis and computational diagnostics may read
Volume~III after becoming familiar with the basic scope formalism,
since the diagnostics there are applications of, rather than
alternatives to, the mathematics of Volume~II. Volume~IV situates the
resulting framework within the broader landscape of philosophy, cognitive
science, and applied mathematics, and states its limitations and
non-claims explicitly. The appendices provide complete proofs, notation
tables, implementation details, and supplementary derivations referenced
throughout the text, and are consulted as needed rather than read in
sequence.

%% ------------------------------------------------------- Notation and Conventions
\chapter*{Notation and Conventions}
\label{sec:notation-and-conventions}
\addcontentsline{toc}{chapter}{Notation and Conventions}

Bold uppercase symbols denote categories or structured spaces, calligraphic
symbols denote operators, constraint systems, or families of sets, and
lowercase Greek letters denote morphisms, transformations, or parameters.
Once a symbol is assigned a formal meaning it is preserved throughout the
monograph and is not reassigned in a later chapter; \cref{app:notation}
collects the complete registry.

Unless stated otherwise, $\Scope=\langle\Boundary,\Interior,\Children\rangle$
denotes a bounded scope geometry, $\ScopeAlphabet$ denotes the primitive
event alphabet, and $\ScopeMorph:\Scope_i\to\Scope_j$ denotes an admissible
scope morphism. These three lines anticipate Chapter~6 and are given here
only so that forward references elsewhere in the front matter are legible;
their full definitions, and the argument for why the alphabet has exactly
four elements, are deferred to Part~II.

The reconstruction of Whitehead's process ontology in Chapters~4 and~5
uses a second, disjoint set of symbols native to that literature: $\WOccasion$
for an actual occasion, $\WData$ for antecedent data, $\WPrehensions$ for a
prehension, $\WValuation$ for valuation polarity, $\WSatisfaction$ for a
satisfaction state, and $\WSubjForm$ for Subjective Form. These are kept
visually distinct from the Scope Calculus notation deliberately: no
identification between the two systems is assumed before Chapter~10, where
the translation functor is constructed and any correspondence between them
is established as a proposition rather than presupposed by shared notation.

Finally, a computational system is written $M=(\Domain,\Trans)$, where
$\Domain$ is a state space and $\Trans$ is the \emph{set} of admissible
transition operators acting on $\Domain$, each individual operator being
written $\tau\in\Trans$ with $\tau:\Domain\to\Domain$. This corrects an
earlier draft in which $\Trans$ was written as if it were itself a single
function; the set reading is what the rest of the text relies on whenever
it writes $\tau_i\in\Trans$.

%%%%%%%%%%%%%%%%%%%%%%%%%%%%%%%%%%%%%%%%%%%%%%%%%%%%%%%%%%%%%%%%%%%%%%%%%%%%%%
%% Main matter
%%%%%%%%%%%%%%%%%%%%%%%%%%%%%%%%%%%%%%%%%%%%%%%%%%%%%%%%%%%%%%%%%%%%%%%%%%%%%%
\mainmatter

%%%%%%%%%%%%%%%%%%%%%%%%%%%%%%%%%%%%%%%%%%%%%%%%%%%%%%%%%%%%%%%%%%%%%%%%%%%%%%
\part{Epistemic and Historical Boundaries}
%%%%%%%%%%%%%%%%%%%%%%%%%%%%%%%%%%%%%%%%%%%%%%%%%%%%%%%%%%%%%%%%%%%%%%%%%%%%%%

%%%%%%%%%%%%%%%%%%%%%%%%%%%%%%%%%%%%%%%%%%%%%%%%%%%%%%%%%%%%%%%%%%%%%%%%%%%%%%
\chapter{Epistemic Framing, Disentanglement, and Primary Research Questions}
\label{ch:epistemic-framing}
%%%%%%%%%%%%%%%%%%%%%%%%%%%%%%%%%%%%%%%%%%%%%%%%%%%%%%%%%%%%%%%%%%%%%%%%%%%%%%

\begin{center}
\fbox{\parbox{0.92\textwidth}{\textbf{Chapter Roadmap.} This chapter
establishes the epistemic foundations of the monograph. It defines the
central object of study, separates structural semantic capacity from
neighboring philosophical concepts, introduces the primary research
questions, summarizes the principal contributions, and outlines the
methodological strategy that governs all subsequent chapters.}}
\end{center}

\vspace{1em}
\begin{flushright}\itshape
The limits of my language mean the limits of my world.\\
\upshape --- Ludwig Wittgenstein, \emph{Tractatus Logico-Philosophicus}
\end{flushright}
\vspace{1em}

\section{Introduction}

Few concepts have become simultaneously more important and more ambiguous
than the notion of \emph{understanding}. Within philosophy, cognitive
science, linguistics, neuroscience, and artificial intelligence, the same
word often refers to fundamentally different phenomena: subjective
conscious awareness in one context, successful behavioral performance in
another, internal representation, semantic competence, or predictive
capability in others. Debates regarding machine intelligence frequently
proceed without agreement concerning the object under discussion itself.

Recent advances in transformer-based language models have intensified
this ambiguity. Systems capable of producing mathematically coherent
proofs, generating software, translating between natural languages,
summarizing scientific literature, and participating in extended
technical dialogue have challenged long-standing intuitions concerning the
relationship between computation and semantic understanding
\cite{Vaswani2017,Brown2020,OpenAI2023}. At the same time, philosophical
objections remain influential, including the Chinese Room \cite{Searle1980},
the distinction between phenomenal and access consciousness
\cite{Block1995}, and various forms of embodied cognition questioning
whether statistical computation alone can constitute genuine
understanding.

The present work does not attempt to resolve the metaphysical question of
consciousness. Instead it addresses a more precise question: is there a
mathematically definable notion of semantic capacity that depends only
upon the preservation of structural relationships during computation? If
such a notion exists, it should admit formal definitions, rigorous proofs,
computational implementations, and empirical evaluation independently of
debates concerning subjective experience. We refer to this proposed
framework throughout as \emph{Structural Semantics}. The central
hypothesis is that semantic capacity may be understood as the preservation
of admissible structural relationships across transformations of bounded
contexts; later chapters develop this proposal using topology, rewriting
systems, category theory, dynamical systems, and signal processing.

\begin{remark}[Terminological note]
\label{sec:terminological-note}
The phrase \emph{structural semantics} has an established meaning in
linguistics, referring to structuralist theories of lexical meaning
developed by Saussure, Hjelmslev, Lyons, Greimas, Coseriu, and related
authors. Throughout this monograph the term is used in a different,
explicitly defined sense: a mathematical framework for the preservation of
admissible structural transformations over bounded contexts. The present
theory is inspired by the general structuralist tradition but is not a
development of lexical field theory, componential analysis, or structural
linguistics; \cref{sec:structuralist-tradition} situates the two usages
against one another in detail.
\end{remark}

\section{Motivating the Problem}

The history of theories of meaning exhibits a recurring pattern:
philosophical accounts generally begin from one of three starting points.
The first begins from conscious experience, regarding understanding as
inseparable from phenomenal awareness or intentionality, as in
phenomenology, process philosophy, and much contemporary discussion of
consciousness. The second begins from observable behavior, as in
functionalist approaches that characterize understanding in terms of
successful interaction, prediction, or problem solving, treating internal
implementation as secondary to externally measurable performance. The
third begins from representation itself, asking whether internal
structures preserve relationships sufficiently to support inference,
abstraction, and generalization independently of their physical
realization.

Each perspective captures an important aspect of cognition, but they
address different questions, and difficulties arise when conclusions
appropriate to one framework are transferred without modification into
another: a behavioral criterion cannot by itself establish phenomenology,
and the absence of phenomenological evidence does not immediately imply
the absence of mathematically characterizable semantic structure. This
observation motivates the methodological separation adopted here. Rather
than collapsing the three viewpoints into a single definition, we
introduce a formal framework that treats structural semantic capacity as
an independent object of investigation, deliberately separated from the
philosophical questions concerning consciousness that remain important
but lie outside its scope.

\section{Epistemic Separation}

The first objective of the monograph is therefore methodological rather
than mathematical. Before introducing scope geometry or category theory,
it is necessary to distinguish several concepts that are frequently
conflated in the literature: phenomenal consciousness, behavioral
competence, veridical truth, intentional agency, and structural semantic
capacity. Only the last of these constitutes the primary object of
mathematical investigation developed in subsequent chapters. The
remaining concepts may interact with structural semantic capacity, but
they are not identified with it; results proved within the Scope Calculus
should not automatically be interpreted as results concerning
consciousness, moral agency, or epistemic truth, which belong to
different theoretical domains and require independent justification.

This separation serves two purposes. First, it permits the mathematical
development to proceed independently of unresolved philosophical debates.
Second, it makes empirical evaluation possible, since structural
properties may be investigated directly through computational observation
without requiring access to subjective experience.

\section{Formal Foundations}

Having established the methodological scope of the monograph, we now
introduce the mathematical objects on which the remainder of the
development is built, identifying only the minimal structures required to
state the principal research questions precisely. Throughout the text a
\emph{system} is understood abstractly, with no assumption regarding its
physical realization: it may be implemented biologically, digitally,
symbolically, mechanically, or by any other computational substrate, and
the definitions below are correspondingly representation-independent.

\begin{definition}[Computational system]
A computational system is an ordered pair $M=(\Domain,\Trans)$, where
$\Domain$ is a non-empty state space and $\Trans$ is a set of admissible
state-transition operators $\tau:\Domain\to\Domain$ acting on $\Domain$.
No assumptions are made regarding continuity, discreteness, determinism,
or stochasticity unless explicitly stated.
\end{definition}

\begin{remark}
The abstraction deliberately encompasses finite-state automata, symbolic
rewriting systems, neural networks, probabilistic graphical models,
dynamical systems, and biological cognitive architectures. The
mathematical framework developed throughout this monograph applies at
this level of generality whenever the stated hypotheses are satisfied.
\end{remark}

\subsection{Representations}

The notion of representation is equally general. Rather than committing
to a particular implementation, we define a representation only through
its structural role.

\begin{definition}[Representation]
Let $M=(\Domain,\Trans)$ be a computational system. A representation is
an element $x\in\Domain$ together with the admissible transition
structure induced by $\Trans$. Two representations need not possess
identical internal encodings in order to support identical structural
behavior.
\end{definition}

This definition intentionally avoids identifying representations with
vectors, symbols, neurons, graphs, or linguistic expressions; a
representation refers only to the mathematical state on which the
transition operators act.

\subsection{Transformation}

Understanding, regardless of philosophical commitment, concerns not
isolated states but the evolution of states under transformation, so
transformations rather than static representations become the primary
object of investigation.

\begin{definition}[Transformation sequence]
Let $M=(\Domain,\Trans)$ be a computational system. A transformation
sequence is a finite ordered family $T=(x_0,x_1,\ldots,x_n)$ such that
$x_{i+1}=\tau_i(x_i)$ for some $\tau_i\in\Trans$. The sequence $T$ records
the structural evolution of a representation through the state space.
\end{definition}

\subsection{Motivation for Scope Geometry}

Transformation sequences alone are insufficient for characterizing
semantic preservation, since two systems may undergo identical numbers of
state transitions while preserving fundamentally different contextual
relationships. A symbolic proof assistant may reject an invalid inference
by enforcing logical typing constraints, whereas a neural language model
may continue generating fluent text despite violating assumptions
introduced earlier in the reasoning process; both systems perform
transformations, yet only one explicitly preserves the contextual
constraints governing the computation.

The central mathematical idea developed in this monograph is that such
constraints may themselves be treated as geometric objects. Rather than
asking whether a representation ``contains meaning,'' we ask whether
successive transformations preserve the admissible boundaries within
which meaning is defined. This observation motivates \emph{bounded scope
geometries}, developed fully in Chapters~6 through~10; for the present
chapter it suffices to observe that semantic preservation appears to
depend not upon individual states but upon the structural integrity of
the permissible transformations connecting them.

\section{The Disentanglement Operator}
\label{sec:disentanglement}

The preceding discussion suggests that much of the disagreement
surrounding semantic capacity originates not from conflicting mathematical
results but from differences in the domains over which the term
\emph{understanding} is defined. Before constructing an algebra of
semantic transformations, it is therefore necessary to specify the
epistemic domain within which subsequent definitions operate. The purpose
of the Disentanglement Operator is precisely this: rather than measuring
semantic capacity, it separates semantic capacity from other properties
that may coexist with it but are not logically implied by it, allowing
later chapters to formulate structural results without simultaneously
making claims concerning consciousness, truth, or agency.

\begin{definition}[Disentanglement operator]
\label{def:disentanglement}
Let $\mathcal{P}=\{\Phi,\mathbf{T},\mathbf{A},\mathbf{F},\mathbf{S}\}$
denote the collection of properties commonly associated with
understanding, where $\Phi$ denotes phenomenal consciousness, $\mathbf{T}$
denotes veridical truth, $\mathbf{A}$ denotes intentional agency,
$\mathbf{F}$ denotes syntactic fluency, and $\mathbf{S}$ denotes
structural semantic capacity. The \emph{Disentanglement Operator}
$\mathfrak{D}:\mathcal{P}\to 2^{\mathcal{P}}$ identifies structural
semantic capacity as an independent object of mathematical investigation
by separating $\mathbf{S}$ from the remaining properties: throughout this
monograph, $\mathbf{S}\not\equiv\Phi$, $\mathbf{S}\not\equiv\mathbf{T}$,
$\mathbf{S}\not\equiv\mathbf{A}$, and $\mathbf{S}\not\equiv\mathbf{F}$.
\end{definition}

\begin{remark}
\Cref{def:disentanglement} establishes only logical independence; it does
not imply statistical, causal, or empirical independence. Systems
possessing high structural semantic capacity may also exhibit high
behavioral competence, and biological systems may simultaneously possess
structural semantic capacity and phenomenal experience. The present
framework merely asserts that these properties are mathematically
distinguishable and should not be identified without additional argument.
\end{remark}

\subsection{Interpretation}

The distinction above serves as a recurring methodological principle
throughout the monograph. Structural Semantics is not presented as a
theory of consciousness: whether subjective awareness accompanies a
particular computational architecture remains outside the scope of the
present formalism and is revisited only in the philosophical discussion
of Chapters~16 and~17. Nor is it a theory of objective truth: a system may
preserve contextual relationships perfectly while reasoning within an
entirely fictional, hypothetical, or mathematically inconsistent universe,
since the framework evaluates the preservation of internal constraint
structures rather than correspondence with external reality. Nor is it a
theory of agency: goal-directed behavior, preference formation, intentional
action, and moral status require additional theoretical machinery that is
deliberately excluded from the present development. Finally, Structural
Semantics is not equivalent to linguistic fluency: contemporary language
models routinely generate syntactically well-formed sentences while
simultaneously violating assumptions established only a few paragraphs
earlier, a discrepancy that motivates the distinction between surface
grammatical regularity and the preservation of admissible contextual
structure.

\begin{example}[Structural coherence without truth]
Consider a narrative set entirely within a fictional world governed by
its own internal rules --- for instance, that a particular journey
between two named locations cannot be completed within a stated span of
time. A system that continues to respect this constraint throughout an
extended narrative, refusing continuations that would violate it, is
exhibiting structural semantic capacity with respect to that narrative's
admissibility conditions, regardless of the fact that neither the
locations nor the constraint correspond to anything in physical reality.
Conversely, a system that abruptly violates a constraint it had itself
established earlier in the same narrative --- introducing a character
incompatible with rules already fixed for that fictional world --- has
produced a boundary violation in the sense of Chapter~6, independently
of whether either version of the narrative is judged more ``realistic.''
Structural Semantics evaluates the second kind of failure. It has
nothing to say about the first.
\end{example}

\section{Boundary Constraints and Bounded Contexts}

The Disentanglement Operator isolates the domain of investigation but does
not yet specify the mathematical structures through which semantic
preservation is expressed. We therefore introduce the notion of a
\emph{boundary constraint}: informally, the conditions under which a
state may enter, remain within, or leave a contextual region. Such
constraints appear in type systems, formal grammars, topological
boundaries, logical theories, and dynamical systems, and the present
formulation deliberately abstracts over these examples, treating a
boundary not as a physical surface but as a collection of admissibility
conditions governing transformations of representations.

\begin{definition}[Boundary constraint]
Let $M=(\Domain,\Trans)$ be a computational system. A boundary constraint
is a predicate $\Boundary:\Domain\times\Trans\to\{0,1\}$, where
$\Boundary(x,\tau)=1$ indicates that the transition $x\xrightarrow{\tau}\tau(x)$
is admissible, and $\Boundary(x,\tau)=0$ indicates that it violates the
structural constraints governing the current context.
\end{definition}

Definition~2.7 intentionally avoids assigning any particular computational
interpretation to $\Boundary$: in different applications it may represent
typing constraints, logical consistency conditions, physical conservation
laws, permission systems, grammatical restrictions, or learned
admissibility relations extracted from data. The subsequent mathematics
depends only on the existence of such a predicate and not on its
implementation.

\subsection{Bounded contexts}

Semantic reasoning rarely occurs within an unrestricted universe of
states. Mathematical proofs introduce hypotheses, computer programs
create lexical environments, conversations establish topics, and
scientific models operate under explicitly stated assumptions; in each
case reasoning proceeds within a bounded context whose admissibility
conditions determine which transformations remain valid.

\begin{definition}[Bounded context]
A bounded context is a pair $(\Domain,\Boundary)$ consisting of a state
domain together with a boundary constraint predicate. Its admissible
region is
\[
\Adm(\Domain,\Boundary)=\{x\in\Domain \mid \Boundary(x,\tau)=1 \text{ for every admissible } \tau\}.
\]
\end{definition}

The notation $\Adm(\Domain,\Boundary)$ recurs throughout the monograph,
denoting informally the collection of states that satisfy every
structural constraint currently imposed by the context.

\subsection{Motivation for scope geometry}

Bounded contexts remain insufficient for describing nested reasoning
processes. Modern computational systems routinely manipulate multiple
interacting contexts simultaneously: a theorem may invoke several nested
lemmas, a compiler maintains independent lexical environments, transformer
models preserve long-range dependencies across dozens of processing
layers, and human reasoning shifts continuously between assumptions,
sub-problems, hypothetical scenarios, and meta-level analyses. These
examples suggest that semantic preservation cannot be understood solely
through individual admissible regions, but requires a hierarchy of
interacting contexts whose relationships are themselves subject to
structural constraints. This motivates \emph{scope geometry}: rather than
viewing context as a single admissible set, it models reasoning as a
collection of nested bounded regions connected by explicit
admissibility-preserving transformations.

\begin{definition}[Scope geometry, preliminary]
A scope geometry is a bounded context together with an explicitly
specified collection of subordinate bounded contexts and the structural
relations connecting them.
\end{definition}

This definition serves only as a conceptual preview; the complete formal
treatment, including the algebra of scope transformations, admissibility
morphisms, and categorical structure, is developed systematically in
Chapters~6 through~10. In particular, the full tuple
$\Scope=\langle\Boundary,\Interior,\Children\rangle$ is deliberately
withheld until Chapter~6, so that it appears there as the culmination of
the intervening development rather than as a symbol introduced before its
mathematical necessity.

\subsection{A preview of Structural Semantics}

The progression established thus far is deliberately incremental:
beginning with computational systems, we introduced representations,
transformation sequences, boundary constraints, bounded contexts, and
finally scope geometries, each concept extending the previous one while
remaining independent of any particular computational architecture. The
central hypothesis explored throughout the remainder of this monograph
may now be stated informally: semantic capacity is not identified with
isolated representations, but is associated with the preservation of
admissible structure across transformations acting on hierarchies of
bounded contexts.

\section{The Research Program}

The preceding sections established the conceptual vocabulary on which the
remainder of this monograph is built, but no substantial mathematical
claims have yet been made; we have established only the domain within
which such claims may be formulated. The objective of the remainder of
this work is to determine whether semantic capacity can be characterized
as a structural property of admissible transformations over bounded
contexts, and this objective separates into four progressively more
specialized research questions, concerning conceptual foundations,
mathematical structure, computational implementation, and empirical
verification respectively.

The first question concerns the object of study itself: can semantic
capacity be defined independently of phenomenal consciousness, intentional
agency, and behavioral performance? The purpose is not to deny the
importance of consciousness but to ask whether semantic structure admits
an independent mathematical description; an affirmative answer would
imply that semantic capacity constitutes a legitimate object of
mathematical investigation irrespective of one's preferred philosophical
theory of mind. Volumes~I and~II develop the formal foundations required
to address this question.

The second question concerns mathematical structure: which mathematical
structures are necessary and sufficient to describe semantic preservation
under transformation? Several candidate frameworks already exist,
including topology, graph theory, rewriting systems, information
geometry, category theory, and dynamical systems; rather than replacing
these disciplines, the present work attempts to synthesize selected ideas
from each into a unified theory of bounded semantic transformations,
developed systematically in Chapters~6 through~10.

The third question concerns implementation: can structural semantic
preservation be extracted from the internal representations of
computational systems? Recent advances in mechanistic interpretability
have demonstrated that large neural networks admit increasingly
sophisticated internal analysis through activation trajectories, sparse
feature models, residual stream decompositions, and circuit-level
investigations \cite{Olah2020,Elhage2021,Nanda2023}, suggesting that
semantic structure may itself become an observable computational quantity
rather than merely a philosophical abstraction. Volume~III develops one
possible mathematical framework for such measurements.

Finally, every mathematical model intended to describe computational
systems should generate consequences capable of empirical evaluation: can
structural measures of semantic preservation generate falsifiable
predictions distinguishing genuine contextual coherence from superficial
statistical fluency? The monograph therefore concludes its technical
development by constructing explicit diagnostic algorithms together with
experimental protocols designed to evaluate the proposed framework, and
these protocols are intentionally formulated so that negative results are
as informative as positive ones; a framework that cannot in principle be
falsified falls outside the scope of the present methodology.

\section{Statement of Contributions}

The mathematical developments presented throughout this monograph are
organized around five principal contributions. The first is a formal
Scope Calculus governing admissible transformations over bounded
contexts, generated minimally by four primitive operations whose
algebraic properties are developed in Chapters~7 and~9. The second is a
structural definition of semantic capacity that depends only on
admissibility-preserving transformations rather than on any particular
computational substrate, permitting semantic equivalence to be
investigated independently of implementation. The third is a categorical
formulation of scope transformations together with a translation functor
relating elements of Whiteheadian process philosophy to the Scope
Calculus, developed in Chapter~10 as a comparative rather than reductive
construction. The fourth is the formulation of the MEM$|$8 analytical
framework for studying activation trajectories within transformer
architectures, which treats transformed activation trajectories as
objects of signal analysis using established tools from harmonic analysis
and digital signal processing rather than proposing a physical wave
ontology. The fifth is an empirical diagnostic program designed to
evaluate the proposed mathematical structures, with explicit falsification
criteria, statistical controls, computational complexity analyses, and
implementation protocols developed in Chapters~14 and~15 and expanded in
the appendices.

\section{Methodological Roadmap}

The remainder of this monograph is organized into four interconnected
volumes. Volume~I examines the historical and philosophical development
of the concept of understanding, culminating in a precise characterization
of the phenomenological boundary separating structural questions from
questions concerning subjective experience. Volume~II develops the
mathematical foundations of Structural Semantics, beginning with bounded
scope geometries and primitive operations and proceeding through
category-theoretic constructions, normalization theorems, admissibility
preservation, and representation invariance. Volume~III shifts from
abstract mathematics to computational application, reviewing established
results from mechanistic interpretability before introducing activation
trajectory analysis, the MEM$|$8 framework, and empirical diagnostic
procedures for contemporary neural architectures. Volume~IV situates the
resulting framework within the broader landscape of philosophy, cognitive
science, artificial intelligence, and applied mathematics, concluding
with explicit limitations, non-claims, and open research questions. The
appendices provide complete proofs, notation tables, implementation
details, and supplementary mathematical derivations referenced throughout
the main text.

\section{Why Structure Rather Than Meaning?}

Before turning to the historical survey, it is worth pausing on a choice
that has been implicit throughout this chapter: the theory developed here
does not begin by asking what meaning \emph{is}. It begins by asking what
relations among representations remain invariant as those representations
are transformed. This is a substantive methodological decision, and it is
worth making the reasoning behind it explicit rather than leaving it to
be inferred from the sequence of definitions.

A direct definition of meaning inherits every disagreement already
present in the literature surveyed in Chapter~2: whether meaning requires
reference, mental content, use, verification conditions, or truth
conditions is itself a live dispute, and any definition chosen at the
outset would import that dispute into the formal development rather than
resolving it. Structure, by contrast, is comparatively tractable. Whether
a transformation preserves a given boundary, whether two admissible
sequences reduce to the same normal form, whether a projection preserves
admissibility --- these are questions with determinate answers once the
relevant objects are specified, regardless of one's antecedent theory of
meaning.

This is not a claim that structure exhausts meaning, or that a full
theory of meaning could be built from structural invariants alone.
Chapter~17 states explicitly that Structural Semantics does not attempt
such a reduction. The claim is narrower and more defensible: whatever
else meaning involves, semantic systems exhibit organization that persists
across representational change, and that organization can be studied on
its own terms, prior to and independently of settling what meaning
ultimately is. Beginning with structure therefore is not a wager that
structure is all there is. It is a decision about where a tractable
mathematical investigation can honestly start.

\section{Concluding Remarks}

The principal objective of this introductory chapter has been
methodological rather than technical. Before constructing a formal theory
of Structural Semantics, it has been necessary to delimit the domain
within which the theory is intended to operate. Beginning with the
observation that the term \emph{understanding} possesses multiple
incompatible meanings across philosophy, cognitive science, and
artificial intelligence, we argued that many contemporary disagreements
arise from treating these distinct meanings as though they referred to a
single mathematical object. Rather than attempting to resolve
longstanding debates concerning phenomenal consciousness or
intentionality, the present work isolates structural semantic capacity as
an independent object of investigation.

The resulting framework proceeds from computational systems to
representations, from representations to transformations, from
transformations to boundary constraints, and finally toward bounded scope
geometries; each concept introduces precisely the mathematical structure
required for the next while avoiding unnecessary assumptions regarding
implementation or physical substrate. Consequently, no substantive claims
concerning transformer architectures, neural representations, or
signal-processing models have yet been made, since those developments
depend on mathematical results not yet established. Maintaining this
logical order is an important methodological principle throughout the
monograph.

\subsection*{Status of Results: Chapter 1}

\paragraph{Established background.} This chapter reviewed established
literature concerning philosophy of mind, artificial intelligence,
cognitive science, and machine learning, including classical discussions
of understanding \cite{Turing1950,Searle1980,Block1995}, predictive
processing \cite{Friston2010,Clark2013}, modern deep learning
\cite{Goodfellow2016,Vaswani2017}, and mechanistic interpretability
\cite{Olah2020,Elhage2021}.

\paragraph{Formal developments.} The chapter introduced the
Disentanglement Operator, Boundary Constraint, Bounded Context, and
preliminary Scope Geometry as original definitions establishing the
mathematical domain of Structural Semantics, and stated the primary
research questions and overall methodological architecture of the
monograph.

\paragraph{Empirical hypotheses.} None tested here. The only
forward-looking claim is that structural semantic preservation should
admit computational measurement independently of behavioral fluency, a
proposal developed formally in Volume~III and subjected to explicit
falsification protocols in Chapter~15.

%%%%%%%%%%%%%%%%%%%%%%%%%%%%%%%%%%%%%%%%%%%%%%%%%%%%%%%%%%%%%%%%%%%%%%%%%%%%%%
\chapter{Historical Foundations}
\label{ch:historical-foundations}
%%%%%%%%%%%%%%%%%%%%%%%%%%%%%%%%%%%%%%%%%%%%%%%%%%%%%%%%%%%%%%%%%%%%%%%%%%%%%%

\begin{center}
\fbox{\parbox{0.92\textwidth}{\textbf{Chapter Roadmap.} This chapter
reconstructs the intellectual lineage from which Structural Semantics
emerges, from Saussure through mechanistic interpretability, identifying
for each tradition what it contributes and what remains to be
formalized.}}
\end{center}

\vspace{1em}
\begin{flushright}\itshape
The value of a sign results solely from the simultaneous presence of
other signs.\\
\upshape --- Ferdinand de Saussure, \emph{Cours de Linguistique G\'en\'erale}
\end{flushright}
\vspace{1em}

\section{Introduction}

The purpose of this chapter is analytical reconstruction rather than mere
historical exposition. No existing framework discussed here is dismissed
as incorrect; rather, each contributes an important structural component
later incorporated into the Scope Calculus of Volume~II. Structural
linguistics contributes relational organization, formal semantics
contributes compositional interpretation, process philosophy contributes
dynamic evolution, category theory contributes compositional abstraction,
and mechanistic interpretability contributes empirical access to internal
representations. The central thesis of this chapter is that these
traditions become considerably more compatible once semantic organization
is treated as a problem of admissible structural transformation rather
than static representation. Throughout, we ask of each tradition which
mathematical object is taken to be primitive, how semantic organization
is represented, and which structural properties remain invariant under
transformation.

\section{The Structuralist Tradition}
\label{sec:structuralist-tradition}

The expression \emph{structural semantics} has a well-established history
within twentieth-century linguistics. Beginning with Ferdinand de
Saussure, structuralist approaches emphasized that the meaning of
linguistic units derives not from isolated reference but from their
position within systems of relational differences \cite{Saussure1916},
motivating a body of work including lexical field theory, componential
analysis, relational semantics, and later structural linguistics
\cite{Lyons1977,Greimas1966,Hjelmslev1961,Coseriu1977}. Despite their
differences, these approaches share the methodological commitment that
semantic phenomena are explained by analyzing the relational organization
of linguistic structures rather than treating individual words as
independent semantic atoms.

The present monograph adopts this structuralist insight while extending
it beyond linguistic objects: rather than studying lexical systems,
semantic fields, or componential features, it investigates admissible
transformations over bounded computational contexts, so that the primary
mathematical objects become transformations, boundary constraints, and
scope geometries rather than lexical relations. The adjective
``structural'' therefore retains its original emphasis on relations, but
the mathematical domain is fundamentally different: where classical
structural semantics studies relations among signs, Structural Semantics
as developed here studies relations among admissible computational
transformations.

\begin{table}[ht]
\centering
\begin{tabular}{ll}
\toprule
Classical Structural Semantics & Structural Semantics (this work) \\
\midrule
Lexical units & Computational states \\
Semantic fields & Bounded scope geometries \\
Sense relations & Admissible transformations \\
Componential features & Boundary constraints \\
Lexical structure & Constraint geometry \\
Meaning relations & Structural preservation \\
\bottomrule
\end{tabular}
\caption{Comparison between classical structural semantics and the
framework developed in this monograph.}
\end{table}

\begin{definition}[Classical structural semantics]
\label{def:classical-structural-semantics}
Classical structural semantics denotes the family of linguistic theories
that analyze meaning through systems of structural relations among
linguistic units, including lexical fields, semantic features, and
relational networks \cite{Saussure1916,Hjelmslev1961,Greimas1966,Lyons1977}.
\end{definition}

\begin{definition}[Structural Semantics, this work]
Throughout this monograph, \emph{Structural Semantics} denotes the study
of semantic preservation under admissible transformations of bounded
computational contexts. Unless stated otherwise, the term refers to this
mathematical framework rather than to \cref{def:classical-structural-semantics}.
\end{definition}

\section{Ferdinand de Saussure and Relational Structure}

Modern structural approaches to language originate largely with
Saussure's posthumously published \emph{Cours de Linguistique
G\'en\'erale} (1916) \cite{Saussure1916}. Saussure rejected the view that
linguistic meaning is intrinsic to isolated words, arguing instead that
linguistic values arise from systems of mutual distinction: a sign
consists of a signifier and a signified, neither possessing independent
meaning in isolation, with meaning emerging through the network of
differences separating one sign from every other sign in the system. From
the perspective developed in this monograph, Saussure's insight may be
read as an early recognition that semantic systems are constraint
systems, in which meaning is preserved not because individual symbols
remain unchanged but because admissible structural relationships remain
stable under transformation.

\begin{observation}[Relational priority]
Saussure's structural linguistics assigns explanatory priority to
relations rather than isolated linguistic entities. The Scope Calculus
extends this principle by treating admissible transformations of those
relations as the primary mathematical objects of investigation.
\end{observation}

Saussure's framework nevertheless remained largely descriptive: it
introduced no explicit algebra governing semantic transformations, no
theory of nested contextual organization, and no mathematical account of
how structural relationships evolve during reasoning or computation.
These omissions lay outside the objectives of early structural
linguistics, but they identify precisely the mathematical questions
addressed later in this monograph, beginning with the formal development
of bounded scope geometries in Chapter~6.

\section{Louis Hjelmslev and Formal Structural Systems}

Louis Hjelmslev sought to give structural linguistics a more rigorous
formal foundation, attempting in his \emph{Prolegomena to a Theory of
Language} (1943) to construct linguistics as an autonomous discipline
based on explicit structural principles rather than psychological or
historical explanation \cite{Hjelmslev1943}. Central to his account is the
distinction between the planes of expression and content, each possessing
its own internal organization, with linguistic meaning arising through
their systematic correlation; and the distinction between form, the
abstract organization governing a system, and substance, the material
realization through which that form is implemented, so that the same
structural organization may be realized through different physical media
without altering the underlying linguistic system. This distinction
foreshadows the Representation Invariance Theorem of Chapter~9, which
shows in a different mathematical setting that semantic equivalence
depends on preserved structural transformations rather than on a
particular representational substrate.

Hjelmslev's theory remains primarily descriptive rather than operational:
it introduces no explicit notion of admissible state transformation, no
rewriting calculus governing structural evolution, and no topology of
contextual boundaries. The Scope Calculus of Volume~II may be understood
as extending structural linguistics from static organization toward a
fully operational theory of constrained semantic dynamics.

\section{John Lyons and Structural Semantics}

The expression \emph{structural semantics} entered modern linguistic
usage primarily through the work of John Lyons \cite{Lyons1963,Lyons1977},
who shifted attention from the internal organization of linguistic signs
toward the systematic organization of lexical meaning itself, arguing
that lexical meaning is determined by the network of semantic relations
within which each lexical item participates rather than by isolated
dictionary definitions. Lexical field theory studies collections of
semantically related expressions whose meanings are mutually constrained,
and Lyons further emphasized systematic sense relations such as synonymy,
antonymy, hyponymy, incompatibility, entailment, and meronymy, which
organize vocabulary into structured semantic networks and may be read as
constraints restricting the admissible configurations of lexical
organization.

The present monograph agrees with Lyons that semantic organization is
relational, but extends this observation in two directions: the primitive
objects of the theory are bounded contextual scopes rather than lexical
items, and the primary concern is not merely the existence of semantic
relations but the admissible transformations that preserve those
relations during computation. Whereas structural semantics in the
linguistic tradition describes the organization of semantic systems, the
Scope Calculus seeks to describe the dynamics of semantic organization
itself.

\section{The Shift from Static Structures to Dynamic Structures}

The three traditions surveyed so far --- Saussure, Hjelmslev, and Lyons
--- share a family resemblance worth naming explicitly before the
chapter turns to Frege, Peirce, and Whitehead. Each studies a
\emph{fixed} relational system: a language at a given moment, whose signs
stand in relations of difference, whose planes of expression and content
are internally organized, whose lexical items occupy a stable network of
sense relations. The object of analysis is the system's synchronic
organization, not its evolution. Saussure himself insisted on this
separation, distinguishing the synchronic study of a language state from
the diachronic study of its historical change, and largely restricting
structural analysis to the former.

This emphasis was appropriate to the questions classical structuralism
asked. It becomes a limitation only relative to a different question,
the one this monograph asks: not how a system of distinctions is
organized at a moment, but how that organization behaves as the system
undergoes transformation. A bounded scope geometry, introduced formally
in Chapter~6, is in one respect a direct descendant of the Saussurean
relational system --- it too specifies a structure of admissible
distinctions rather than treating entities as meaningful in isolation.
What it adds is that the geometry is defined together with the
transformations that act on it, so that admissibility becomes a property
of transitions rather than a property of a single static configuration.

The remaining sections of this chapter --- Frege's compositional
combination of meanings, Peirce's semiosis as a chain of successive
interpretants, and above all Whitehead's replacement of substance with
process --- can be read as successive steps away from purely synchronic
description and toward treating transformation itself as the primary
object of study. None of them frames the shift in exactly Saussure's
terms, and none anticipates the Scope Calculus specifically. But by the
time the chapter reaches Whitehead, the historical trajectory it traces
has already moved a considerable distance from the static relational
systems with which it began, and Part~II inherits that trajectory rather
than starting it.

\section{Gottlob Frege and Compositional Semantics}

Although Frege predates twentieth-century structural linguistics, his
work established many of the foundations of modern formal semantics
\cite{Frege1892}. His distinction between sense (\emph{Sinn}) and
reference (\emph{Bedeutung}) observes that two expressions, such as
``Morning Star'' and ``Evening Star,'' may refer to the same object while
presenting it under different conceptual descriptions, so that semantic
analysis cannot be reduced solely to reference. Frege's Principle of
Compositionality further holds that the meaning of a complex expression
depends on the meanings of its constituents and the rule by which they
are combined, a principle underlying formal semantics, type theory,
programming-language semantics, and modern logic.

The Scope Calculus accepts that semantic organization is compositional,
but generalizes compositionality beyond syntactic construction: rather
than asking only how meanings combine, it investigates whether the
transformations acting on those meanings preserve admissible structural
constraints, so that composition becomes one instance of a broader
mathematical problem concerning the preservation of semantic organization
under transformation. Frege's framework primarily concerns the
interpretation of completed expressions; it provides no explicit theory
of evolving contextual states and no algebra governing semantic
transformations, gaps that the Scope Calculus is built to address while
building directly on the compositional insight.

\section{Charles Sanders Peirce and Triadic Semiotics}

Charles Sanders Peirce developed a broader semiotic framework intended to
explain how signs function across reasoning, communication, scientific
inquiry, and cognition \cite{Peirce1931,Short2007}, rejecting the view
that signs exist merely as binary relations between words and objects.
Every act of signification instead involves a triadic relation of sign,
object, and interpretant, and because an interpretant may itself function
as a new sign, interpretation may propagate through a chain Peirce called
\emph{unlimited semiosis}: meaning develops through an evolving network
of successive interpretations rather than a single isolated semantic
event.

Peirce's semiotics anticipates several themes developed later in this
monograph, particularly its emphasis on process over static
representation, but his interpretant remains a philosophical concept
describing semantic mediation rather than an explicit mathematical
operator. The Scope Calculus instead describes the structural constraints
under which interpretation may evolve without producing contextual
inconsistency, so that interpretation becomes one possible realization of
a more general theory of admissible semantic evolution. Peirce's emphasis
on process provides an important bridge to Whitehead, examined next, whose
metaphysics treats process itself as the primitive constituent of
reality.

\section{Alfred North Whitehead and Process Philosophy}

Whitehead's \emph{Process and Reality} (1929) sought to replace
traditional substance metaphysics with a comprehensive process ontology
\cite{Whitehead1929}. Where classical ontologies treat objects as
persisting substances undergoing accidental modification, Whitehead
argued that the fundamental constituents of reality are events of
becoming, with stability emerging from recurring patterns of organized
process rather than from immutable substance --- an inversion from
objects to processes that becomes central to the categorical translation
developed in Chapter~10.

The primitive unit of Whitehead's ontology is the \emph{actual entity} or
\emph{actual occasion}: a transient event that arises, integrates
information from prior occasions through \emph{prehension}, achieves
completion (\emph{satisfaction}), and contributes to future processes.
Whitehead distinguishes positive prehensions, which contribute information
to the developing entity, from negative prehensions, which exclude it,
and calls the process by which multiple prehensions integrate into a
single completed entity \emph{concrescence}. Among these components,
\emph{subjective form} --- the manner in which a datum is experienced
during prehension --- is explicitly phenomenal, and cannot simply be
removed without altering the ontology; \cref{ch:whitehead,ch:phenomenological-boundary}
ask whether the structural organization of prehensions can be formalized
independently of subjective form, and answer affirmatively.

Whitehead's framework shares with the present work the conviction that
transformation is more fundamental than static state and that complex
structures emerge from compositions of simpler relations, but it develops
a metaphysical ontology intended to describe the general structure of
reality, whereas the Scope Calculus develops a mathematical theory of
admissible transformations acting on bounded contextual structures. The
present work therefore neither adopts nor rejects Whitehead's
metaphysical conclusions, extracting instead those structural components
capable of rigorous mathematical formalization, a project undertaken in
full in Chapters~4, 5, and~10.

\section{John Searle and the Chinese Room}

Among the most influential critiques of computational theories of mind is
Searle's Chinese Room argument \cite{Searle1980}, directed against the
claim that the correct execution of a computer program is itself
sufficient for understanding. Searle imagines a person who does not
understand Chinese, enclosed in a room with a rule book, manipulating
Chinese symbols according to purely formal instructions and producing
output indistinguishable from that of a fluent speaker; to external
observers the room appears to understand Chinese, but Searle argues that
no genuine understanding exists, since the individual manipulates symbols
syntactically without attaching semantic significance to them. Syntax
alone, on this view, is insufficient for semantics.

The present monograph neither accepts nor rejects Searle's conclusions
regarding phenomenal consciousness, arguing instead that the Chinese Room
addresses a different question from the one investigated here: Searle
asks whether formal symbol manipulation is sufficient for subjective
understanding, while the Scope Calculus asks whether semantic
organization can be mathematically characterized through the preservation
of admissible structural transformations. These questions are logically
independent, and this distinction motivates the epistemic separation
introduced in Chapter~1 and developed formally in Chapter~5. Rather than
answering Searle's question directly, Structural Semantics replaces the
binary question of whether a machine understands with the more
mathematically tractable question of which structural properties of a
computational system remain invariant under admissible semantic
transformation.

\section{Ned Block and the Separation of Consciousness}

Ned Block's distinction between \emph{phenomenal consciousness} and
\emph{access consciousness} \cite{Block1995} clarified that many
philosophical disagreements concerning artificial intelligence arise
because these fundamentally different notions of consciousness are
routinely conflated. Phenomenal consciousness refers to subjective
qualitative experience, what Nagel called ``what it is like'' to be a
conscious organism \cite{Nagel1974}\footnote{The expression itself
predates Nagel's 1974 essay: it appears in B.~A.~Farrell's 1950 paper
\emph{Experience}, and earlier still in Wittgenstein's remarks of
1946--47 and Russell's 1926 \emph{Encyclopaedia Britannica} entry on
psychology. Nagel's essay is credited here, as it standardly is, for
making the phrase central to the philosophy of mind and the explanatory
gap, not for originating it.}; access consciousness, by contrast,
concerns the availability of information for reasoning, planning, memory,
verbal report, and behavioral control, and is defined functionally,
admitting computational description far more readily than phenomenal
consciousness.

Structural Semantics intersects Block's framework at a carefully
delimited point: the Scope Calculus is not a theory of phenomenal
consciousness, nor a complete theory of access consciousness, but
investigates a third mathematical object, the geometry of admissible
semantic transformations. Information accessibility may facilitate
structural semantic organization, but structural semantic organization is
defined independently, through the preservation of admissible boundaries
(\cref{def:disentanglement}), so that the present framework neither
requires nor denies phenomenal consciousness and does not require any
particular cognitive architecture for access consciousness.

\section{Predictive Processing and the Free Energy Principle}

During the early twenty-first century, cognitive science increasingly
modeled intelligence as a continuous process of prediction, error
correction, and dynamical adaptation. Predictive Processing, developed by
Clark, Hohwy, and others \cite{Clark2013,Hohwy2013}, and the closely
related Free Energy Principle proposed by Friston \cite{Friston2010,
Friston2019}, both treat cognition as active inference whose objective is
the maintenance of coherent internal models under uncertain sensory
information, describing cognition through probabilistic inference,
dynamical systems, and Bayesian optimization rather than explicit logical
manipulation. Perception is understood not as passive reception but as
continual generation of predictions compared against observations, with
cognition inferring the posterior $P(x\mid y)$ over latent causes $x$
given observations $y$. Friston generalizes this into the minimization of
an upper bound on Bayesian surprise called variational free energy,
\[
F(q)=D_{\mathrm{KL}}\!\left(q(x)\,\middle\|\,P(x\mid y)\right)-\log P(y),
\]
where $q(x)$ is an internal recognition density approximating the true
posterior; minimizing $F$ approximately minimizes surprise while
improving the accuracy of the internal generative model, unifying
perception, learning, action, and homeostatic regulation. Active
Inference extends the framework so that an organism may also act on its
environment to reduce prediction error, replacing stimulus-response
models with a closed feedback loop between organism and environment.

Predictive Processing unifies perception, learning, memory, and motor
control within a common inferential framework, accommodates uncertainty
naturally, scales to hierarchical architectures, and admits quantitative
modeling; but although it explains how internal models evolve under
prediction error, it provides comparatively little formal machinery
describing the organization of semantic structure itself, since
prediction error measures numerical discrepancy rather than the
admissibility of contextual transformations or the preservation of
semantic boundaries. Structural Semantics shares with Predictive
Processing an emphasis on internal organization and stability under
continual transformation, but its primary mathematical object differs:
Predictive Processing minimizes variational free energy over probability
distributions, while Structural Semantics studies admissible
transformations over bounded scope geometries, with prediction error
becoming only one possible signal influencing boundary operations rather
than the defining quantity itself. The two approaches are therefore
potentially complementary, a relationship revisited systematically in
Chapter~16.

\section{Mechanistic Interpretability}

The rapid success of transformer architectures introduced a new
scientific challenge: while large language models demonstrated remarkable
behavioral competence, the internal mechanisms producing this behavior
remained largely opaque. Mechanistic interpretability emerged as a
research program seeking causal explanations grounded directly in the
internal organization of the model, rather than treating neural networks
as inscrutable black boxes evaluated only on input-output performance.

Modern transformer architectures consist of repeated blocks of multi-head
attention, feed-forward networks, normalization, and residual connections
\cite{Vaswani2017}; if the residual stream at layer $l$ is
$\mathbf{x}_l\in\mathbb{R}^{d_{\mathrm{model}}}$, the forward computation
is the iterative transformation $\mathbf{x}_{l+1}=T_l(\mathbf{x}_l)$,
forming a trajectory through a high-dimensional representation space
rather than a sequence of explicitly symbolic operations. Recent work by
Olah, Elhage, Nanda, and collaborators suggests that many such
computations can be understood in terms of interacting computational
circuits --- induction heads, copy circuits, name-mover circuits,
arithmetic circuits, factual retrieval pathways, and long-range attention
mechanisms among them \cite{Olah2020,Elhage2021,Nanda2023} --- indicating
that distributed neural computation often possesses substantial internal
organization despite the absence of explicit symbolic programming.
Techniques including activation patching, causal tracing, feature
visualization, sparse autoencoders, attention intervention, linear
probing, and circuit analysis \cite{Olah2020,Elhage2021,Marks2024} attempt
to identify stable computational features whose causal manipulation
produces predictable behavioral changes, with the overall objective
explanatory rather than merely descriptive.

Despite substantial progress, current interpretability techniques remain
fragmented: many methods identify individual features, neurons, or
circuits without a unified mathematical description of contextual
organization across an entire computational trajectory, and no generally
accepted algebra currently exists for describing admissible semantic
transformations across changing internal representations. Structural
Semantics builds on the objectives of mechanistic interpretability while
proposing a different abstraction: individual neurons, attention heads,
sparse features, or circuits are not themselves treated as semantic
units, but constitute one possible substrate whose state transitions may
be projected onto bounded scope geometries, so that the primary
mathematical object shifts from individual network components to
admissible transformations preserving contextual organization. This
motivates the Representation Invariance Theorem of Chapter~9, which
establishes conditions under which different computational substrates may
be considered structurally equivalent within the Scope Calculus.

\section{Historical Synthesis}

The historical survey reveals three broad explanatory traditions. The
first emphasizes subjective experience as the defining feature of
understanding, including phenomenology, portions of process philosophy,
and accounts centered on consciousness or intentionality. The second
emphasizes functional organization, identifying understanding through
behavioral competence or successful interaction with the environment, as
in cybernetics, functionalism, and contemporary AI benchmarks. The third
emphasizes structural organization, investigating the internal
organization of representations and the transformations that preserve
them; although elements of this viewpoint appear throughout structural
linguistics, information theory, category theory, and mechanistic
interpretability, no single unified mathematical framework currently
develops this approach systematically, and the present monograph is
situated within this third tradition.

Existing theories successfully explain many individual aspects of
semantic organization --- structural linguistics describes relational
organization within languages, category theory studies structure-
preserving mappings, predictive processing models adaptive inference, and
mechanistic interpretability analyzes internal computation --- but none
develops a unified algebra describing how semantic contexts themselves
transform while preserving admissible organizational constraints. This
gap motivates the formal developments beginning in Part~II. Structural
Semantics does not seek to replace the traditions surveyed in this
chapter, but to synthesize insights from several of them: the primacy of
relational organization from structural linguistics, the study of
structure-preserving transformations from category theory, the emphasis
on dynamic evolution from process philosophy, the importance of continual
adaptation from predictive processing, and the commitment to analyzing
internal computational organization from mechanistic interpretability.

\subsection*{Status of Results: Chapter 2}

\paragraph{Established background.} Historical developments in structural
linguistics, process philosophy, philosophy of mind, predictive
processing, active inference, and mechanistic interpretability
\cite{Saussure1916,Trier1931,Lyons1977,Whitehead1929,Searle1980,Block1995,
Clark2013,Friston2010,Vaswani2017,Elhage2021,Nanda2023}.

\paragraph{Formal developments.} A comparative historical synthesis
identifying three major explanatory traditions and motivating the need
for a formal theory of Structural Semantics.

\paragraph{Empirical hypotheses.} None. This chapter is historical and
conceptual, providing the intellectual foundations for the formal
developments beginning in Chapter~3.

%%%%%%%%%%%%%%%%%%%%%%%%%%%%%%%%%%%%%%%%%%%%%%%%%%%%%%%%%%%%%%%%%%%%%%%%%%%%%%
\chapter{Three Definitions of Understanding}
\label{ch:three-definitions}
%%%%%%%%%%%%%%%%%%%%%%%%%%%%%%%%%%%%%%%%%%%%%%%%%%%%%%%%%%%%%%%%%%%%%%%%%%%%%%

\vspace{1em}
\begin{flushright}\itshape
Every definition illuminates one aspect of a phenomenon while obscuring
another. The scientific task is therefore not merely to define, but to
identify which definitions remain appropriate for which explanatory
questions.
\end{flushright}
\vspace{1em}

The preceding chapters established two complementary foundations.
Chapter~1 isolated the epistemic domain of Structural Semantics by
separating semantic capacity from phenomenal consciousness, truth,
agency, and syntactic fluency; Chapter~2 surveyed the historical
development of theories of meaning, cognition, and representation from
structural linguistics through mechanistic interpretability. The
historical record reveals that disagreements concerning artificial
intelligence frequently arise because the word \emph{understanding} is
used to denote fundamentally different properties, so that researchers
often disagree not because they evaluate the same property differently,
but because they evaluate different properties under a common name. The
objective of this chapter is therefore taxonomic rather than polemical:
rather than arguing that one definition is correct and the others
mistaken, we distinguish three conceptually independent notions of
understanding that have emerged throughout the literature ---
phenomenological, functional, and structural --- each capturing a genuine
explanatory question, each admitting different methodologies, and each
possessing different mathematical requirements.

\section{Motivating the Taxonomy}

Throughout the history surveyed in Chapter~2, philosophers have asked
whether machines understand language, whether symbolic manipulation is
sufficient for intelligence, whether consciousness is necessary for
meaning, and whether behavioral competence implies semantic knowledge.
Although these questions appear closely related, they address distinct
levels of explanation: some concern subjective experience, others
observable behavior, still others the internal organization of
representations. Asking whether an artificial system ``understands''
without specifying the intended definition therefore produces an
ill-posed scientific question, and a precise taxonomy becomes a
prerequisite for rigorous discussion. The classification introduced below
is not an exhaustive ontology of cognition, but identifies three broad
explanatory frameworks that recur throughout contemporary philosophy of
mind, cognitive science, and artificial intelligence.

\begin{definition}[Understanding, neutral form]
Let $\System$ denote a system operating over a domain $\Domain$. An
\emph{understanding relation} is a mapping
$U:\System\times\Domain\to\mathcal{P}$, where $\mathcal{P}$ denotes a
collection of explanatory properties whose interpretation depends on the
theoretical framework under consideration.
\end{definition}

The codomain $\mathcal{P}$ is intentionally left abstract: for
phenomenological theories it concerns subjective experience, for
functional theories behavioral competence, and for Structural Semantics
the admissible preservation of contextual organization. This neutrality
allows the three theories to be compared without presupposing the
correctness of any one of them.

The three definitions that follow represent distinct research programs
rather than competing philosophical dogmas, and are not mutually
exclusive: a biological organism may satisfy all three, an artificial
system may satisfy one or two, and the objective here is clarification of
explanatory questions rather than classification of particular systems.
This taxonomy functions as the conceptual interface between the
historical literature of Chapter~2 and the mathematical framework of
Chapters~6 through~10; without it, later formal definitions of Scope
Geometry and structural semantic capacity would remain disconnected from
the broader philosophical literature.

\section{Phenomenological Understanding}

The oldest and perhaps most intuitive conception of understanding
identifies it with conscious experience itself: not merely the successful
manipulation of symbols or the production of appropriate behavioral
responses, but something inseparable from the existence of a first-person
perspective through which meaning is directly experienced. This
perspective appears throughout Husserl's phenomenology, Heidegger's
existential analysis, Merleau-Ponty's embodied cognition, Whitehead's
process philosophy, Nagel's analysis of subjective experience, Searle's
critique of computational intelligence, Chalmers' formulation of the hard
problem, and Block's distinction between phenomenal and access
consciousness \cite{Husserl1913,Heidegger1927,MerleauPonty1945,
Whitehead1929,Nagel1974,Searle1980,Chalmers1996,Block1995}; despite
differences in terminology and metaphysical commitment, these theories
converge on the claim that understanding is inseparable from phenomenal
experience.

\begin{definition}[Phenomenological understanding]
Let $\System$ be a cognitive system operating over a semantic domain
$\Domain$. The system possesses \emph{phenomenological understanding}
over $\Domain$ if and only if successful semantic activity is
accompanied by subjective phenomenal experience, formally
$U_{\mathrm{phen}}:\System\times\Domain\to\Phi$, where $\Phi$ denotes the
domain of phenomenal conscious experience. The precise ontology of $\Phi$
is left unspecified, as different phenomenological traditions provide
different accounts of subjective experience.
\end{definition}

Phenomenological theories generally hold that semantic understanding is
intrinsically first-person, that meaning is not exhausted by behavioral
competence, and that subjective experience is explanatorily fundamental
rather than derivative; consequently they distinguish appearing
intelligent from genuinely understanding, and their best-known
contemporary defense is Searle's Chinese Room, discussed already in
Chapter~2. Within Whitehead's process philosophy, phenomenology enters the
formal structure through Subjective Form, $\WSubjForm$, the manner in
which an occasion experiences its own becoming, formalized fully in
Chapter~4.

These theories capture an undeniable feature of ordinary experience and
raise questions that remain largely unresolved by current neuroscience
and artificial intelligence, but they encounter a substantial
methodological difficulty when applied to computational systems: because
phenomenal experience is intrinsically first-person, no generally
accepted procedure exists for measuring it objectively, so that if
subjective experience is adopted as part of the definition of
understanding, any system lacking independently established phenomenology
is excluded by definition rather than by empirical investigation. This
observation does not demonstrate that artificial systems lack
phenomenology; it indicates only that phenomenological theories alone
provide no operational criterion by which the question can be resolved,
an observation developed into a general argument in Chapter~5. The
present monograph neither accepts nor rejects phenomenological accounts,
isolating them instead as one legitimate explanatory framework among
several, and asks in the chapters that follow whether semantic
organization can also be characterized independently through formal
structural properties.

\section{Functional Understanding}

A second tradition identifies understanding not with subjective
experience but with what a system can reliably do, treating the internal
substrate responsible for cognition as secondary to whether the system
exhibits appropriate causal organization, behavioral competence, adaptive
capacity, or information-processing ability. This perspective has
profoundly influenced cybernetics, cognitive science, artificial
intelligence, robotics, and computational neuroscience, forming the
dominant methodological framework for modern AI research
\cite{Turing1950,Putnam1967,Fodor1975,NewellSimon1976}.

\begin{definition}[Functional understanding]
Let $\System$ be a system operating over a semantic domain $\Domain$. The
system possesses \emph{functional understanding} if its observable
behavior satisfies the operational requirements associated with
successful performance over $\Domain$, formally
$U_{\mathrm{func}}:\System\times\Domain\to\mathcal{B}$, where
$\mathcal{B}$ denotes an appropriate space of behavioral or computational
performance measures. Unlike phenomenological understanding, no
assumptions are made concerning subjective experience or phenomenal
awareness.
\end{definition}

Turing proposed replacing metaphysical questions concerning thought with
operational tests based on conversational behavior \cite{Turing1950};
cybernetics emphasized feedback, control, and adaptive behavior
\cite{Wiener1948,Ashby1956}; and Putnam and Fodor developed forms of
functionalism identifying mental states with their causal roles rather
than their physical realization \cite{Putnam1967,Fodor1975}, naturally
supporting the possibility of multiple realizability, in which distinct
physical systems instantiate the same cognitive organization. Most
contemporary evaluations of large language models remain largely
functionalist, measuring question answering, mathematical reasoning, code
generation, planning, factual recall, and conversational competence
through observable outputs rather than subjective experience, a
methodology that has enabled rapid empirical progress because behavioral
metrics admit objective measurement and direct comparison across systems.

Functional theories provide operational definitions suitable for
empirical investigation, remain largely independent of metaphysical
assumptions concerning consciousness, naturally accommodate substrate
independence, and rely on criteria that are reproducible and
experimentally testable. However, behavior alone does not uniquely
determine internal organization: distinct computational mechanisms may
produce nearly identical external behavior while employing fundamentally
different internal representational structures, so that behavioral
equivalence does not necessarily imply structural equivalence. Consider,
for instance, two theorem-proving systems, one using symbolic logical
deduction and the other a transformer trained on formal proof corpora,
which produce identical formal proofs for every theorem in a benchmark
suite; from a functional perspective the two exhibit equivalent
understanding, yet their internal organizations differ substantially, and
functional evaluation alone does not determine whether they preserve
semantic context through comparable structural mechanisms.

Structural Semantics accepts the functionalist insight that semantic
organization need not depend on any particular physical substrate, but
proposes that behavioral competence alone is insufficient for a complete
mathematical theory of understanding, and instead investigates whether
successive internal state transitions preserve bounded semantic
organization throughout a computation, treating behavior as evidence for
understanding rather than as its complete definition.

\section{Structural Understanding}

The third conception developed in this monograph begins from a different
premise than either the phenomenological or the functional: rather than
identifying understanding with subjective experience or observable
behavior, it treats understanding as a property of the internal
organization of representations and, specifically, of the transformations
that preserve them. This does not deny the importance of consciousness or
reject behavioral evaluation; it isolates a third object of scientific
investigation whose mathematical properties can be studied independently
of both. The central question is therefore not what a system experiences,
nor how successfully it behaves, but which structural relations remain
invariant as it transforms its internal representations --- a shift that
moves the study of semantic organization from psychology toward
mathematics.

\begin{definition}[Structural understanding]
Let $\System$ be a computational or cognitive system operating over a
semantic domain $\Domain$. The system possesses \emph{structural
understanding} if there exists a family of admissible transformations
$\Trans=\{T_i\}$ acting on bounded scope geometries such that every
transformation preserves the admissibility constraints governing
contextual organization; equivalently,
$U_{\mathrm{struct}}:\System\times\Domain\to\CatScope$, where semantic
capacity is evaluated through admissible morphisms in the category
$\CatScope$ rather than through phenomenal experience or behavioral
observation.
\end{definition}

The defining property of structural understanding is invariance: although
internal representations may change continuously throughout computation,
certain organizational relationships remain preserved, and these
invariant relationships constitute the semantic content of the
representation, precisely characterized in Parts~II and~III. Structural
Semantics proposes that semantic organization is fundamentally a problem
of scope preservation --- of maintaining coherent contextual organization
while new information is introduced, previous assumptions are revised,
and nested contexts evolve --- which motivates the bounded scope
geometries of Chapter~6 and the primitive operations of Chapter~7. A
defining feature of the framework is its independence from particular
representational substrates: whether information is encoded symbolically,
neurally, graphically, or through another computational medium is
secondary to the induced transformation over bounded scope geometries, a
principle formalized later as the Representation Invariance Theorem
(Theorem~9.7).

Consider again two reasoning systems proving the same mathematical
theorem, one by symbolic deduction and the other by continuous
optimization over neural activation vectors: although their internal
representations differ radically, both may maintain identical
assumptions, preserve identical logical dependencies, introduce
hypotheses in the same admissible order, and discharge those assumptions
according to equivalent structural rules. From the perspective of
Structural Semantics, both systems instantiate the same semantic
organization despite their differing computational substrates, since
their semantic equivalence depends on preserved contextual structure
rather than on identical internal encodings.

\section{Comparing the Three Conceptions}

The three conceptions of understanding introduced in this chapter may now
be compared directly. Phenomenological understanding asks whether
semantic activity is accompanied by subjective experience; functional
understanding asks whether semantic activity produces appropriate
observable behavior; structural understanding asks whether semantic
activity preserves the organization of contextual constraints throughout
computation. These questions are logically independent: a biological
organism may satisfy all three, a theorem prover may satisfy the latter
two while remaining silent with respect to phenomenology, and a
hypothetical conscious system suffering severe cognitive impairment might
possess phenomenology while exhibiting poor functional or structural
organization. The taxonomy therefore separates explanatory questions
rather than ranking competing theories.

\begin{table}[h]
\centering
\begin{tabular}{p{3.6cm}p{3cm}p{3cm}p{4.4cm}}
\toprule
Framework & Primary object & Evaluation & Representative literature \\
\midrule
Phenomenological & Subjective experience & First-person awareness &
Husserl, Heidegger, Whitehead, Nagel, Chalmers \\
Functional & Behavioral competence & Observable performance &
Turing, Putnam, Fodor, Newell, Simon \\
Structural & Contextual organization & Admissible transformations &
Scope Calculus, category theory, mechanistic interpretability \\
\bottomrule
\end{tabular}
\caption{Three complementary conceptions of understanding.}
\end{table}

\section{Chapter Summary}

This chapter established a conceptual taxonomy separating three distinct
research programs concerning understanding: phenomenological theories
identify understanding with subjective experience, functional theories
with successful behavioral organization, and Structural Semantics
introduces a third conception in which semantic capacity is identified
with the preservation of admissible contextual organization under
transformation. The remainder of this monograph develops this third
perspective systematically, beginning in Chapter~4 with a mathematical
reconstruction of Whitehead's process ontology that later permits
comparison with the Scope Calculus through the translation functor of
Chapter~10.

\subsection*{Status of Results: Chapter 3}

\paragraph{Established background.} Phenomenological, functional, and
structural traditions in philosophy of mind, cognitive science, and
artificial intelligence \cite{Turing1950,Searle1980,Block1995,
Whitehead1929}.

\paragraph{Formal developments.} Definitions of Phenomenological
Understanding, Functional Understanding, and Structural Understanding,
together with their comparative taxonomy.

\paragraph{Empirical hypotheses.} Structural semantic capacity
constitutes an independently measurable property that is neither
reducible to subjective phenomenology nor fully captured by behavioral
performance alone. Subsequent chapters develop the formal mathematical
machinery necessary to evaluate this hypothesis.

%%%%%%%%%%%%%%%%%%%%%%%%%%%%%%%%%%%%%%%%%%%%%%%%%%%%%%%%%%%%%%%%%%%%%%%%%%%%%%
\chapter{The Whiteheadian Framework}
\label{ch:whitehead}
%%%%%%%%%%%%%%%%%%%%%%%%%%%%%%%%%%%%%%%%%%%%%%%%%%%%%%%%%%%%%%%%%%%%%%%%%%%%%%

\vspace{1em}
\begin{flushright}\itshape
The many become one, and are increased by one.\\
\upshape --- Alfred North Whitehead, \emph{Process and Reality}
\end{flushright}
\vspace{1em}

The preceding chapter established that discussions of understanding
separate into phenomenological, functional, and structural frameworks.
This chapter returns to one of the most influential systems within the
phenomenological tradition: Whitehead's process philosophy. Although the
present monograph ultimately develops a different formal framework,
Whitehead's emphasis on dynamic organization, relational structure, and
process makes his work unusually compatible with modern mathematical
approaches to computation, and this chapter is accordingly presented not
as a critique but as a careful mathematical reconstruction. Its objective
is twofold: to provide precise formal definitions of Whitehead's
principal concepts using contemporary notation, and to identify exactly
which components admit structural formalization and which remain
explicitly phenomenological, preparing the translation functor developed
in Chapter~10.

\section{Process Rather Than Substance}

Classical metaphysics generally treats enduring substances as the primary
objects of reality, with processes and changes regarded as properties or
modifications of an underlying entity. Whitehead reversed this
relationship: reality consists fundamentally of events, and objects are
abstractions constructed from sufficiently stable patterns within those
events. If processes are primitive, transformation becomes more
fundamental than state and relations more fundamental than isolated
entities, an orientation shared by category theory, dynamical systems,
network theory, and process algebra, though Whitehead developed his
system decades before these fields matured.

\section{Actual Occasions}

The primitive object of Whitehead's ontology is the \emph{actual
occasion}: not a material object but a single completed event of
becoming, from which all larger structures --- organisms, mountains,
societies, even physical particles --- are understood as persistent
patterns generated by vast collections of interconnected occasions.

\begin{definition}[Actual occasion]
An actual occasion is represented by an ordered tuple
$\WOccasion=\langle \WData,\WPrehensions,\WSatisfaction\rangle$, where
$\WData$ denotes the collection of antecedent data inherited from prior
occasions, $\WPrehensions$ denotes the collection of prehensions
performed on that data, and $\WSatisfaction$ denotes the completed
satisfaction state produced after concrescence.
\end{definition}

This definition is intentionally schematic; the remaining sections
formalize each component individually.

\section{Prehension}

Whitehead replaces classical perception with the broader concept of
\emph{prehension}, the fundamental relational operation through which one
occasion incorporates aspects of previous occasions into its own becoming.
Unlike ordinary perception, prehension is not restricted to conscious
organisms: every actual occasion prehends antecedent reality, so that
prehension functions as the primitive relational operator of Whitehead's
ontology.

\begin{definition}[Prehension]
A prehension is represented by the ordered triple
$\WPrehensions=\langle \WData,\WValuation,\WSubjForm\rangle$, where
$\WData$ denotes the datum being prehended, $\WValuation\in\{+,-\}$
denotes the valuation polarity, with positive valuation indicating
inclusion and negative valuation indicating exclusion, and $\WSubjForm$
denotes the Subjective Form associated with the prehension.
\end{definition}

Whitehead distinguishes positive prehensions, which contribute directly
to the formation of the new occasion, from negative prehensions, which
exclude data from further participation. Formally, if
$\Domain=\{d_1,\ldots,d_n\}$ denotes the available antecedent data, the
process of becoming partitions this collection into disjoint subsets
$\Domain=P^{+}\cup P^{-}$ with $P^{+}\cap P^{-}=\varnothing$: positive
prehensions contribute to the evolving occasion, negative prehensions
remove information from subsequent construction. This distinction will
later admit a natural comparison with the $\Bind$ and $\Refuse$
operations of the Scope Calculus, although no such identification is
assumed at this stage.

\section{Concrescence}

The collection of individual prehensions does not remain an unordered
set; instead it undergoes a process of integration Whitehead calls
\emph{concrescence}, transforming many inherited influences into one
completed occasion, symbolically
$\{P_1,\ldots,P_n\}\longrightarrow \WSatisfaction$, summarized in
Whitehead's own phrase: ``the many become one, and are increased by
one.'' From a mathematical perspective, concrescence resembles an
aggregation operator acting on a structured family of relational inputs;
precisely which algebraic properties such an operator possesses is a
question examined later through the Scope Calculus.

\section{Subjective Form}

Among the components of Whitehead's ontology, Subjective Form occupies a
special position. Whereas antecedent data and valuation polarity describe
observable structural relations, Subjective Form describes the manner in
which the occasion experiences its own process of becoming, and unlike
positive and negative prehension it is not introduced as an externally
observable property but concerns the intrinsic character accompanying
each prehension. Consequently Subjective Form functions as the explicitly
phenomenological component of Whitehead's ontology; the present monograph
seeks to determine which components of his system admit purely structural
reconstruction and which remain irreducibly phenomenological, preserving
Subjective Form explicitly rather than eliminating it, while carefully
distinguishing it from the structural operations that surround it.

\section{Chapter Summary}

This chapter reconstructed the fundamental concepts of Whitehead's
process ontology using contemporary mathematical notation: actual
occasions as primitive events of becoming; prehensions as ordered
triples of antecedent data, valuation polarity, and Subjective Form;
positive and negative prehensions as an explicit set-theoretic
decomposition; and concrescence as the aggregation process transforming
many prehensions into one completed satisfaction state. Most importantly,
Subjective Form was isolated as the explicitly phenomenological component
of the framework. The following chapter examines this distinction in
detail, showing that Subjective Form functions as an independent axiom
within Whitehead's ontology and forms the critical bridge between
phenomenological theories of understanding and the structural framework
developed throughout the remainder of the monograph.

\subsection*{Status of Results: Chapter 4}

\paragraph{Established background.} Whitehead's process ontology, actual
occasions, prehensions, concrescence, and Subjective Form
\cite{Whitehead1929}.

\paragraph{Formal developments.} Mathematical reconstruction of actual
occasions, prehensions, valuation polarity, and concrescence using
explicit definitions suitable for later categorical translation.

\paragraph{Empirical hypotheses.} None. This chapter reconstructs an
existing philosophical framework in preparation for the structural
analysis developed in Chapters~5--10.

%%%%%%%%%%%%%%%%%%%%%%%%%%%%%%%%%%%%%%%%%%%%%%%%%%%%%%%%%%%%%%%%%%%%%%%%%%%%%%
\chapter{The Phenomenological Boundary}
\label{ch:phenomenological-boundary}
%%%%%%%%%%%%%%%%%%%%%%%%%%%%%%%%%%%%%%%%%%%%%%%%%%%%%%%%%%%%%%%%%%%%%%%%%%%%%%

\vspace{1em}
\begin{flushright}\itshape
Everything is vague to a degree you do not realize till you have tried to
make it precise.\\
\upshape --- Bertrand Russell
\end{flushright}
\vspace{1em}

The preceding chapter reconstructed Whitehead's process ontology without
either criticism or reinterpretation, identifying its primitive
components in their own terms and observing that the architecture
naturally separates into structural entities --- actual occasions,
prehensions, valuation polarity, and concrescence, which admit explicit
mathematical representation --- and Subjective Form, introduced not as an
externally observable relation but as the intrinsic manner in which an
occasion experiences its own becoming. The present chapter examines the
consequences of this distinction. The question is not whether Subjective
Form exists, nor whether Whitehead's metaphysics is correct, but
methodological: if Subjective Form is treated as an indispensable
prerequisite for semantic capacity, the evaluation of artificial systems
becomes independent of any mathematical or empirical investigation, and
the conclusion that such systems lack understanding follows directly from
the initial definition rather than from subsequent analysis. The purpose
of this chapter is to isolate the precise logical boundary separating
structural questions from phenomenological axioms, so that the
mathematical framework developed in subsequent chapters can be evaluated
independently of broader metaphysical commitments.

\section{Structural and Phenomenological Components}

The reconstruction of Chapter~4 naturally partitions Whitehead's ontology
into two classes. The first consists of objects whose behavior may be
described through explicit structural relations: $\WData$, $\WPrehensions$,
$\WValuation$, and $\WSatisfaction$ each admit mathematical description,
whether as sets, graphs, tensors, categories, discrete operators, or
aggregation functionals, and each possesses an externally describable
operational role. The second class contains only Subjective Form,
$\WSubjForm$: unlike every other component, it is not introduced through
observable transformation rules but denotes the intrinsic character
accompanying each prehension, so that its inclusion within the theory is
fundamentally different from every other primitive introduced thus far.

\begin{definition}[Structural component]
A primitive object is \emph{structural} if its contribution to the
behavior of a system is completely specified by externally observable
relations and transformation laws.
\end{definition}

\begin{definition}[Phenomenological component]
A primitive object is \emph{phenomenological} if its defining properties
concern intrinsic experience rather than externally observable structural
relations.
\end{definition}

Under these definitions $\WData$, $\WPrehensions$, $\WValuation$, and
$\WSatisfaction$ are structural components, while $\WSubjForm$ is
phenomenological. The distinction is purely methodological: it does not
deny the existence of phenomenological properties, but identifies which
primitives participate directly in the structural mathematics developed
throughout the remainder of the monograph.

\section{Axiomatic Dependence}

The central logical observation of this chapter concerns the role played
by Subjective Form within theories of understanding. Suppose semantic
capacity is defined by the predicate $\mathrm{Understand}(M)$. If one
adopts the axiom $\WSubjForm\neq\varnothing\Rightarrow\mathrm{Understand}(M)$
together with its converse
$\WSubjForm=\varnothing\Rightarrow\neg\,\mathrm{Understand}(M)$, then the
semantic status of every system is determined immediately by the presence
or absence of Subjective Form, and no subsequent behavioral experiment,
mathematical analysis, or empirical measurement can alter this
conclusion.

\begin{observation}
If Subjective Form is introduced as a necessary condition for semantic
capacity, then every system lacking Subjective Form is excluded by
definition rather than by empirical investigation.
\end{observation}

This observation concerns logical dependence rather than philosophical
correctness; one may certainly choose such an axiom, but the resulting
theory belongs to a different methodological category from one whose
conclusions depend on mathematical or empirical evidence.

\section{Definitions and Empirical Criteria}

Scientific theories generally distinguish between definitions, which
specify the concepts under consideration, and measurements, which
determine whether particular systems satisfy those definitions; when
these two roles are conflated, empirical investigation becomes
impossible. If one defines intelligence as human intelligence, no
experiment comparing humans and machines could demonstrate machine
intelligence, since the conclusion follows directly from the definition;
similarly, if semantic capacity is defined as requiring Subjective Form,
every non-conscious computational architecture is excluded before its
internal organization is examined. Structural Semantics adopts a
different strategy, defining semantic capacity through structural
preservation and treating whether phenomenology accompanies such
preservation as a distinct philosophical question.

\section{The Structural Research Program}

Separating structural organization from phenomenology produces a
well-defined mathematical research program whose central question is
which structural invariants must be preserved for semantic organization
to remain stable. Unlike questions concerning subjective experience, this
question admits explicit mathematical treatment: one may define
admissible transformations, characterize boundary operators, prove
preservation theorems, derive computational diagnostics, and compare
different representational architectures through common structural
invariants. Consequently the mathematical framework developed throughout
Volumes~II and~III does not depend on resolving the philosophical debate
concerning consciousness, studying instead the preservation of
structured relational organization.

This methodological separation is compatible with, while differing from,
several existing traditions: classical functionalism evaluates systems
primarily through behavioral performance \cite{Putnam1967,Fodor1975},
phenomenological theories through subjective experience or intrinsic
intentionality \cite{Nagel1974,Searle1980}, and mechanistic
interpretability through internal computational organization
\cite{Olah2020,Elhage2021}. Structural Semantics occupies a different
position, identifying semantic capacity with neither behavior nor
phenomenology but investigating whether semantic organization can be
characterized by mathematically definable structural invariants, an
orientation that motivates the algebraic development beginning in
Chapter~6.

\section{Chapter Summary}

This chapter identified the precise methodological boundary separating
phenomenological and structural theories of semantic capacity,
distinguishing structural primitives from phenomenological primitives and
showing that Subjective Form occupies a unique position within
Whitehead's ontology: theories requiring Subjective Form as an axiom
classify machine understanding by definition rather than by empirical
investigation. Structural Semantics adopts a complementary research
program, developing a mathematical theory of structural organization
whose predictions may be evaluated independently of phenomenological
commitments. This completes Volume~I. Beginning with Chapter~6, the
monograph transitions from philosophical analysis to the formal
mathematical development of Scope Geometry.

\subsection*{Status of Results: Chapter 5}

\paragraph{Established background.} Whitehead's distinction between
structural process and Subjective Form, together with classical debates
concerning consciousness and intentional content \cite{Whitehead1929,
Nagel1974,Searle1980}.

\paragraph{Formal developments.} Definitions of structural and
phenomenological components, together with the logical separation between
axiomatic phenomenological theories and structural mathematical theories.

\paragraph{Empirical hypotheses.} None. This chapter establishes
methodological distinctions that motivate the formal development of
Scope Geometry in Volume~II.

%%%%%%%%%%%%%%%%%%%%%%%%%%%%%%%%%%%%%%%%%%%%%%%%%%%%%%%%%%%%%%%%%%%%%%%%%%%%%%
\part{Scope Geometry, Algebra, and Categories}
%%%%%%%%%%%%%%%%%%%%%%%%%%%%%%%%%%%%%%%%%%%%%%%%%%%%%%%%%%%%%%%%%%%%%%%%%%%%%%

%%%%%%%%%%%%%%%%%%%%%%%%%%%%%%%%%%%%%%%%%%%%%%%%%%%%%%%%%%%%%%%%%%%%%%%%%%%%%%
\chapter{Scope Geometry}
\label{ch:scope-geometry}
%%%%%%%%%%%%%%%%%%%%%%%%%%%%%%%%%%%%%%%%%%%%%%%%%%%%%%%%%%%%%%%%%%%%%%%%%%%%%%

\vspace{1em}
\begin{flushright}\itshape
The art of reasoning consists in getting hold of the subject at the right
end, of seizing the few general ideas that illuminate the whole, and of
persistently organizing all subsidiary facts around them.\\
\upshape --- Alfred North Whitehead
\end{flushright}
\vspace{1em}

Volume~I established the philosophical and epistemic foundations of
Structural Semantics. Its principal conclusion was not that phenomenology
is false or unnecessary, but that the mathematical analysis of semantic
organization can proceed independently of questions concerning subjective
experience. Having isolated the structural domain, we now begin
constructing the mathematical objects that will carry the remainder of
the theory.

The central claim of this volume is that semantic organization is
fundamentally geometric. Every computational system, biological or
artificial, operates within collections of constraints that determine
which transformations preserve coherence and which produce contradiction
or collapse. These collections of constraints are neither arbitrary sets
nor merely symbolic contexts: they possess internal organization,
hierarchical nesting, admissible boundaries, and lawful mechanisms for
interaction. This chapter introduces the mathematical object that
captures these properties, the \emph{bounded scope geometry}, which
provides the primitive space on which the Scope Calculus of subsequent
chapters operates. Primitive operations such as Pop, Refuse, Bind, and
Collapse will later be defined as transformations acting on these
geometries (Chapter~7), while Chapters~8--10 establish their categorical,
algebraic, and functorial properties. Throughout this chapter we
deliberately avoid introducing operational rules, defining first the
spaces on which operations act before studying the operations themselves
--- mirroring the development of modern topology and differential
geometry, where spaces are characterized independently of the mappings
subsequently defined between them.

\begin{remark}[Terminological note]
The formalism developed in this and the following four chapters was
originally introduced and published under the name \emph{Spherepop}. The
present exposition is self-contained and does not presuppose that
earlier work; it is renamed and reorganized here as the Scope Calculus to
foreground its role within Structural Semantics rather than its original
programming-language framing. Readers already familiar with Spherepop
will recognize $\Pop$, $\Refuse$, $\Bind$, and $\Collapse$ as the same
four primitives under new emphasis; readers meeting the material for the
first time can treat the two names as synonymous.
\end{remark}

\section{Motivation}

Nearly every formal treatment of semantics assumes, explicitly or
implicitly, the existence of context. Natural language interpretation,
programming languages, logical deduction, theorem proving, biological
regulation, and distributed computation all depend on determining which
objects belong to a given context and which transformations preserve that
membership; yet context is rarely treated as a mathematical object in its
own right. In linguistics it is often represented informally through
discourse structure or lexical fields \cite{Lyons1977,Saussure1916}; in
formal logic it appears as assumptions within derivations
\cite{Gentzen1935}; programming-language semantics typically represents
it through environments, symbol tables, or typing judgments
\cite{Winskel1993,Pierce2002}; and category theory studies morphisms
between objects while generally leaving contextual admissibility to
external structure \cite{MacLane1998}. These approaches have been
extraordinarily successful within their respective domains, but none was
developed specifically to study semantic stability under arbitrary
transformations across multiple representational substrates.

Structural Semantics therefore introduces an explicit geometry of
contexts. Rather than viewing context merely as auxiliary information
attached to a computation, we regard context itself as the primary
mathematical object: semantic transformations become transformations of
context geometries, and semantic preservation becomes the preservation of
admissible geometric structure. This shift is analogous to the historical
transition from Euclidean coordinate calculations to differential
geometry, in which geometry became primary while coordinates became
representations; likewise, scope geometry treats context as primary while
individual representations become coordinates within that context.

\section{Primitive Intuition}

Before introducing formal definitions, it is useful to develop geometric
intuition. Consider a mathematical proof: at any stage there exists a
collection of assumptions, definitions, variables, and established
propositions that remain valid, forming a coherent environment in which
introducing an undefined symbol, violating an assumption, or applying a
theorem outside its domain produces inconsistency. Consider similarly
lexical interpretation: a discussion of planetary motion establishes a
context in which the word ``orbit'' possesses one meaning, while a
discussion of computing establishes a different context in which the same
lexical item acquires an entirely different interpretation, and the
semantic validity of subsequent expressions depends on remaining inside
the appropriate context unless an explicit contextual transition occurs.
Programming languages exhibit the same phenomenon: variables exist only
within declared scopes, references outside these scopes become invalid,
and nested functions inherit some contextual information while
introducing new local constraints.

Although these examples originate in different disciplines, they share a
common structural pattern: each consists of a bounded region of
admissible entities, explicit constraints governing admissibility,
hierarchical nesting, and lawful transitions between nested regions.
Scope Geometry abstracts exactly these common structural features.

\section{Why Boundaries Are Fundamental}

The formal definitions beginning in the next section take a boundary,
not a symbol, a vector, or an entity, as the primitive notion from which
everything else is built. This choice is worth motivating before the
definitions arrive, since it is not the only one available: a theory
could equally well begin with a space of representations and add
constraints afterward, treating admissibility as a derived property
rather than a starting point.

The reason for beginning with boundaries instead is that the examples of
the preceding section share nothing at the level of representation. A
proof's assumptions, a word's discourse context, and a variable's
lexical scope have no common representational format --- one is a set of
propositions, another a topic under discussion, the third a region of
program text. What they share is only the pattern of admissibility
itself: a bounded region within which certain moves are licensed and
others are not, independently of what occupies that region. A theory
that started from representations would need a separate account of
admissibility for every representational format it encountered. A theory
that starts from the boundary can be stated once and instantiated by
whatever representation a given application supplies, which is exactly
the substrate independence pursued formally in Chapter~9.

This also explains why the definitions that follow introduce
$\Boundary$ before $\Interior$ rather than the reverse. The interior of a
scope, Definition~\ref{def:scope-geometry} below, is defined \emph{as}
the states satisfying the boundary, so that the boundary is prior in the
order of definition as well as in the order of motivation: there is no
admissible interior to speak of until the boundary that determines it has
already been specified.

\section{State Domains}

Every computational or cognitive system evolves through collections of
states. We therefore begin with the underlying state domain.

\begin{definition}[State domain]
A \emph{state domain} is a nonempty set $\Domain$ whose elements represent
admissible configurations of a computational, logical, or cognitive
system.
\end{definition}

The internal structure imposed on $\Domain$ depends on the application: for
symbolic reasoning it may consist of logical derivation states, for
transformer architectures $\Domain\subset\mathbb{R}^d$ may represent
activation vectors, and for programming languages it may denote
environments consisting of variable bindings and execution frames. The
theory developed here does not depend on the specific realization of
$\Domain$; only the existence of distinguishable states is assumed.

\section{Boundaries}

State domains alone are insufficient. Semantic organization requires
constraints determining which transitions remain valid; these constraints
are represented by boundaries.

\begin{definition}[Boundary]
Let $\Domain$ be a state domain. A \emph{boundary} $\Boundary$ is a
collection of admissibility predicates $\Boundary=\{\beta_i\}_{i\in I}$,
where each $\beta_i:\Domain\to\{0,1\}$ determines whether a given state
satisfies a particular structural constraint.
\end{definition}

Boundaries need not be purely logical: they may represent semantic
consistency, typing constraints, physical conservation laws, security
permissions, resource limits, or probabilistic admissibility criteria,
and the abstraction intentionally encompasses all such cases. A boundary
in this sense also admits a topological reading, in which it separates
admissible from inadmissible regions of a semantic space in exactly the
sense that a topological boundary separates a region from its
complement; this reading is developed rigorously in Appendix~\ref{app:topology}.

\section{Interior Spaces}

Having introduced boundaries, we define the region they enclose.

\begin{definition}[Interior space]
Given a state domain $\Domain$ and boundary $\Boundary$, the associated
interior space is $\Interior=\{x\in\Domain:\beta(x)=1\text{ for every
}\beta\in\Boundary\}$.
\end{definition}

Thus $\Interior$ contains precisely those states satisfying every active
admissibility constraint; states outside $\Interior$ may still exist within
the global domain $\Domain$, but they are not locally admissible within the
current scope.

\section{Nested Scope Geometries}

Contexts rarely exist in isolation. Most systems construct hierarchies:
programming languages contain nested blocks, logical proofs contain
nested assumptions, human discourse continually opens and closes
subordinate contexts, and transformer attention similarly constructs
localized contextual substructures. To model this hierarchy we introduce
nested scope geometries.

\begin{definition}[Bounded scope geometry]
\label{def:scope-geometry}
A bounded scope geometry is the ordered tuple
$\Scope=\langle\Boundary,\Interior,\Children\rangle$, where $\Boundary$ is a
boundary, $\Interior$ is the associated admissible interior, and $\Children$
is a finite or countable collection of child scope geometries.
\end{definition}

The recursive appearance of $\Children$ permits arbitrarily deep
hierarchical structures while preserving a uniform mathematical
description, and this recursive definition will later permit induction
over scope depth and the derivation of decomposition theorems in
Chapter~9.

\section{Hierarchy, Containment, and Inheritance}

The recursive definition of bounded scope geometries naturally induces a
hierarchical organization: rather than existing as isolated objects,
scopes form partially ordered containment structures in which local
contexts inherit constraints from larger environments while introducing
additional local restrictions. This hierarchical organization appears
across numerous disciplines: block-structured programming languages
inherit variable environments from enclosing functions while introducing
local bindings \cite{Pierce2002}, formal proof systems inherit
assumptions from outer derivations while temporarily extending them with
auxiliary hypotheses \cite{Gentzen1935}, and linguistic discourse
similarly maintains a persistent conversational context while permitting
nested quotations, counterfactuals, and hypothetical reasoning
\cite{Lyons1977,Saussure1916}. The purpose of Scope Geometry is to provide
a common mathematical framework for each of these situations.

\begin{definition}[Parent and child scope]
Let $\Scope=\langle\Boundary,\Interior,\Children\rangle$ be a bounded scope
geometry. If $\Scope'\in\Children$, then $\Scope$ is the \emph{parent
scope} of $\Scope'$ and $\Scope'$ is a \emph{child scope} of $\Scope$.
\end{definition}

Parenthood defines only an immediate containment relation; ancestor and
descendant relationships are obtained by repeated application of the
parent relation.

\begin{definition}[Scope containment]
Let $\Scope_1=\langle\Boundary_1,\Interior_1,{\Children}_1\rangle$ and
$\Scope_2=\langle\Boundary_2,\Interior_2,{\Children}_2\rangle$. We write
$\Scope_2\preceq\Scope_1$ if and only if $\Scope_2$ is contained within
$\Scope_1$.
\end{definition}

The relation $\preceq$ is interpreted as geometric containment rather than
set inclusion: two scopes may contain overlapping semantic information
while remaining distinct objects within the hierarchy, so containment
captures organizational structure rather than membership alone.

A central property of hierarchical contexts is inheritance: child scopes
do not arise independently, but inherit structural constraints
established by their parents unless those constraints are explicitly
modified through admissible transformations.

\begin{definition}[Boundary inheritance]
Let $\Scope_2\preceq\Scope_1$. The boundary of the child scope inherits from
the parent if $\Boundary_1\subseteq\Boundary_2$; additional local
constraints may be introduced, but inherited constraints remain active
unless explicitly discharged.
\end{definition}

Boundary inheritance guarantees that entering a child scope cannot erase
the structural commitments established by its ancestors; instead, deeper
scopes become progressively more constrained. This monotonic refinement
of admissibility will later prove essential for the decomposition
theorems of Chapter~9.

\subsection{Local and Global Admissibility}

The distinction between local and global validity is fundamental: a state
may be perfectly admissible within one scope while violating the
constraints of another. A temporary variable inside a programming block is
valid locally but ceases to exist after leaving the block, and an
assumption introduced during an indirect proof is admissible only within
the subproof in which it was declared. This motivates two complementary
notions.

\begin{definition}[Local admissibility]
A state $x\in\Domain$ is \emph{locally admissible} with respect to a scope
$\Scope$ if $x\in\Interior$.
\end{definition}

\begin{definition}[Global admissibility]
A state is \emph{globally admissible} if it satisfies the boundary
constraints of every ancestor scope containing the current scope.
\end{definition}

Every globally admissible state is locally admissible, but the converse
need not hold.

\subsection{Boundary Violations}

Boundary constraints divide state transitions into two classes: some
preserve admissibility, others produce structural inconsistency.

\begin{definition}[Boundary violation]
Let $x\in\Interior$ and let $T:\Domain\to\Domain$ be a state
transformation. The transformation is a \emph{boundary violation} if
$T(x)\notin\Interior$.
\end{definition}

Boundary violations need not correspond to logical contradictions;
depending on the application they may represent type errors, inconsistent
assumptions, invalid semantic references, resource exhaustion, security
violations, unstable activation trajectories, or loss of contextual
coherence. The Scope Calculus introduced in the next chapter exists
precisely to characterize the permissible responses to such violations.

\subsection{The Scope Hierarchy as a Partially Ordered Structure}

The containment relation naturally equips the collection of scopes with a
partial order.

\begin{proposition}
The relation $\preceq$ over bounded scope geometries is reflexive,
antisymmetric, and transitive.
\end{proposition}

\begin{proof}
Reflexivity follows immediately since every scope contains itself.
Antisymmetry follows because if $\Scope_1\preceq\Scope_2$ and
$\Scope_2\preceq\Scope_1$, then both scopes possess identical containment
relationships and are identified as the same scope geometry. Transitivity
follows because containment through an intermediate scope implies
containment within every ancestor. Hence $(\CatScope,\preceq)$ forms a
partially ordered set.
\end{proof}

This order will later support induction over scope depth, recursive
repair procedures, and the categorical constructions of Chapters~8--10.

\section{Discussion}

The geometry introduced thus far deliberately remains static: we have
defined the spaces on which semantic organization is represented, but not
yet described how those spaces evolve. This distinction mirrors the
historical development of geometry itself --- classical topology first
characterizes spaces and only subsequently studies continuous mappings
between them, and differential geometry first defines manifolds before
introducing vector fields and flows. Scope Geometry follows the same
methodological order: Chapters~6 and~7 separate geometric structure from
dynamical evolution, so that once the spaces have been rigorously defined,
transformations may be introduced without ambiguity. In particular, the
four primitive operations of the Scope Calculus will be shown to generate
every admissible structural transformation required by the framework,
providing the dynamical counterpart to the static geometries defined here.

\section{Chapter Summary}

This chapter introduced the fundamental geometric objects underlying
Structural Semantics. Beginning with state domains and boundary
constraints, we constructed bounded scope geometries as recursively
nested contexts possessing admissible interiors and hierarchical
organization, defined parent-child relationships, boundary inheritance,
local and global admissibility, and established that the collection of
scope geometries forms a partially ordered hierarchy under containment.
These constructions intentionally remain independent of any operational
mechanism, defining the spaces on which semantic organization is
represented without prescribing how those spaces evolve. The next chapter
introduces the minimal operational language required to manipulate these
geometries: the four primitive operations Pop, Refuse, Bind, and Collapse.

\subsection*{Status of Results: Chapter 6}

\paragraph{Established background.} Hierarchical contexts in formal
logic, programming-language semantics, topology, and lexical semantics
\cite{Gentzen1935,Pierce2002,Lyons1977,MacLane1998}.

\paragraph{Formal developments.} Definitions of parent and child scopes,
scope containment, boundary inheritance, local and global admissibility,
boundary violations, and the partial ordering of scope hierarchies.

\paragraph{Empirical hypotheses.} None. This chapter establishes the
static geometric framework on which the operational Scope Calculus will
be developed in Chapter~7.

%%%%%%%%%%%%%%%%%%%%%%%%%%%%%%%%%%%%%%%%%%%%%%%%%%%%%%%%%%%%%%%%%%%%%%%%%%%%%%
\chapter{Primitive Operations and the Scope Calculus}
\label{ch:scope-calculus}
%%%%%%%%%%%%%%%%%%%%%%%%%%%%%%%%%%%%%%%%%%%%%%%%%%%%%%%%%%%%%%%%%%%%%%%%%%%%%%

\vspace{1em}
\begin{flushright}\itshape
Mathematics possesses not only truth, but supreme beauty --- a beauty cold
and austere, like that of sculpture.\\
\upshape --- Bertrand Russell
\end{flushright}
\vspace{1em}

Chapter~6 introduced the static geometry underlying Structural Semantics:
bounded scope geometries, boundary constraints, hierarchical containment,
and admissibility relations, without specifying how these objects evolve.
The present chapter supplies that missing component by introducing the
primitive operations that generate all admissible transformations within
Scope Geometry.

The guiding principle is one of structural minimality: rather than
constructing a large collection of specialized operators, we seek the
smallest operational alphabet capable of expressing every admissible
transformation required by semantic organization, paralleling how group
theory seeks minimal generating sets, formal language theory studies
finite alphabets generating infinite languages, and universal algebra
investigates operations sufficient to construct entire classes of
algebraic structures \cite{MacLane1998,Hopcroft2006,Burris1981}. The claim
advanced in this chapter is that four primitive operations,
$\ScopeAlphabet$, are sufficient to describe the evolution of bounded
scope geometries. These primitives are defined independently of any
particular computational substrate: they are not programming-language
instructions, neural-network operations, or logical inference rules, but
abstract transformations of context itself. Subsequent chapters show that
this primitive alphabet admits a rich algebraic structure, supports
canonical normal forms, and provides the foundation for categorical and
topological analysis.

\section{Motivation}

Any sufficiently expressive computational system must introduce new
contextual information, reject information incompatible with existing
constraints, establish relationships between previously independent
entities, and eventually resolve temporary structures into stable
results. Programming languages accomplish these tasks through scope
creation, exception handling, variable binding, and function evaluation;
mathematical proofs introduce assumptions, reject contradictions,
construct dependencies, and discharge temporary hypotheses; conversation
continually opens new topics, rejects inappropriate interpretations,
links ideas, and eventually resolves discussion into shared conclusions.
Although the implementations differ dramatically, the underlying
structural transformations remain remarkably similar, and the Scope
Calculus abstracts these transformations into four primitive operations
acting directly on bounded scope geometries.

\section{Primitive Alphabet}

\begin{definition}[Primitive alphabet]
The primitive operational alphabet of the Scope Calculus is
$\ScopeAlphabet$. Each primitive denotes a partial transformation acting
on bounded scope geometries.
\end{definition}

The primitives are interpreted operationally rather than symbolically:
they describe transformations of context rather than transformations of
syntax, and every admissible computation considered throughout this
monograph is represented as a finite composition of these primitive
operations. The remainder of this chapter examines each primitive
individually before constructing their algebraic interactions.

\subsection{The Pop Operator}

The first primitive governs contextual escalation. Many computations
require information generated within a local scope to become available
within an enclosing scope --- returning a function value, proving an
intermediate lemma, exporting a variable, or promoting a local conclusion
into a global argument.

\begin{definition}[Pop]
Let $\Scope_c\preceq\Scope_p$ be a child scope contained within a parent
scope. The Pop operator $\Pop:\Scope_c\to\Scope_p$ promotes an admissible
entity from the child scope into its parent while preserving all
inherited boundary constraints.
\end{definition}

Pop does not destroy the child scope; it merely transfers designated
structural information across the boundary between nested contexts. The
admissibility of such promotion depends entirely on the boundary
constraints of Chapter~6: if promotion violates an inherited boundary, the
operation is not admissible. The interaction between Pop and Refuse is
examined later in this chapter.

\subsection{The Refuse Operator}

The second primitive governs boundary enforcement. Every coherent context
must distinguish admissible transformations from inadmissible ones,
without which semantic organization degenerates into unrestricted
accumulation of incompatible information.

\begin{definition}[Refuse]
Let $x\in\Domain$ be a proposed state transition. The Refuse operator
$\Refuse$ rejects the transition whenever $x\notin\Interior$; equivalently,
$\Refuse$ acts whenever the proposed transformation violates one or more
active boundary predicates.
\end{definition}

Refuse is not an exceptional condition added to the calculus, but a
primitive structural mechanism preserving admissibility. In programming
languages it corresponds to type violations or access restrictions, in
logical deduction to invalid inference, in conversational reasoning to
rejecting incompatible interpretations, and in the transformer
diagnostics of Chapters~13--15 it will be interpreted as the detection of
structural context collapse rather than low statistical confidence.

\subsection{The Bind Operator}

Whereas Pop transfers existing entities between contexts, Bind creates new
structural relationships within a context.

\begin{definition}[Bind]
Let $a,b\in\Interior$. The Bind operator $\Bind$ creates an admissible
dependency relation $a\leftrightarrow b$ within the current scope
geometry.
\end{definition}

The precise mathematical nature of this dependency depends on the
application: within programming languages it may represent variable
association, within logical systems dependence between hypotheses and
conclusions, and within neural architectures it may later be interpreted
as persistent activation coupling. The calculus intentionally abstracts
over these realizations, concerning itself with the creation of lawful
structural dependence rather than its physical implementation. Readers
who know the $\lambda$-calculus may already suspect a connection between
$\Bind$ and variable binding; that connection is made precise, rather
than left as an analogy, in Appendix~\ref{app:lambda}.

\subsection{The Collapse Operator}

The final primitive governs contextual resolution. Temporary structures
cannot remain indefinitely unresolved: at some stage subordinate
computations terminate and their results become available to surrounding
contexts.

\begin{definition}[Collapse]
Let $\Scope=\langle\Boundary,\Interior,\Children\rangle$. The Collapse
operator $\Collapse$ maps the current scope geometry onto a canonical
admissible result while discharging temporary internal structure.
\end{definition}

Collapse should not be interpreted as destruction but as evaluation:
temporary internal organization is summarized by a stable outcome that
can subsequently participate in larger computations, such as evaluation
of an expression, completion of a proof, termination of recursion,
resolution of a dialogue, or completion of an optimization step.

\subsection{Operational Completeness}

\begin{theorem}[Operational completeness]
\label{thm:operational-completeness}
Every admissible transformation of bounded scope geometries admits a
finite presentation as a composition of the primitive operations
$\{\Pop,\Refuse,\Bind,\Collapse\}$.
\end{theorem}

\begin{proof}
Proceed by structural induction over admissible context transformations.
Every transformation either introduces contextual information, rejects
inadmissible information, creates admissible dependency, or resolves
temporary structure; complex transformations decompose recursively into
finite compositions of these elementary cases. The complete inductive
argument is deferred to Appendix~\ref{app:proofs}.
\end{proof}

\section{Composition of Primitive Operations}

The primitive operators acquire their mathematical significance only
through composition: individual operations describe isolated contextual
events, while computation, reasoning, and semantic interpretation arise
from finite sequences of such events. Unlike arithmetic operations, the
primitive transformations of the Scope Calculus are inherently
contextual, so the algebra they generate is generally non-commutative.

\begin{definition}[Primitive composition]
Let $\alpha,\beta\in\ScopeAlphabet$. The composition $\beta\circ\alpha$
denotes sequential execution of $\alpha$ followed by $\beta$, provided the
intermediate scope geometry remains admissible.
\end{definition}

Whenever admissibility fails after the first operation, the composed
transformation is undefined, so the Scope Calculus naturally forms a
partial algebra rather than a total one.

A finite computation over a bounded scope geometry $\Scope_0$ is
represented by an ordered sequence
$\ScopeMorph=\alpha_n\circ\cdots\circ\alpha_1$ with $\alpha_i\in\ScopeAlphabet$,
executed right to left and generating an evolution
$\Scope_0\to\Scope_1\to\cdots\to\Scope_n$ in which each transition
satisfies the admissibility conditions of Chapter~6.

\begin{definition}[Admissible computation]
A finite primitive sequence $\ScopeMorph=\alpha_n\circ\cdots\circ\alpha_1$ is
an \emph{admissible computation} if every intermediate transition
$\Scope_i\to\Scope_{i+1}$ preserves local boundary admissibility.
\end{definition}

Admissibility is therefore a property of the entire computation rather
than merely its initial and final states.

\begin{proposition}[Associativity]
Let $\alpha,\beta,\gamma\in\ScopeAlphabet$. Whenever both sides are
defined, $(\gamma\circ\beta)\circ\alpha=\gamma\circ(\beta\circ\alpha)$.
\end{proposition}

\begin{proof}
Primitive transformations are partial functions acting on bounded scope
geometries, so associativity follows directly from associativity of
functional composition, the only restriction being that every
intermediate transformation remains admissible.
\end{proof}

Associativity permits omission of parentheses in finite primitive
sequences; henceforth $\alpha_1\circ\alpha_2\circ\cdots\circ\alpha_n$
denotes the unique associative composition.

Although composition is associative, it is generally not commutative: the
order in which contextual transformations occur affects the resulting
scope geometry, since binding an entity before it exists differs
fundamentally from binding it after introduction.

\begin{proposition}[Non-commutativity]
There exist $\alpha,\beta\in\ScopeAlphabet$ such that
$\alpha\circ\beta\neq\beta\circ\alpha$.
\end{proposition}

\begin{proof}
Consider a local entity $x\in\Interior$. Applying $\Bind$ followed by $\Pop$
first establishes a dependency before promoting the resulting structure;
reversing the order promotes the entity prior to creating the dependency.
Unless the parent scope already contains the corresponding relational
structure, the resulting scope geometries differ, so
$\Pop\circ\Bind\neq\Bind\circ\Pop$.
\end{proof}

Non-commutativity is therefore an intrinsic consequence of contextual
organization rather than an implementation artifact.

\begin{definition}[Identity operation]
The identity operation $\mathrm{id}$ is the unique primitive sequence of
length zero satisfying $\mathrm{id}(\Scope)=\Scope$ for every bounded scope
geometry.
\end{definition}

Consequently $\ScopeMorph\circ\mathrm{id}=\ScopeMorph=\mathrm{id}\circ\ScopeMorph$;
this identity transformation later becomes the identity morphism of the
category $\CatScope$ introduced in Chapter~8.

The four primitives interact in characteristic ways: some compositions
are always admissible, others require boundary verification, and still
others are undefined whenever structural constraints are violated. Table
\ref{tab:primitive-interaction} summarizes these interactions
informally; the precise rewrite rules follow in the next section.

\begin{table}[h]
\centering
\begin{tabular}{lll}
\toprule
First & Second & Typical interpretation \\
\midrule
$\Bind$ & $\Pop$ & Promote established relation \\
$\Pop$ & $\Bind$ & Bind after promotion \\
$\Refuse$ & $\Pop$ & Promotion denied \\
$\Pop$ & $\Refuse$ & Boundary validation after promotion \\
$\Collapse$ & $\Pop$ & Promote evaluated result \\
$\Collapse$ & $\Bind$ & Bind canonical representative \\
$\Bind$ & $\Collapse$ & Evaluate constructed dependency \\
\bottomrule
\end{tabular}
\caption{Typical interactions among primitive operators. Formal rewrite
rules follow in the next section.}
\label{tab:primitive-interaction}
\end{table}

\section{The Primitive Rewriting System}

We now equip the algebra with a rewriting system that identifies
operationally equivalent computations. Its purpose is twofold: first, to
eliminate redundant or vacuous transformations, yielding shorter and more
canonical descriptions of computations; second, to provide the foundation
for the normal-form results of Chapter~9. The rewriting system is
deliberately conservative --- it does not alter the semantic content of a
computation, but replaces one admissible presentation with another that
is structurally simpler while preserving boundary constraints and
contextual relationships. The elementary treatment given here is
sufficient for the results of this chapter; the general theory of
abstract reduction systems that this construction is an instance of,
including the full termination and confluence arguments, is developed
independently in Appendix~\ref{app:reduction}.

\begin{definition}[Primitive word]
A \emph{primitive word} is a finite sequence
$w=\alpha_1\alpha_2\cdots\alpha_n$ with $\alpha_i\in\ScopeAlphabet$. The
empty word is denoted $\varepsilon$.
\end{definition}

The collection of all primitive words forms the free monoid
$\ScopeAlphabet^{*}$ under concatenation, containing every syntactically
valid computation, including many that are operationally redundant or
semantically equivalent; the purpose of rewriting is to identify these
equivalent representations.

\begin{definition}[Rewrite relation]
A \emph{rewrite relation} $\longrightarrow$ is a binary relation on
$\ScopeAlphabet^{*}$ such that $u\longrightarrow v$ whenever the primitive
word $u$ may be replaced by the simpler word $v$ without altering its
admissible structural action on bounded scope geometries.
\end{definition}

Repeated application of rewrite rules produces a reduction sequence
$w\longrightarrow w_1\longrightarrow w_2\longrightarrow\cdots$; a word is
\emph{irreducible} if no rewrite rule applies. The primitive algebra
admits four elementary reductions, motivated by operational redundancy
rather than arbitrary symbolic simplification: identity elimination,
$w\varepsilon\longrightarrow w$ and $\varepsilon w\longrightarrow w$ (R1);
consecutive refusal, $\Refuse\Refuse\longrightarrow\Refuse$ (R2), since
repeated refusals without an intervening admissible transformation are
equivalent to a single refusal; consecutive collapse,
$\Collapse\Collapse\longrightarrow\Collapse$ (R3), since further collapse
operations after a scope has already been collapsed into its canonical
representative produce no additional effect; and empty promotion,
$\Pop\varepsilon\longrightarrow\varepsilon$ (R4), since promotion of an
empty computation is redundant. These elementary rules intentionally
avoid introducing semantic assumptions regarding application-specific
behavior; more specialized rewrite rules may be added in concrete
implementations provided they preserve admissibility.

\begin{definition}[Operational equivalence]
Two primitive words $u,v\in\ScopeAlphabet^{*}$ are \emph{operationally
equivalent}, written $u\equiv v$, if one can be transformed into the other
through a finite sequence of rewrite rules and inverse rewrite rules.
\end{definition}

Operational equivalence partitions the free monoid into equivalence
classes whose members describe identical admissible computations, so the
Scope Calculus distinguishes between the syntactic form of a computation
and its structural effect.

\begin{lemma}[Termination of elementary rewriting]
Every finite sequence of elementary rewrite rules terminates after a
finite number of reductions.
\end{lemma}

\begin{proof}
Each elementary rewrite rule strictly decreases either the length of the
primitive word or the number of redundant operators it contains. Since
every primitive word has finite length, no infinite descending chain
exists, so every reduction sequence terminates.
\end{proof}

Termination alone is insufficient: different rewrite orders must also
lead to compatible results.

\begin{definition}[Local confluence]
A rewriting system is \emph{locally confluent} if whenever
$w\longrightarrow u$ and $w\longrightarrow v$, there exists a primitive
word $z$ such that $u\longrightarrow^{*}z$ and $v\longrightarrow^{*}z$.
\end{definition}

For the elementary rewriting rules above, overlapping rewrite patterns
occur only in a finite number of cases, each resolving consistently; the
detailed verification appears in Appendix~\ref{app:proofs}. Termination and local
confluence together motivate the search for unique canonical
representatives of operational equivalence classes; rather than proving
uniqueness immediately, we postpone the complete argument to Chapter~9,
where the Primitive Normal Form Theorem is established after the
algebraic and categorical foundations have been completed. For the
remainder of this chapter a primitive word is called \emph{reduced} if no
elementary rewrite rule applies; every admissible computation considered
in later chapters may be represented by one or more primitive words, each
reducible to a common reduced representative.

\section{Algebraic Properties of the Scope Calculus}

Having introduced the primitive alphabet and its rewriting system, we now
examine the algebraic structure it generates. Unlike many familiar
algebraic systems, the Scope Calculus is neither a group nor a ring:
primitive operations are generally irreversible, context-dependent, and
only partially defined. Nevertheless they exhibit a rich collection of
structural regularities sufficient to support canonical reduction,
categorical interpretation, and representation independence.

\begin{definition}[Scope monoid]
The pair $(\ScopeAlphabet^{*},\circ)$ consisting of all finite primitive
words together with concatenation is called the \emph{Scope Monoid}.
\end{definition}

\begin{proposition}
The Scope Monoid is closed under concatenation, associative, and possesses
the identity element $\varepsilon$.
\end{proposition}

\begin{proof}
Closure follows because concatenation of two finite primitive words
produces another finite primitive word; associativity follows from
concatenation of finite sequences; and $\varepsilon$ acts as both left and
right identity.
\end{proof}

This monoid is intentionally syntactic; operational semantics arise only
after admissibility constraints and rewriting relations are introduced.

\begin{definition}[Operational quotient]
Let $\equiv$ denote operational equivalence under the rewrite system. The
quotient algebra $\mathcal{Q}=\ScopeAlphabet^{*}/\equiv$ is the \emph{Scope
Quotient Algebra}.
\end{definition}

Elements of $\mathcal{Q}$ are equivalence classes of primitive
computations rather than individual primitive words, so the quotient
algebra captures computational behavior instead of syntactic
representation.

\begin{definition}[Partial composition]
Let $u,v\in\ScopeAlphabet^{*}$. The composition $v\circ u$ is defined only
if execution of $u$ terminates in a scope geometry satisfying the
admissibility conditions required for the initial application of $v$.
\end{definition}

Thus the Scope Calculus forms a partial algebra over admissible
computations; this partiality is not an inconvenience but one of the
defining features of the theory, since invalid compositions correspond
precisely to structural violations of context.

\begin{proposition}
Under the elementary rewrite rules, $\Refuse^2=\Refuse$ and
$\Collapse^2=\Collapse$.
\end{proposition}

\begin{proof}
Both identities follow immediately from rewrite rules (R2) and (R3):
repeated refusal introduces no additional structural restriction beyond
the first refusal, and once a scope has been collapsed into its canonical
representative, further collapse operations leave the resulting geometry
unchanged.
\end{proof}

These identities illustrate that certain operations are naturally
stabilizing, with repeated application eventually reaching a fixed point.

\begin{proposition}
The Scope Calculus is not, in general, a group.
\end{proposition}

\begin{proof}
Suppose $\Collapse$ possessed an inverse; then every collapsed computation
could be reconstructed uniquely into its complete internal history.
However, Collapse intentionally replaces temporary internal structure by
a canonical representative, so information regarding intermediate
computation is generally lost, and $\Collapse$ admits no universal
inverse. An identical argument applies to $\Refuse$, whose rejection of
inadmissible transitions does not uniquely determine the rejected
proposal. Hence the primitive algebra cannot satisfy the group axioms.
\end{proof}

This non-invertibility distinguishes Scope Geometry from reversible
computation while aligning naturally with evaluation, abstraction, and
proof construction.

\begin{definition}[Operational fixed point]
A bounded scope geometry $\Scope$ is a \emph{fixed point} of a primitive
transformation $\alpha$ if $\alpha(\Scope)=\Scope$.
\end{definition}

Examples include a fully evaluated computation under $\Collapse$, a
boundary already enforcing every relevant admissibility constraint under
$\Refuse$, and an established dependency under repeated $\Bind$ where no
additional relations are introduced. Fixed points provide the algebraic
foundation for convergence arguments developed in Chapter~9 and later
become central to the diagnostic theory of stable semantic trajectories
in Chapters~13 and~15.

\section{Discussion}

The algebra introduced in this chapter deliberately occupies a middle
ground between abstract algebra and operational semantics. On one hand,
the primitive operators generate an algebraic structure supporting
composition, quotient construction, and canonical reduction; on the
other, the algebra remains closely connected to concrete computational
behavior through admissibility constraints and contextual boundaries.
This dual character distinguishes the Scope Calculus from traditional
rewriting systems: rather than manipulating symbols in isolation, every
rewrite reflects a structural transformation of bounded semantic
contexts. The following chapter abstracts these operational ideas one
level further, examining the category whose objects are bounded scope
geometries and whose morphisms are admissible context-preserving
transformations; the algebra developed here becomes the operational
substrate for the categorical constructions of Chapter~8.

\section{Chapter Summary}

This chapter completed the construction of the Scope Calculus. Beginning
with the four primitive operations Pop, Refuse, Bind, and Collapse, we
developed their composition laws, established the elementary rewriting
system, introduced operational equivalence, and constructed the quotient
algebra governing admissible computations, showing that the primitive
alphabet generates a partial monoid, identifying idempotent operators,
proving the absence of global inverses, and introducing operational fixed
points. These results provide the complete algebraic machinery required
for the categorical development that follows, in which Chapter~8 elevates
these operational constructions into the category $\CatScope$.

\subsection*{Status of Results: Chapter 7}

\paragraph{Established background.} Free monoids, rewriting systems,
quotient algebras, partial algebras, and operational semantics
\cite{Hopcroft2006,Baader1998,MacLane1998}.

\paragraph{Formal developments.} Definition of the Scope Monoid, Scope
Quotient Algebra, partial composition, operational equivalence,
idempotence, non-invertibility, and operational fixed points.

\paragraph{Empirical hypotheses.} None. Chapter~7 develops the algebraic
machinery underlying the Scope Calculus; empirical application begins
only in Volume~III after the categorical and topological foundations have
been completed.

%%%%%%%%%%%%%%%%%%%%%%%%%%%%%%%%%%%%%%%%%%%%%%%%%%%%%%%%%%%%%%%%%%%%%%%%%%%%%%
\chapter{Category-Theoretic Foundations of Scope Geometry}
\label{ch:category}
%%%%%%%%%%%%%%%%%%%%%%%%%%%%%%%%%%%%%%%%%%%%%%%%%%%%%%%%%%%%%%%%%%%%%%%%%%%%%%

\vspace{1em}
\begin{flushright}\itshape
Mathematics is the art of giving the same name to different things.\\
\upshape --- Henri Poincar\'e
\end{flushright}
\vspace{1em}

Category theory was introduced to mathematics by Eilenberg and Mac Lane as
a language for describing mathematical structure independently of any
particular realization \cite{Eilenberg1945,MacLane1998}, and during the
twentieth century it became one of the central organizing frameworks of
modern algebra, topology, geometry, logic, and computer science. The
motivation for introducing it into Structural Semantics is considerably
more modest. The previous chapters established two complementary
viewpoints: bounded scopes as mathematical objects possessing boundaries,
interiors, and nested child scopes, and admissible computations as finite
compositions of primitive operations governed by the Scope Calculus.
These developments naturally suggest a higher level of abstraction, in
which we study the transformations between scope geometries rather than
reasoning directly about primitive words or individual scope objects, and
category theory provides precisely the mathematical language required for
this transition.

Throughout this chapter we adopt a principle of \emph{structural
minimalism}: only those categorical constructions required by subsequent
chapters are introduced. We do not assume that the resulting category is
monoidal, cartesian closed, complete, cocomplete, or possesses universal
limits; such structures may eventually arise as derived properties, but
they are never assumed a priori.

\section{Motivation}

Many mathematical theories begin by studying individual objects before
discovering that the important information lies not within isolated
objects but in the relationships between them: linear algebra studies
vector spaces together with linear maps, topology studies spaces together
with continuous maps, and group theory studies groups together with
homomorphisms. Likewise, Structural Semantics ultimately studies scope
geometries together with admissible context-preserving transformations.
This shift of emphasis is essential, since a semantic system rarely
remains static --- programs execute, proofs evolve, conversations
accumulate context, reasoning systems revise hypotheses, and large
language models continuously transform latent representations from one
layer to the next --- so semantic organization should be understood
primarily as a family of transformations rather than as isolated states.
The purpose of category theory is therefore not to replace the Scope
Calculus, but to provide an abstract language describing the algebra
constructed in Chapter~7.

\begin{definition}[Structural minimalism principle]
A categorical construction shall be introduced only if every required
axiom is explicitly verified, the construction is necessary for
subsequent mathematical development, and no stronger categorical
structure is assumed than is required.
\end{definition}

This principle ensures that the categorical framework grows naturally
from the algebra rather than being imposed externally.

\section{Objects}

The objects of the category are precisely the bounded scope geometries
introduced earlier, reinterpreted from the categorical perspective.

\begin{definition}[Objects of $\CatScope$]
The objects of $\CatScope$ are all well-formed bounded scope geometries
$\Scope=\langle\Boundary,\Interior,\Children\rangle$ satisfying the
admissibility conditions of Chapter~6.
\end{definition}

No assumptions are made concerning the internal structure of the state
space: the interior may consist of symbolic expressions, graphs, neural
activations, logical formulas, computational processes, or any other
representation satisfying the boundary conditions, so the objects of
$\CatScope$ are intentionally substrate independent.

\section{Morphisms}

The morphisms of the category represent admissible semantic
transformations. Operationally these correspond to computations generated
by the primitive alphabet; abstractly they are structure-preserving maps
between bounded scope geometries.

\begin{definition}[Scope morphism]
Let $\Scope_1,\Scope_2\in\Obj(\CatScope)$. A morphism
$\ScopeMorph:\Scope_1\to\Scope_2$ is an admissible transformation preserving
the boundary relations of the source scope throughout execution.
\end{definition}

This definition is deliberately abstract: a morphism is characterized by
what it preserves, not by how it is presented syntactically. The
connection to the primitive alphabet of Chapter~7 is then a fact to be
proved rather than a feature built into the definition, which keeps the
category free to accommodate other presentations of the same
transformations later.

\begin{proposition}[Primitive presentation]
\label{prop:primitive-presentation}
Every scope morphism $\ScopeMorph\in\Hom(\Scope_1,\Scope_2)$ admits a
presentation as a finite sequence of primitive operations drawn from
$\ScopeAlphabet$.
\end{proposition}

\begin{proof}
By Definition~\ref{def:scope-geometry}, the structural difference between
$\Scope_1$ and $\Scope_2$ is confined to finite modifications of
$\Boundary$, $\Interior$, and $\Children$. Theorem~\ref{thm:operational-completeness}
established that every admissible transformation of a bounded scope
geometry decomposes into a finite composition of $\Pop$, $\Refuse$,
$\Bind$, and $\Collapse$; applying this to each of the finitely many
modifications yields a finite primitive presentation of $\ScopeMorph$.
\end{proof}

Morphisms need not preserve every internal detail: temporary
computational states may change dramatically, and only admissibility
constraints are required to remain valid throughout the transformation.
This observation foreshadows the Representation Invariance Theorem of
Chapter~9. The categorical structure introduced in this chapter is kept
deliberately minimal; the stronger question of whether $\CatScope$
supports the full Cartesian closed structure needed to interpret typed
$\lambda$-terms is addressed, under explicit hypotheses, in
Appendix~\ref{app:ccc}.

\begin{definition}[Identity morphism]
For every bounded scope geometry $\Scope$ there exists an identity
morphism $\operatorname{id}_{\Scope}:\Scope\to\Scope$ defined by
$\operatorname{id}_{\Scope}(x)=x$ for every admissible state $x\in\Interior$.
\end{definition}

Operationally, the identity morphism corresponds to the empty primitive
computation of Chapter~7: no boundaries change, no scope relationships
are modified, and no computational history is generated. Identity
therefore represents structural persistence rather than computational
inactivity.

\begin{definition}[Composition]
Suppose $\ScopeMorph_1:\Scope_1\to\Scope_2$ and
$\ScopeMorph_2:\Scope_2\to\Scope_3$. Their composition
$\ScopeMorph_2\circ\ScopeMorph_1:\Scope_1\to\Scope_3$ is defined whenever the
codomain of $\ScopeMorph_1$ equals the domain of $\ScopeMorph_2$.
\end{definition}

This definition mirrors ordinary functional composition; operationally it
represents the execution of one admissible context transformation
immediately after another while preserving the boundary constraints
inherited from intermediate scopes.

\begin{theorem}
The pair $(\Obj(\CatScope),\Hom(\CatScope))$ forms a category.
\end{theorem}

\begin{proof}
Identity morphisms exist by construction, and composition is defined
whenever codomain and domain agree. Associativity follows immediately
from associativity of sequential composition established for the Scope
Calculus in Chapter~7. Finally,
$\operatorname{id}_{\Scope_2}\circ\ScopeMorph=\ScopeMorph=\ScopeMorph\circ\operatorname{id}_{\Scope_1}$
for every admissible morphism $\ScopeMorph:\Scope_1\to\Scope_2$. Therefore all
category axioms are satisfied.
\end{proof}

The importance of this theorem is conceptual rather than computational:
it establishes that semantic transformations possess exactly the
structural regularity required for categorical reasoning.

\section{Interpretation}

The category $\CatScope$ should not be interpreted as describing
cognition, consciousness, or intelligence; instead it provides an
abstract mathematical language for discussing families of admissible
transformations independently of their concrete implementation. Whether a
computation is carried out by symbolic rewriting, biological neurons,
silicon processors, or distributed agents is irrelevant at the
categorical level; only the preservation of admissible scope structure
remains visible. This abstraction is precisely what later permits
comparison between Whiteheadian process philosophy and Scope Geometry
through a translation functor developed in Chapter~10. Readers wondering
whether the single notion of composition used throughout this chapter is
the only one Scope Geometry supports may wish to know that it is not: a
second, independent notion of composition, arising from refining a scope
rather than executing it, is developed as a conditional extension in
Appendix~\ref{app:higher}, though nothing in the chapters that follow
depends on it.

\section{Chapter Summary}

This chapter introduced the category $\CatScope$ using the principle of
structural minimalism. Objects were defined as bounded scope geometries,
morphisms as admissible context-preserving transformations, and
sequential execution was formalized through categorical composition.
Verification of the category axioms established that Scope Geometry
possesses a mathematically well-defined categorical structure while
avoiding unnecessary assumptions concerning monoidal or higher
categorical properties. The next chapter develops the principal
structural theorems governing this category: closure, decomposition,
canonical normal forms, admissibility preservation, reconstruction, and
the Representation Invariance Theorem.

\subsection*{Status of Results: Chapter 8}

\paragraph{Established background.} Elementary category theory: objects,
morphisms, identity morphisms, and composition
\cite{Eilenberg1945,MacLane1998,Awodey2010}.

\paragraph{Formal developments.} Definition of the category $\CatScope$,
the Structural Minimalism Principle, categorical interpretation of
bounded scope geometries, admissible morphisms, and verification of the
category axioms.

\paragraph{Empirical hypotheses.} None. This chapter develops only the
abstract mathematical framework required for subsequent structural
theorems.

%%%%%%%%%%%%%%%%%%%%%%%%%%%%%%%%%%%%%%%%%%%%%%%%%%%%%%%%%%%%%%%%%%%%%%%%%%%%%%
\chapter{Foundational Theorems of Structural Semantics}
\label{ch:foundational-theorems}
%%%%%%%%%%%%%%%%%%%%%%%%%%%%%%%%%%%%%%%%%%%%%%%%%%%%%%%%%%%%%%%%%%%%%%%%%%%%%%

\vspace{1em}
\begin{flushright}\itshape
The essence of mathematics lies in its freedom.\\
\upshape --- Georg Cantor
\end{flushright}
\vspace{1em}

The preceding chapters introduced the fundamental components of the Scope
Calculus: bounded scope geometries as the primary semantic objects,
primitive operations as generators of admissible computations, and
Chapter~8 demonstrated that these objects and transformations naturally
form the category $\CatScope$. The present chapter instead establishes the
principal structural results on which the remainder of the monograph
depends, providing the internal consistency of the Scope Calculus and
justifying its interpretation as a framework for describing semantic
organization independently of any particular computational substrate.
Several proofs are necessarily lengthy; to preserve continuity of the
main exposition, complete technical proofs are deferred to Appendix~\ref{app:proofs},
and only the principal arguments are presented here.

\section{Motivation}

Every mathematical theory requires internal guarantees: linear algebra
requires that matrix multiplication be associative, topology requires
continuity to behave well under composition, and group theory requires
closure under multiplication. Structural Semantics requires analogous
guarantees, without which the primitive operations of Chapter~7 would
remain merely descriptive rather than forming a coherent mathematical
system. The goals of this chapter are therefore to demonstrate that
admissible computations remain within the space of admissible scope
geometries; that every admissible computation possesses a canonical
decomposition; that the rewriting system converges to unique normal
forms; that semantic preservation can be characterized in terms of
admissibility; and finally to establish the Representation Invariance
Theorem, which provides the mathematical basis for substrate independence.

\section{Preliminary Assumptions}

Throughout this chapter we assume the framework established previously:
bounded scope geometries as in Definition~\ref{def:scope-geometry},
primitive operations satisfying the operational rules of Chapter~7,
admissible morphisms forming the category $\CatScope$ of Chapter~8, and
finiteness of every primitive computation. The finiteness assumption is
important: without it, uniqueness of normal forms and termination of
rewriting would generally require additional hypotheses concerning
infinitary rewriting systems.

\section{Closure of Scope Geometry}

\begin{theorem}[Closure theorem]
\label{thm:closure}
Let $\Scope\in\Obj(\CatScope)$ and let $\ScopeMorph:\Scope\to\Scope'$ be an
admissible morphism. Then $\Scope'\in\Obj(\CatScope)$.
\end{theorem}

\begin{proof}
Every admissible morphism preserves the boundary relations defining the
target scope. Because admissibility requires every primitive operation to
satisfy the boundary tensor $\Boundary$, the resulting object satisfies the
same well-formedness conditions required by
Definition~\ref{def:scope-geometry}, so the codomain remains a bounded
scope geometry. A complete inductive proof appears in Appendix~\ref{app:proofs}.
\end{proof}

The Closure Theorem guarantees that semantic computation cannot escape
the mathematical universe of discourse, analogous to closure properties
throughout algebra.

\section{Canonical Primitive Decomposition}

\begin{theorem}[Scope decomposition]
\label{thm:decomposition}
Every admissible morphism $\ScopeMorph:\Scope_1\to\Scope_2$ admits a finite
decomposition into primitive operations drawn from $\ScopeAlphabet$.
\end{theorem}

\begin{proof}
Proceed by structural induction on the length of admissible computations.
The base case consists of identity morphisms. The inductive step follows
because every admissible transformation changes only one of $\Boundary$,
$\Interior$, or $\Children$, and Chapter~7 established that each such
modification is generated by one of the primitive operations. Hence every
admissible computation possesses a primitive decomposition.
\end{proof}

The significance of this theorem extends beyond implementation: it states
that the four primitive operations are operationally complete with
respect to admissible semantic transformations.

\section{Canonical Normal Forms}

Primitive decompositions are generally not unique, since different
computational histories may represent the same semantic transformation.
The next theorem resolves this ambiguity.

\begin{theorem}[Primitive normal form]
\label{thm:normal}
Every admissible primitive sequence reduces, through the rewriting rules
of Chapter~7, to a unique canonical normal form $\ScopeMorph^*$.
\end{theorem}

\begin{proof}
Termination follows because every rewrite decreases a well-founded
complexity measure on primitive expressions. Local confluence follows
from pairwise verification of overlapping rewrite rules. Since the
rewriting system is terminating and locally confluent, Newman's Lemma
implies global confluence, so every admissible computation possesses a
unique normal form. The complete proof occupies Appendix~\ref{app:proofs}.
\end{proof}

Without canonical normal forms, semantic comparison between different
computations would remain fundamentally ambiguous; normal forms therefore
provide the canonical representatives of semantic equivalence classes.

\section{Admissibility Preservation}

\begin{theorem}[Admissibility preservation]
\label{thm:admissibility}
A computational system possesses structural semantic capacity precisely
when every intermediate transformation preserves admissibility.
\end{theorem}

\begin{proof}
Suppose admissibility fails at some intermediate computation; then at
least one boundary relation becomes invalid, and the resulting
computation no longer represents a legal transformation within $\CatScope$.
Conversely, if every intermediate transformation preserves admissibility,
every composed morphism remains well-defined by the Closure Theorem, so
semantic organization is preserved throughout the computation.
\end{proof}

This theorem formalizes one of the central philosophical claims of
Chapter~1: semantic capacity is characterized by preservation of
admissible structure rather than by phenomenological or behavioral
criteria.

\begin{theorem}[Composition integrity]
Suppose $\ScopeMorph_1:\Scope_1\to\Scope_2$ and
$\ScopeMorph_2:\Scope_2\to\Scope_3$ are admissible. Then
$\ScopeMorph_2\circ\ScopeMorph_1$ is likewise admissible.
\end{theorem}

\begin{proof}
Because both morphisms individually preserve the boundary tensor
$\Boundary$, every intermediate scope remains admissible. Associativity of
composition established in Chapter~8 then implies that their sequential
execution also preserves admissibility.
\end{proof}

This theorem allows arbitrarily long semantic computations to be
constructed from locally admissible components, providing the
compositional principle underlying subsequent chapters.

\begin{theorem}[State reconstruction]
Let $\Scope$ be a bounded scope geometry satisfying the admissibility
conditions of Chapter~6. If every collapse operation records its boundary
history, then the parent scope can be reconstructed uniquely.
\end{theorem}

\begin{proof}
Each collapse records precisely those admissible boundary modifications
required to reconstruct the parent scope. Because admissibility excludes
ambiguous boundary transitions, the reconstruction algorithm is
deterministic. Full details are presented in Appendix~\ref{app:proofs}.
\end{proof}

This theorem plays an important role later in the development of
diagnostic algorithms for transformer models (Chapter~14), where
observable trajectories are used to infer latent scope evolution. The
underlying pattern --- assembling a global object from a family of
locally consistent pieces --- recurs throughout mathematics under the
name of gluing; Appendix~\ref{app:sheaf} develops this reading of
reconstruction explicitly, in terms general enough to state precisely
when local reconstruction data fails to determine a unique global scope.

\section{Representation Invariance}
\label{sec:representation-invariance}

The preceding theorems remain entirely internal: they establish that the
calculus is mathematically coherent, but do not yet address its
principal philosophical motivation. Structural Semantics was introduced
in Chapter~1 as a substrate-independent account of semantic organization.
To justify this claim mathematically, we require a precise criterion
under which two computational systems possessing different internal
architectures may nevertheless instantiate the same semantic computation.
Unlike many discussions of representation in philosophy of mind, the
theorem below makes no claims concerning consciousness, subjective
experience, intentionality, or ontology; it concerns only the
preservation of admissible transformation structure within $\CatScope$.

\begin{definition}[Admissibility-preserving projection]
\label{def:projection}
Let $\Domain$ be a computational state space. A mapping
$h:\Domain\to\Obj(\CatScope)$ is an \emph{admissibility-preserving
projection} if it is surjective; for every admissible transition
$x\rightsquigarrow y$ in $\Domain$, the induced scope transition
$h(x)\rightsquigarrow h(y)$ is admissible in $\CatScope$; and boundary
relations are preserved, $\Boundary(x)=\Boundary(y)\Rightarrow\Boundary(h(x))=\Boundary(h(y))$.
\end{definition}

The definition deliberately does not require injectivity: many
computational states may project onto the same scope geometry, since
different implementations frequently possess redundant internal degrees
of freedom that play no role in semantic organization. The projection
should not be interpreted as a decoder recovering hidden semantic
content, but as an abstraction analogous to quotient mappings in algebra
or projection operators in topology, identifying the structural
information retained by the Scope Calculus while ignoring
implementation-specific variation.

\begin{remark}
The admissibility clause of Definition~\ref{def:projection} is a
compatibility condition, not a free choice: it requires
$x\sim_{\mathrm{adm}}y$ in $\Domain$ to imply
$h(x)\sim_{\mathrm{adm}}h(y)$ in $\CatScope$ for every pair of states
connected by an admissible transition, not merely for some convenient
subset of them. Theorem~\ref{thm:representation} below depends on this
condition holding for the specific $h_1,h_2$ under consideration; the
theorem is conditional on the existence of such projections; it does not
assert that admissibility-preserving projections exist for every pair of
systems one might wish to compare.
\end{remark}

\begin{definition}[Semantic equivalence]
\label{def:semantic_equivalence}
Let $\ScopeMorph_1,\ScopeMorph_2$ be admissible morphisms in $\CatScope$.
They are \emph{semantically equivalent}, $\ScopeMorph_1\sim_{\mathrm{sem}}\ScopeMorph_2$,
whenever their canonical normal forms coincide: $\ScopeMorph_1^*=\ScopeMorph_2^*$.
\end{definition}

Because canonical normal forms are unique (Theorem~\ref{thm:normal}),
this relation is well defined, and it is straightforward to check that
$\sim_{\mathrm{sem}}$ is an equivalence relation on the class of
admissible morphisms; the admissible morphisms therefore partition into
equivalence classes, each represented uniquely by its canonical normal
form.

\begin{theorem}[Representation invariance]
\label{thm:representation}
Let $\mathcal M_1$ and $\mathcal M_2$ be computational systems with state
spaces $\Domain_1,\Domain_2$. Suppose there exist admissibility-preserving
projections $h_1:\Domain_1\to\Obj(\CatScope)$ and
$h_2:\Domain_2\to\Obj(\CatScope)$ such that every admissible computation
induces a morphism of $\CatScope$; corresponding induced morphisms possess
identical canonical normal forms; and corresponding boundary tensors are
preserved under $h_1$ and $h_2$. Then $\mathcal M_1$ and $\mathcal M_2$ are
semantically equivalent relative to the Scope Calculus.
\end{theorem}

\begin{proof}
Let $T_1=(x_0,\ldots,x_n)$ be an admissible computation inside
$\mathcal M_1$, and let $T_2=(y_0,\ldots,y_m)$ be the corresponding
computation inside $\mathcal M_2$. Applying the projections yields
admissible morphisms $\ScopeMorph_1=h_1(T_1)$ and $\ScopeMorph_2=h_2(T_2)$. By
Definition~\ref{def:projection}, both morphisms preserve admissibility.
By Theorem~\ref{thm:normal}, each possesses a unique canonical normal
form, $\ScopeMorph_1^*,\ScopeMorph_2^*$. The second hypothesis states
$\ScopeMorph_1^*=\ScopeMorph_2^*$, hence
$\ScopeMorph_1\sim_{\mathrm{sem}}\ScopeMorph_2$ by
Definition~\ref{def:semantic_equivalence}. Preservation of boundary
tensors guarantees that the equivalence extends beyond the terminal
states to every admissible intermediate transition. Therefore both
systems determine precisely the same semantic computation inside
$\CatScope$, and $\mathcal M_1$, $\mathcal M_2$ are semantically
equivalent relative to the Scope Calculus.
\end{proof}

The theorem deliberately avoids stronger claims: it does not assert that
two computational systems are physically equivalent, nor that they have
equal computational complexity, energy consumption, learning dynamics, or
phenomenological equivalence. It establishes only that the admissible
semantic structure recognized by the Scope Calculus is invariant under
the stated projection hypotheses.

\begin{corollary}
\label{cor:substrate}
Structural semantic capacity depends only on admissible transformation
structure and not on the physical substrate implementing that structure.
\end{corollary}

\begin{proof}
Immediate from Theorem~\ref{thm:representation}.
\end{proof}

The Representation Invariance Theorem serves as the principal
mathematical bridge between the abstract theory of Volumes~I and~II and
the empirical investigations of transformer representations developed in
Volume~III: rather than searching for semantic organization directly
within vectors, tensors, symbolic graphs, or biological neurons,
subsequent chapters investigate whether these systems induce equivalent
admissible transformation structures after projection into the common
categorical framework established by $\CatScope$.

\subsection{Interpretation of Representation Invariance}

The theorem should not be understood as claiming that all computational
architectures are equivalent, but as specifying a precise criterion under
which equivalence may be established --- one based not on physical
realization, implementation language, computational complexity, or
biological plausibility, but solely on the preservation of admissible
transformation structure within $\CatScope$. This mirrors analogous
developments throughout mathematics: the same abstract group may be
represented by matrices, permutations, or symmetries of geometric
objects, and although these realizations differ substantially in
implementation, they are regarded as equivalent because they preserve the
same algebraic structure. Structural Semantics adopts an analogous
perspective, treating the organization of admissible transformations
rather than the implementation or coordinate system as the object of
interest.

\begin{definition}[Substrate independence]
A semantic property is \emph{substrate independent} whenever its validity
depends only on admissible structural transformations within $\CatScope$
and not on the physical realization of those transformations.
\end{definition}

This definition intentionally avoids stronger metaphysical claims: it
does not imply that all substrates are equally efficient, expressive, or
capable of supporting arbitrary computations, only that once admissible
transformations have been projected into the Scope Calculus, semantic
evaluation depends on those transformations rather than on their physical
implementation. Consider, for instance, a compiler translating a
high-level programming language into two different machine
architectures: although the resulting executable instructions differ
substantially, both preserve identical control flow, scope hierarchy, and
variable-binding relationships, and within the Scope Calculus these
executions correspond to the same canonical admissible morphism despite
differing binary representations. Similarly, if a symbolic reasoning
system built on graph rewriting and a transformer neural network induce
identical admissible primitive transformations after projection into
$\CatScope$, Theorem~\ref{thm:representation} establishes their semantic
equivalence relative to the Scope Calculus, without implying identical
learning mechanisms, computational cost, training dynamics, or
phenomenological properties.

Representation Invariance occupies an intermediate position between
several established traditions: denotational equivalence in formal
semantics identifies programs whose observable behaviors coincide despite
differing implementations \cite{Scott1970,Winskel1993}, quotient
constructions in algebra identify objects modulo equivalence relations,
and naturally isomorphic constructions in category theory are routinely
regarded as equivalent despite differing internal presentations
\cite{MacLane1998,Awodey2010}. The present theorem shares this structural
philosophy while preserving a different invariant: instead of numerical
values, logical formulas, or observable program outputs, it preserves
admissible context transformations. This distinction becomes
particularly important in Chapters~13--15, where the empirical question
investigated is not whether latent vectors resemble human concepts
directly, but whether they preserve admissible semantic organization
under projection into Scope Geometry.

\begin{remark}[A parallel in philosophy of physics]
A structurally similar move appears in the philosophy of physics.
\citet{IsmaelVanFraassen2003} argue that a symmetry transformation
identifies surplus theoretical structure precisely when it leaves every
observable prediction unchanged: two models related by such a
transformation, $M\sim M'$, should be treated as one physical state
rather than two, so that the physically meaningful object is the
quotient $\mathcal{M}/\!\sim$ rather than the original space of models.
The admissibility-preserving projections of
Definition~\ref{def:projection} play an analogous role here: two
computational systems related by such a projection are treated, relative
to the Scope Calculus, as one semantic computation rather than two,
regardless of how their underlying representations differ. The two
constructions are not the same --- theirs identifies structure that is
dynamically inert, ours identifies structure that is
admissibility-irrelevant --- but both replace a space of concrete
representations with a coarser space of what those representations are
for. Separately, \citet{Ismael1997} argues that a deterministic process
cannot introduce a symmetry-breaking distinction into its effect that was
not already present, in some form, in its cause: an asymmetric outcome
requires an asymmetric antecedent. The Closure Theorem
(Theorem~\ref{thm:closure}) and Admissibility Preservation Theorem
(Theorem~\ref{thm:admissibility}) make a comparable claim about
admissible transformations specifically: no admissible sequence of
primitive operations can produce a scope geometry violating constraints
that were not already violable given the boundary structure it started
from. Neither parallel is used to derive anything in this monograph; both
are noted because the same structural intuition --- that transformation
can reveal or reorganize distinctions but not manufacture them from
nothing --- recurs across fields that otherwise share little vocabulary.
\end{remark}

\subsection{Representations as Coordinates Rather Than Meaning}

It is worth stating plainly what Theorem~\ref{thm:representation}
implies about the status of representations themselves, since the
implication is easy to state informally but easy to misread formally.
The theorem does not say that vectors, symbols, neural activations, and
graphs are all secretly the same kind of thing. It says that, once an
admissibility-preserving projection into $\CatScope$ has been specified,
the differences among them become differences in \emph{coordinates} on a
common invariant object rather than differences in the object itself ---
in the same sense that a point in the plane has different numerical
descriptions in Cartesian and polar coordinates without being two
different points.

This reframing has a consequence for how the question ``what does this
representation mean?'' should be posed within this framework. Asked of a
coordinate system directly, the question is comparatively hard: a
$4096$-dimensional activation vector, considered as raw numbers, carries
no more intrinsic meaning than a pair of polar coordinates carries
intrinsic direction independent of the convention that assigns it one.
Asked of the admissible transformation the vector participates in, the
question becomes tractable, because admissibility is exactly the
structure the projection is built to preserve. The shift in emphasis
from Chapters~11 through~13, from cataloguing what activation space
\emph{looks like} to tracking what transformations it \emph{admits}, is
this reframing applied to a specific representational format.

The reframing has a limit worth stating alongside it. A choice of
coordinates is never entirely arbitrary in practice: some coordinate
systems make a given computation efficient and others make it
intractable, exactly as Cartesian coordinates simplify some geometric
problems that polar coordinates complicate. Representation Invariance
establishes that semantic content, as this monograph defines it, does not
depend on the choice; it does not establish that the choice is
inconsequential for anything else, including learnability, computational
cost, or the ease with which a projection satisfying
Definition~\ref{def:projection} can be found or verified in the first
place.

\subsection{Limitations of the Present Results}

The Representation Invariance Theorem establishes only a conditional
equivalence: its conclusions depend on the hypotheses stated explicitly
in Theorem~\ref{thm:representation}, and if admissibility-preserving
projections do not exist, or if canonical normal forms fail to coincide,
no semantic equivalence follows. Nor does the theorem provide a
constructive algorithm for discovering such projections; the existence of
admissibility-preserving mappings is a mathematical assumption whose
verification depends on the computational architecture under
investigation, and developing practical procedures for constructing these
projections is one of the principal objectives of Volume~III. The
distinction between existential theorems and constructive algorithms is
fundamental: Chapter~9 proves that semantic equivalence follows if
suitable projections exist, while Chapter~14 investigates algorithmic
procedures for approximating such projections within transformer
activation spaces. The mathematical theorem should therefore not be
interpreted as an implemented diagnostic method, but as the theoretical
target toward which the diagnostic algorithms are directed.

\section{Levels of Structural Equivalence}

The equivalence established by Theorem~\ref{thm:representation} is exact:
two systems are semantically equivalent relative to the Scope Calculus
only if their projected morphisms reduce to the identical canonical
normal form. Part~III will repeatedly ask a looser question --- whether
two activation trajectories are similar enough, in some measurable sense,
to indicate comparable semantic organization --- and it is worth
distinguishing, before that material begins, several notions of
equivalence that exact identity does not exhaust. Conflating them is a
natural source of overclaiming in exactly the kind of empirical work
Chapters~13 through~15 undertake.

\emph{Exact equivalence} is the relation of
Definition~\ref{def:semantic_equivalence}: two morphisms are equivalent
only if $\ScopeMorph_1^*=\ScopeMorph_2^*$, with no tolerance. This is the
only notion of equivalence any theorem in this monograph establishes.

\emph{Approximate equivalence} weakens exact equivalence by a distance
threshold: two morphisms might be considered equivalent if their
canonical normal forms differ by less than some $\epsilon$ under a
suitable metric on $\ScopeAlphabet^{*}$. Nothing in Chapters~6--9
constructs such a metric or proves that any candidate metric interacts
well with composition, so approximate equivalence in this sense remains
an open extension rather than an established consequence of the present
theorems.

\emph{Probabilistic equivalence} replaces a fixed morphism with a
distribution over morphisms, as arises naturally once retrieval or
routing in a real architecture is stochastic rather than deterministic;
two systems might then be called equivalent if their induced
distributions over canonical normal forms coincide or are close in some
divergence. Appendix~G's measure-theoretic constructions provide some of
the vocabulary such a notion would need, but the monograph does not
develop it into a theorem.

\emph{Functional equivalence} is weaker still, and closest to what
Chapters~14 and~15 can actually test: two systems are functionally
equivalent if they produce statistically indistinguishable diagnostic
behavior --- comparable Phase Locking Values, comparable admissibility
estimates --- without any claim that their underlying admissible
transformations coincide, approximately or otherwise. This is the notion
implicitly in play whenever an experiment in Chapter~15 compares model
architectures rather than proving a theorem about them.

The four levels form a strict hierarchy: exact equivalence implies every
weaker notion, but none of the implications reverse. An empirical finding
of functional equivalence between two systems is evidence bearing on,
but not a proof of, probabilistic, approximate, or exact equivalence in
the senses above. Keeping this hierarchy in view is part of what the
Methodological Statement of the front matter asks of every claim in this
monograph: an empirical result belongs to the empirical-hypothesis
category regardless of how strongly it is worded, and dressing a
functional-equivalence finding in the language of Theorem~\ref{thm:representation}
would misstate what has actually been shown.

\section{Chapter Summary}

This chapter established the principal mathematical foundations of
Structural Semantics. Beginning with the Closure Theorem, we demonstrated
that admissible computations remain within the category $\CatScope$. The
Scope Decomposition and Primitive Normal Form Theorems showed that every
admissible computation possesses a canonical structural representation.
The Admissibility Preservation and Composition Integrity Theorems
characterized semantic organization as the preservation of boundary
constraints under sequential transformation, while the State
Reconstruction Theorem established that admissible computational
histories can be recovered from recorded boundary events. These internal
consistency results culminated in the Representation Invariance Theorem,
which demonstrated that semantic equivalence depends on admissible
transformation structure rather than the underlying computational
substrate, providing the mathematical foundation for treating symbolic
systems, neural architectures, graph-based reasoners, and other
computational models within a common structural framework. The next
chapter builds on these results by introducing a translation functor from
Whiteheadian process ontology into Scope Geometry.

\subsection*{Status of Results: Chapter 9}

\paragraph{Established background.} Canonical rewriting systems,
equivalence relations, quotient constructions, denotational semantics,
and elementary category theory
\cite{Scott1970,Winskel1993,MacLane1998,Awodey2010}.

\paragraph{Formal developments.} The Closure, Scope Decomposition,
Primitive Normal Form, Admissibility Preservation, Composition Integrity,
and State Reconstruction Theorems, together with the Representation
Invariance Theorem and its corollary establishing substrate-independent
structural semantic capacity.

\paragraph{Empirical hypotheses.} None. All results established in this
chapter are purely mathematical, conditional on the stated hypotheses.
The construction of practical admissibility-preserving projections is
deferred to the computational framework developed in Volume~III.
%%%%%%%%%%%%%%%%%%%%%%%%%%%%%%%%%%%%%%%%%%%%%%%%%%%%%%%%%%%%%%%%%%%%%%%%%%%%%%
\chapter{Comparative Ontology and the Translation Functor}
\label{ch:translation-functor}
%%%%%%%%%%%%%%%%%%%%%%%%%%%%%%%%%%%%%%%%%%%%%%%%%%%%%%%%%%%%%%%%%%%%%%%%%%%%%%

\vspace{1em}
\begin{flushright}\itshape
The many become one, and are increased by one.\\
\upshape --- Alfred North Whitehead, \emph{Process and Reality}
\end{flushright}
\vspace{1em}

The preceding chapter established the internal mathematical foundations
of Structural Semantics: closure, canonical decomposition, admissibility
preservation, and representation invariance demonstrated that the Scope
Calculus forms a coherent mathematical system independent of any
particular computational substrate. The present chapter serves a
different purpose. Rather than proving additional properties of the
calculus itself, we construct a formal correspondence between the Scope
Calculus and an existing philosophical framework, asking whether the
structural relationships described by Whitehead's process ontology admit
a mathematically precise translation into the language of bounded scope
geometries. The answer developed here is affirmative, although carefully
qualified: we define a functor $F:\CatProcess\to\CatScope$ that preserves
the operational structure of process interactions while explicitly
identifying the point at which the two theories diverge, namely
Whitehead's notion of Subjective Form.

\section{Motivation}

A recurring theme of twentieth-century philosophy is that process
precedes substance: rather than viewing reality as composed of static
objects possessing accidental properties, process philosophies understand
stable objects as relatively persistent organizations of ongoing events.
Although Whitehead developed this view for metaphysical purposes,
remarkably similar ideas appear throughout modern theoretical computer
science --- programming languages describe evolving computational states,
distributed systems describe interacting processes, category theory
emphasizes morphisms before objects, dynamical systems focus on
trajectories rather than isolated points, and large language models
themselves perform successive transformations of latent representations
rather than manipulating static symbolic objects. These parallels
motivate a careful mathematical comparison. The purpose of this chapter is
therefore not to defend Whiteheadian metaphysics, nor to derive the Scope
Calculus from process philosophy, but to investigate the extent to which
the operational structure of one framework may be translated into the
other.

\section{The Category of Processes}

Throughout this chapter we use $\CatProcess$ to denote an abstract category
whose objects represent Whiteheadian actual occasions and whose morphisms
represent admissible process transitions. This construction is not a
complete formalization of Whitehead's philosophy, but a mathematical
abstraction capturing only those structures required for comparison with
Scope Geometry.

\begin{definition}[Category of processes]
The category $\CatProcess$ consists of objects representing actual
occasions, morphisms representing admissible transitions between
occasions, ordinary categorical composition, and identity morphisms
corresponding to persistence of an occasion.
\end{definition}

This definition intentionally avoids introducing unnecessary metaphysical
structure; concepts such as eternal objects, creativity, and the
extensive continuum remain outside the scope of the present construction.

\begin{definition}[Actual occasion, categorical form]
An actual occasion is represented by the ordered triple
$\WPrehensions=\langle\WData,\WValuation,\WSubjForm\rangle$, where $\WData$
denotes the inherited data, $\WValuation$ the polarity of prehension, and
$\WSubjForm$ the Subjective Form associated with the occasion.
\end{definition}

This representation follows the operational interpretation of Chapter~4
and serves as the domain of the translation functor; no assumptions are
yet made concerning the interpretation of $\WSubjForm$, whose treatment
emerges naturally from the construction of the functor.

\section{Operational Correspondence}

Inspection of Chapters~4 through~9 reveals a striking operational
parallel between the two theories: positive prehensions accumulate
structural relations, negative prehensions exclude incompatible relations,
transduction propagates contextual information, and concrescence produces
completed organizations, while the Scope Calculus possesses analogous
operations in binding, refusal, promotion, and canonical evaluation. This
observation motivates the correspondence summarized in
Table~\ref{tab:translation}.

\begin{table}[h]
\centering
\begin{tabular}{ll}
\toprule
Whiteheadian process & Scope Calculus \\
\midrule
Positive prehension & $\Bind$ \\
Negative prehension & $\Refuse$ \\
Transduction & $\Pop$ \\
Concrescence & $\Collapse$ \\
\bottomrule
\end{tabular}
\caption{Operational correspondence between process ontology and Scope
Geometry.}
\label{tab:translation}
\end{table}

The purpose of the remainder of the chapter is to determine whether this
informal correspondence extends to a genuine functor.

\begin{definition}[Translation functor]
The translation functor $F:\CatProcess\to\CatScope$ acts on objects
according to $F(\WPrehensions)=\Scope$, where
$\Scope=\langle\Boundary,\Interior,\Children\rangle$ is the corresponding
bounded scope geometry, and on morphisms maps process transitions to
admissible scope morphisms preserving boundary structure.
\end{definition}

The operational component of the translation is determined by the
assignments $F(P^{+})=\Bind$, $F(P^{-})=\Refuse$, $F(T)=\Pop$, and
$F(C)=\Collapse$, motivated by operational rather than philosophical
similarity: each pair performs analogous structural functions despite
arising from different theoretical traditions, and the interpretation
should be understood as structural rather than metaphysical. The
remaining question concerns Subjective Form, which unlike the preceding
components has no direct analogue inside the Scope Calculus and requires
separate consideration.

\begin{example}[The functor applied to a single occasion]
Let $\WPrehensions=\langle\WData,\WValuation,\WSubjForm\rangle$ describe an
occasion prehending two antecedent data, $d_1$ positively and $d_2$
negatively. Under $F$, the positive prehension of $d_1$ becomes an
instance of $\Bind$, establishing an admissible dependency between $d_1$
and the developing occasion within the corresponding scope; the negative
prehension of $d_2$ becomes an instance of $\Refuse$, excluding $d_2$
from the interior of that scope. If the occasion's data derive from an
antecedent occasion in a different scope, the transduction carrying that
data forward becomes an instance of $\Pop$; the occasion's concrescence
into a completed satisfaction becomes an instance of $\Collapse$, the
scope's resolution into a canonical result. Every component of
$\WPrehensions$ except $\WSubjForm$ is accounted for by one of these four
assignments; nothing in the translation invents a fifth operator or maps
prehension onto a notion --- such as a labeled ``inclusion vector'' or
``exclusion predicate'' --- outside the primitive alphabet
$\ScopeAlphabet$ already fixed in Chapter~7.
\end{example}

\section{Functorial Preservation}

Having defined the action of the translation functor on both objects and
primitive operations, we must verify that it satisfies the defining
axioms of a functor: preservation of identity morphisms and of
composition. Recall from Chapter~8 that a functor between two categories
preserves the categorical structure, so the validity of the present
construction depends on ordinary categorical reasoning
\cite{MacLane1998,Awodey2010} rather than philosophical interpretation.
This is the second functor constructed in this monograph rather than the
first: Appendix~\ref{app:ccc} constructs one between the simply typed
$\lambda$-calculus and $\CatScope$ under a Cartesian closure hypothesis,
and the categorical vocabulary used to verify $F$ here is the same
vocabulary developed rigorously there.

\begin{proposition}[Identity preservation]
\label{prop:identity_preservation}
For every object $P\in\CatProcess$, the translation functor satisfies
$F(\operatorname{id}_{P})=\operatorname{id}_{F(P)}$.
\end{proposition}

\begin{proof}
Let $P=\WPrehensions$ be an arbitrary process object. The identity
morphism $\operatorname{id}_{P}$ produces no modification of inherited
data, process polarity, or admissible transition structure. Under the
action of the translation functor, $P\mapsto\Scope$; since no structural
modification occurs, the corresponding scope geometry also remains
unchanged, so $F(\operatorname{id}_{P})=\operatorname{id}_{\Scope}=\operatorname{id}_{F(P)}$.
\end{proof}

\begin{theorem}[Composition preservation]
\label{thm:functor_composition}
Let $f:P_1\to P_2$ and $g:P_2\to P_3$ be morphisms in $\CatProcess$. Then
$F(g\circ f)=F(g)\circ F(f)$.
\end{theorem}

\begin{proof}
Suppose $f$ induces the primitive sequence $(\alpha_1,\ldots,\alpha_n)$
while $g$ induces $(\beta_1,\ldots,\beta_m)$. Composition in $\CatProcess$
corresponds to concatenation of these transition sequences. The
translation functor maps each primitive process operation to its
corresponding primitive scope operation according to the assignments
above, so $F(g\circ f)$ produces exactly the concatenated admissible
primitive sequence
$F(\beta_1),\ldots,F(\beta_m),F(\alpha_1),\ldots,F(\alpha_n)$, identical to
first applying $F(f)$ followed by $F(g)$. Hence
$F(g\circ f)=F(g)\circ F(f)$.
\end{proof}

\section{The Image of Subjective Form}

The previous sections established a nearly complete structural
correspondence between process ontology and Scope Geometry. One
component, Subjective Form ($\WSubjForm$), has remained untouched, and
this absence is not accidental: it marks the precise mathematical
location at which the two frameworks diverge.

\begin{definition}[Image of the functor]
The image of $F:\CatProcess\to\CatScope$ consists of every object and
morphism generated solely through admissible scope operations.
\end{definition}

\begin{theorem}[Phenomenological boundary]
\label{thm:phenomenological_boundary}
Subjective Form $\WSubjForm$ lies outside the image of the translation
functor: $\WSubjForm\notin\operatorname{Im}(F)$.
\end{theorem}

\begin{proof}
The Scope Calculus is generated entirely by the primitive operational
alphabet $\{\Pop,\Refuse,\Bind,\Collapse\}$. Every object and morphism in
$\CatScope$ is constructed by finite compositions of these operations
together with their induced boundary relations. Subjective Form
$\WSubjForm$ does not participate in any primitive operational rule: it
neither modifies admissibility conditions nor induces additional scope
transformations, so no admissible sequence of primitive operations
generates an object corresponding to it. Therefore $\WSubjForm$ does not
belong to the image of the functor.
\end{proof}

Theorem~\ref{thm:phenomenological_boundary} should not be interpreted as
a refutation of Whitehead's philosophy, but as identifying the precise
boundary between two theories: the translation functor preserves every
operational component of process organization represented within the
Scope Calculus, while Subjective Form remains external because it is
introduced as an additional primitive of the philosophical theory rather
than as an operational component of the structural calculus. This
distinction mirrors the methodological boundary established in Chapter~5
between formal mathematical constructions and phenomenological
interpretations.

\section{Interpretive Consequences}

The existence of the translation functor demonstrates that substantial
portions of Whitehead's process ontology admit an entirely structural
interpretation: positive and negative prehensions become admissible
binding and refusal operations, concrescence becomes canonical collapse,
and propagation of process becomes scope transduction. The resulting
correspondence is remarkably extensive, but not complete, and its
incompleteness is mathematically informative rather than philosophically
problematic. Instead of attempting to eliminate Subjective Form or
reinterpret it as a hidden computational variable, the present framework
identifies the limits of the translation: the image of the functor
consists precisely of those aspects of process ontology that admit
operational realization within the Scope Calculus. Structural Semantics
therefore neither affirms nor denies the existence of phenomenological
experience; it demonstrates only that the structural organization of
semantic transformations can be developed without requiring
phenomenological primitives, while whether additional phenomenological
structure exists remains an independent philosophical question lying
outside the mathematical scope of the present theory.

\section{Chapter Summary}

This chapter constructed the translation functor $F:\CatProcess\to\CatScope$,
verified that it preserves identity morphisms and composition, and
established that Subjective Form lies outside its image, giving a formal
mathematical expression of the phenomenological boundary established in
Chapter~5. The next chapter leaves philosophy behind and turns toward
contemporary machine learning, examining the mathematical structure of
transformer architectures before introducing the novel dynamical and
spectral analyses developed later in Volume~III.

\subsection*{Status of Results: Chapter 10}

\paragraph{Established background.} Elementary category theory, functors,
and Whitehead's process ontology as reconstructed in Chapters~4--5
\cite{MacLane1998,Awodey2010,Whitehead1929}.

\paragraph{Formal developments.} The category $\CatProcess$, the
translation functor $F:\CatProcess\to\CatScope$, proofs of identity and
composition preservation, and the Phenomenological Boundary Theorem
establishing that Subjective Form lies outside the image of $F$.

\paragraph{Empirical hypotheses.} None. This chapter is a purely
mathematical comparative construction.

%%%%%%%%%%%%%%%%%%%%%%%%%%%%%%%%%%%%%%%%%%%%%%%%%%%%%%%%%%%%%%%%%%%%%%%%%%%%%%
\part{Computational Semantics: Transformer Diagnostics}
%%%%%%%%%%%%%%%%%%%%%%%%%%%%%%%%%%%%%%%%%%%%%%%%%%%%%%%%%%%%%%%%%%%%%%%%%%%%%%

%%%%%%%%%%%%%%%%%%%%%%%%%%%%%%%%%%%%%%%%%%%%%%%%%%%%%%%%%%%%%%%%%%%%%%%%%%%%%%
\chapter{The Mathematical Structure of Transformer Architectures}
\label{ch:transformer-structure}
%%%%%%%%%%%%%%%%%%%%%%%%%%%%%%%%%%%%%%%%%%%%%%%%%%%%%%%%%%%%%%%%%%%%%%%%%%%%%%

\vspace{1em}
\begin{flushright}\itshape
Attention is all you need.\\
\upshape --- Ashish Vaswani et al., 2017
\end{flushright}
\vspace{1em}

The previous volume developed the abstract mathematical framework of
Structural Semantics independently of any particular computational
substrate: Scope Geometry, admissibility, categorical structure, and the
Representation Invariance Theorem were formulated without reference to
neural networks, symbolic systems, or biological cognition. Having
established this substrate-independent foundation, we now turn to one
important computational realization, modern transformer architectures.
This chapter differs fundamentally from those preceding it: no new
mathematical framework is introduced, and no novel definitions of
Structural Semantics are proposed. It instead serves as a survey of
established mathematical results concerning transformer architectures,
establishing a common notation and mathematical vocabulary for the
remainder of Volume~III. Maintaining this separation matters
methodologically: the constructions introduced later, most notably the
activation trajectory framework of Chapter~12 and the resonant analytic
signal model of Chapter~13, are intended as original proposals that
should be evaluated against an accurate presentation of the current state
of transformer theory rather than replacing or reinterpreting it.

\section{Historical Development}

The transformer architecture was introduced by \citet{Vaswani2017} as an
alternative to recurrent and convolutional architectures for sequence
modeling. Rather than processing sequences strictly in temporal order,
transformers employ attention mechanisms that permit every token to
interact directly with every other token within the available context
window. This architectural change increased computational parallelism,
since token representations could be processed simultaneously rather than
recursively; made long-range dependencies easier to model, since every
token could attend directly to distant contextual information; and proved
highly scalable, with increasing parameter count, model depth, and
training corpus size producing remarkably predictable improvements across
a wide variety of language tasks \cite{Kaplan2020,Hoffmann2022}. These
developments established transformers as the dominant architecture for
contemporary large language models.

\section{Residual Stream Representation}

Mechanistic interpretability has converged on the residual stream as one
of the principal mathematical objects of transformer analysis
\cite{Elhage2021,Olah2020,Nanda2023}. Throughout this monograph we adopt
the standard notation $\mathbf{X}\in\mathbb{R}^{(L+1)\times T\times
d_{\mathrm{model}}}$, where $L$ denotes the number of transformer layers,
$T$ the sequence length, and $d_{\mathrm{model}}$ the embedding
dimension; each element $\mathbf{x}_{l,t}\in\mathbb{R}^{d_{\mathrm{model}}}$
represents the residual vector corresponding to token $t$ after
completion of layer $l$. The residual stream functions as a persistent
information channel through which every attention block and multilayer
perceptron contributes incremental modifications, providing a natural
mathematical object for analysing evolving semantic organization.

\section{Attention Mechanisms}

Self-attention constitutes the defining computational mechanism of the
transformer architecture. Given an input matrix $X\in\mathbb{R}^{T\times
d}$, linear projections generate query, key, and value matrices
$Q=XW_Q$, $K=XW_K$, $V=XW_V$, and scaled dot-product attention is computed
as
\[
\operatorname{Attention}(Q,K,V)
=
\operatorname{softmax}\!\left(\frac{QK^{\top}}{\sqrt{d_k}}\right)V,
\]
where $d_k$ denotes the key dimension. Multi-head attention repeats this
computation across several independent projection spaces before
concatenating the resulting representations. From the perspective of
Structural Semantics, attention itself should be regarded simply as an
established computational mechanism; no structural interpretation is
imposed at this stage, and such interpretations are deferred until the
activation trajectory framework of Chapter~12.

\section{Feed-Forward Networks}

Each transformer layer also contains a position-wise multilayer
perceptron, typically represented as
$\operatorname{MLP}(x)=W_2\,\phi(W_1x+b_1)+b_2$, where $\phi$ denotes a
nonlinear activation function such as GELU or SwiGLU. Although attention
is often viewed as the defining component of transformers, mechanistic
interpretability studies increasingly recognize MLP layers as major
repositories of learned features \cite{Geva2022}, so the residual stream
reflects the combined effects of attention and feed-forward computation
rather than either mechanism individually.

\section{Activation Manifolds}

Each residual vector $\mathbf{x}_{l,t}$ may be regarded as a point within
the high-dimensional space $\mathbb{R}^{d_{\mathrm{model}}}$, and as
successive transformer layers are evaluated these vectors trace
continuous trajectories through activation space. The collection of all
such vectors defines a high-dimensional activation manifold whose
geometry has become a major topic of mechanistic interpretability
research, investigated through cosine similarity, Euclidean distance,
principal-component structure, singular-vector decompositions, sparse
feature representations, clustering, activation patching, and causal
mediation. These methods have substantially improved understanding of
transformer representations, but remain primarily geometric descriptions
of isolated layers or local neighborhoods. The following chapter
introduces a complementary viewpoint in which activation vectors are
interpreted as continuous trajectories evolving through depth rather than
as independent static points.

\section{Current Directions in Mechanistic Interpretability}

Recent work in mechanistic interpretability has shifted from identifying
individual neurons toward understanding distributed computational
circuits, including induction heads, composition of attention circuits,
feature superposition, sparse autoencoders, causal tracing, and
activation patching \cite{Elhage2021,Anthropic2023,Templeton2024}. These
studies share the common objective of seeking explicit mathematical
explanations for the computations performed by internal representations
rather than treating neural networks as opaque statistical devices.
Structural Semantics aligns naturally with this broader objective, though
the framework developed in the present monograph differs in emphasis:
whereas mechanistic interpretability often studies the functional role of
particular neurons or circuits, subsequent chapters focus on
admissibility-preserving trajectories through activation space, shifting
the object of study from isolated computational units to the global
organization of transformations.

\section{Relationship to Scope Geometry}

At this point no mathematical identification is made between transformer
representations and Scope Geometry; this restraint is deliberate.
Chapter~9 established that semantic equivalence requires explicit
admissibility-preserving projections into $\CatScope$, and the existence
of such projections cannot simply be assumed for neural activation
spaces. Instead the remainder of Volume~III develops these constructions
gradually: Chapter~12 introduces activation trajectories as continuous
dynamical systems, Chapter~13 constructs analytic signal representations
of those trajectories, Chapter~14 develops computational algorithms for
estimating admissibility within latent representations, and Chapter~15
proposes empirical experiments capable of evaluating these constructions.

\section{Chapter Summary}

This chapter reviewed the established mathematical foundations of
transformer architectures, introducing the residual stream tensor,
attention mechanisms, feed-forward networks, activation manifolds, and
current approaches in mechanistic interpretability while adopting the
standard notation used throughout the contemporary literature.
Importantly, no novel theoretical claims were introduced; the purpose has
been exclusively expository. The next chapter begins the original
developments of Volume~III by replacing the conventional view of
activation vectors as isolated points with a dynamical interpretation in
which semantic organization is studied through continuous trajectories
evolving across transformer depth.

\subsection*{Status of Results: Chapter 11}

\paragraph{Established background.} Transformer architecture,
self-attention, residual streams, feed-forward networks, activation
manifolds, scaling laws, and mechanistic interpretability
\cite{Vaswani2017,Kaplan2020,Hoffmann2022,Elhage2021,Olah2020,Nanda2023}.

\paragraph{Formal developments.} None. This chapter establishes notation,
terminology, and mathematical background only.

\paragraph{Empirical hypotheses.} None. Novel empirical hypotheses begin
in Chapter~12 with the dynamical analysis of activation trajectories.

%%%%%%%%%%%%%%%%%%%%%%%%%%%%%%%%%%%%%%%%%%%%%%%%%%%%%%%%%%%%%%%%%%%%%%%%%%%%%%
\chapter{Dynamic Geometry of Activation Trajectories}
\label{ch:trajectory-dynamics}
%%%%%%%%%%%%%%%%%%%%%%%%%%%%%%%%%%%%%%%%%%%%%%%%%%%%%%%%%%%%%%%%%%%%%%%%%%%%%%

\vspace{1em}
\begin{flushright}\itshape
The mathematics of change is not merely the mathematics of states, but of
the paths connecting them.\\
\upshape --- inspired by the dynamical systems tradition
\end{flushright}
\vspace{1em}

The previous chapter introduced the standard mathematical formalism used
to describe transformer architectures, establishing a common vocabulary
without introducing new theory. This chapter begins the original
developments of Volume~III. Rather than studying activation vectors as
isolated points in a high-dimensional space, we examine the trajectories
generated as those vectors evolve throughout the forward computation of a
transformer, a shift from static to dynamical geometry paralleling a
well-established transition throughout mathematics and physics --
differential geometry studies curves rather than isolated coordinates,
dynamical systems investigate flows rather than individual states, and
control theory analyzes trajectories rather than instantaneous
configurations. The constructions developed in this chapter remain
entirely within conventional mathematics, employing only linear algebra,
multivariable calculus, and classical dynamical systems theory; the
analytic signal formulation of Chapter~13 is deliberately deferred until
this ordinary dynamical framework has been established.

\section{Activation Trajectories}

Consider a transformer of $L$ layers. For a fixed input token $t$, the
residual stream produces a sequence of vectors
$\mathbf{x}_{0,t},\mathbf{x}_{1,t},\ldots,\mathbf{x}_{L,t}$, with
$\mathbf{x}_{l,t}\in\mathbb{R}^{d_{\mathrm{model}}}$. Rather than treating
these vectors independently, we interpret them as a trajectory through
activation space.

\begin{definition}[Activation trajectory]
For a fixed token $t$, the activation trajectory is the ordered sequence
$\gamma_t=(\mathbf{x}_{0,t},\mathbf{x}_{1,t},\ldots,\mathbf{x}_{L,t})$;
equivalently, $\gamma_t:\{0,\ldots,L\}\to\mathbb{R}^{d_{\mathrm{model}}}$ is
a discrete curve parameterized by transformer depth.
\end{definition}

This trajectory represents the progressive refinement of the latent
representation associated with a single token as information propagates
through the network. Where Chapter~11 emphasized the geometry of
activation space itself, the present chapter emphasizes the geometry of
motion within that space.

\section{Discrete Dynamical Systems Interpretation}

Transformer evaluation naturally induces a discrete dynamical system:
each layer defines a transition operator
$F_l:\mathbb{R}^{d_{\mathrm{model}}}\to\mathbb{R}^{d_{\mathrm{model}}}$
such that $\mathbf{x}_{l+1}=F_l(\mathbf{x}_l)$, and the complete forward
computation is the composition $F=F_{L-1}\circ\cdots\circ F_1\circ F_0$.
Unlike autonomous dynamical systems, transformer dynamics are generally
layer-dependent, since each layer possesses distinct learned parameters,
so $F_l\neq F_k$ for $l\neq k$ and the resulting system is non-autonomous.
Nevertheless, many classical concepts from dynamical systems theory
remain applicable.

\section{Layerwise Jacobians and Local Linearization}

\begin{definition}[Layerwise Jacobian]
For each layer $F_l$, the layerwise Jacobian is
$J_l=\partial\mathbf{x}_{l+1}/\partial\mathbf{x}_l$.
\end{definition}

The Jacobian describes the local sensitivity of subsequent activations to
small perturbations of the current activation vector; large singular
values correspond to locally expansive directions, small singular values
to locally contractive directions, and the eigenspectrum of $J_l$
therefore provides valuable information concerning trajectory stability.
Classical dynamical systems theory studies nonlinear trajectories through
their local linear approximations: for a sufficiently small perturbation
$\delta\mathbf{x}_l$, a first-order Taylor expansion yields
$\delta\mathbf{x}_{l+1}=J_l\,\delta\mathbf{x}_l+O(\|\delta\mathbf{x}_l\|^2)$,
so that to first order $J_l$ completely determines the local evolution of
perturbations and local trajectory stability can be analyzed
independently of the global nonlinear dynamics.

\section{Curvature and Trajectory Length}

Activation trajectories rarely follow straight paths through latent
space; instead, successive layers continually redirect the
representation, motivating the introduction of curvature.

\begin{definition}[Discrete curvature]
Let $\Delta_l=\mathbf{x}_{l+1}-\mathbf{x}_l$. The discrete curvature at
layer $l$ is $\kappa_l=\|\Delta_{l+1}-\Delta_l\|$.
\end{definition}

Curvature measures the rate at which the direction of the activation
trajectory changes between successive layers: large curvature indicates
rapid representational restructuring, small curvature relatively stable
semantic evolution. Unlike the spectral quantities of the following
chapter, curvature depends only on elementary Euclidean geometry.

\begin{definition}[Trajectory length]
The length of an activation trajectory is
$\mathcal{L}(\gamma)=\sum_{l=0}^{L-1}\|\mathbf{x}_{l+1}-\mathbf{x}_l\|$.
\end{definition}

Trajectory length measures the total representational displacement
experienced during forward computation. Different transformer
architectures frequently exhibit substantially different trajectory
lengths despite achieving similar downstream task performance,
consequently providing a useful architecture-independent geometric
statistic.

\section{Local Stability and Lyapunov Analysis}

One of the central concepts of dynamical systems theory is the Lyapunov
exponent. Although transformers are finite discrete systems rather than
continuous flows, an analogous local quantity may still be defined.

\begin{definition}[Local layerwise Lyapunov exponent]
For each layer $l$, define $\lambda_l=\log\sigma_{\max}(J_l)$, where
$\sigma_{\max}(J_l)$ denotes the largest singular value of the Jacobian.
\end{definition}

Positive values indicate local expansion, negative values contraction,
and values near zero approximately neutral dynamics; these quantities
describe only local behavior and should not be confused with global
Lyapunov exponents defined for autonomous dynamical systems. The singular
value decomposition $J_l=U_l\Sigma_lV_l^{\top}$ provides an orthogonal
decomposition of the local tangent space, in which the right singular
vectors describe perturbation directions entering the layer, the left
singular vectors the corresponding output directions, and the singular
values the local amplification, permitting activation trajectories to be
analyzed in terms of dominant expansion and contraction modes. Later
chapters reinterpret these modes using analytic signal representations.

\section{Relationship to Scope Geometry}

The mathematical constructions introduced thus far remain entirely within
ordinary dynamical systems theory; no assumptions have yet been made
regarding admissibility, boundary-preserving projections, or Scope
Geometry. Nevertheless, the motivation for introducing activation
trajectories should now be apparent. The Representation Invariance
Theorem (Theorem~\ref{thm:representation}) established that semantic
equivalence depends on admissible structural transformations rather than
raw coordinate values, and activation trajectories provide the continuous
objects from which such transformations may ultimately be extracted.
Chapter~13 develops the mathematical machinery required for this step by
constructing analytic signal representations of these trajectories,
allowing phase relationships and synchronization phenomena to be studied
directly.

\section{Chapter Summary}

This chapter reformulated transformer computation as a discrete
dynamical system evolving through activation space. We introduced
activation trajectories, layerwise Jacobians, local linearization,
trajectory curvature, cumulative path length, local Lyapunov exponents,
and principal tangent directions, each derived directly from standard
linear algebra and dynamical systems theory. No spectral or wave-based
interpretation has yet been introduced; the purpose here has been to
establish a rigorous dynamical foundation on which the subsequent
analytic signal framework may be constructed.

\subsection*{Status of Results: Chapter 12}

\paragraph{Established background.} Discrete dynamical systems, Jacobian
matrices, singular value decomposition, Lyapunov analysis, differential
geometry, and local linearization.

\paragraph{Formal developments.} Activation trajectories, discrete
trajectory curvature, trajectory length, local layerwise Lyapunov
exponents, and their application to transformer residual streams.

\paragraph{Empirical hypotheses.} Semantic organization within transformer
models is reflected more faithfully by the geometry of activation
trajectories than by isolated activation vectors. This hypothesis
motivates the spectral analysis developed in Chapter~13.

%%%%%%%%%%%%%%%%%%%%%%%%%%%%%%%%%%%%%%%%%%%%%%%%%%%%%%%%%%%%%%%%%%%%%%%%%%%%%%
\chapter{Resonant Analysis of Activation Trajectories}
\label{ch:resonant-analysis}
%%%%%%%%%%%%%%%%%%%%%%%%%%%%%%%%%%%%%%%%%%%%%%%%%%%%%%%%%%%%%%%%%%%%%%%%%%%%%%

\vspace{1em}
\begin{flushright}\itshape
Every signal contains both amplitude and phase. Understanding often lies
not merely in what changes, but in how changes remain synchronized.\\
\upshape --- author's formulation
\end{flushright}
\vspace{1em}

The previous chapter reformulated transformer computation as a discrete
dynamical system whose principal mathematical object is the activation
trajectory. The present chapter extends that analysis by constructing an
analytic signal representation of activation trajectories using the
discrete Hilbert transform, a construction well established throughout
signal processing and communications theory, where analytic signals
separate amplitude and phase components of a real signal. The novelty
proposed here is not the mathematics of analytic signals themselves, but
their application to transformer activation trajectories. Throughout this
chapter, every reference to oscillation, resonance, standing waves,
synchronization, or phase should be understood as an analytic description
of latent trajectory geometry after transformation into analytic signal
space; no claim is made that physical waves exist inside neural networks.

\section{Motivation}

Many existing interpretability techniques analyze activation magnitude,
feature direction, neuron attribution, or attention weights, generally
treating successive layer activations as independent observations or
static geometric objects. Activation trajectories suggest a different
viewpoint: once activations are interpreted as ordered signals evolving
through layer depth, additional mathematical quantities become
available, including instantaneous phase, phase synchronization, local
frequency, amplitude envelopes, coherence, and spectral stability. These
quantities are standard throughout signal analysis but have been only
minimally explored within transformer representations.

\section{Analytic Signal Representation}

Consider an activation trajectory
$\gamma=(\mathbf{x}_0,\mathbf{x}_1,\ldots,\mathbf{x}_L)$. Each coordinate
defines a discrete real-valued signal indexed by layer depth; applying
the discrete Hilbert transform produces a corresponding analytic
representation.

\begin{definition}[Analytic activation signal]
Let $\mathbf{x}_l\in\mathbb{R}^{d_{\mathrm{model}}}$ denote the residual
activation at layer $l$. The associated analytic signal is
$\mathbf{z}_l=\mathbf{x}_l+i\,\mathcal{H}(\mathbf{x}_l)$, where
$\mathcal{H}$ denotes the discrete Hilbert transform applied
coordinate-wise.
\end{definition}

The analytic signal contains the original activation as its real part and
its Hilbert transform as the imaginary part, together defining a
complex-valued trajectory $\mathbf{z}_0,\ldots,\mathbf{z}_L$ that will
serve as the principal object of analysis throughout the remainder of
this chapter.

\section{Amplitude, Instantaneous Phase, and Synchronization}

The analytic representation permits every activation to be decomposed
into polar coordinates: the instantaneous amplitude $A_l=|\mathbf{z}_l|$
and the instantaneous phase $\phi_l=\arg(\mathbf{z}_l)$. Amplitude
measures the local magnitude of the activation trajectory, while phase
describes its relative orientation within analytic signal space; neither
quantity individually characterizes semantic organization, but meaningful
structure is hypothesized to emerge through the interaction of phase
relationships across layers and tokens.

Signal processing frequently studies synchronization between multiple
analytic signals, a construction we adapt to activation trajectories.
Given two trajectories $\gamma_i$ and $\gamma_j$, their instantaneous
phase difference is $\Delta\phi_l=\phi_l^{(i)}-\phi_l^{(j)}$; perfect
synchronization corresponds to constant phase difference, while rapidly
varying phase indicates desynchronization. The principal synchronization
statistic employed throughout this monograph is the Phase Locking Value
(PLV), widely used in neuroscience and signal processing.

\begin{definition}[Phase Locking Value]
Given $N$ paired observations,
$\operatorname{PLV}=\left|\frac1N\sum_{k=1}^{N}e^{\,i\Delta\phi_k}\right|$.
\end{definition}

The Phase Locking Value satisfies $0\le\operatorname{PLV}\le1$: values
approaching $1$ indicate strong phase coherence, values near $0$ phase
dispersion, and within the present framework PLV is interpreted as a
quantitative measure of synchronization between evolving activation
trajectories.

The derivative of instantaneous phase defines the local frequency; since
transformer layers are discrete we employ finite differences,
$\omega_l=\phi_{l+1}-\phi_l$. Large values correspond to rapid phase
evolution, small values to relatively stable semantic progression, and
$\omega_l$ will later contribute to the detection of abrupt context
transitions.

\section{Semantic Coherence Hypothesis}

Having introduced the principal spectral quantities, we state the central
hypothesis motivating the remainder of Volume~III.

\begin{hypothesis}[Semantic coherence]
Activation trajectories associated with successful reasoning tasks
exhibit higher phase coherence than trajectories associated with
hallucination, context collapse, or inconsistent reasoning.
\end{hypothesis}

This hypothesis is intentionally empirical: unlike the theorems of
Chapters~8--10, it is not proved mathematically, but motivates the
experimental protocols of Chapter~15.

\section{The Marine Operator}

The previous quantities characterize local spectral properties. We now
introduce a derived operator intended to summarize the overall spectral
organization of an activation trajectory. Unlike the amplitude, phase,
and PLV quantities above, which are standard signal-processing
constructions applied to a new object, the Marine Operator is an original
proposal of this monograph: a specific way of combining those quantities
into a single filtering functional, motivated by acoustic and radar
salience filtering.

\begin{definition}[Marine Operator, abstract form]
Let $\gamma$ be an activation trajectory. The Marine Operator
$\mathcal{M}(\gamma)$ is any functional combining instantaneous
amplitude, phase evolution, local synchronization, and trajectory
stability into a single diagnostic operator
$\mathcal{M}:\Gamma\to\mathbb{R}^{k}$, where $\Gamma$ denotes the space of
activation trajectories.
\end{definition}

We now give one concrete instance of this abstract form, built from the
quantities already introduced. Nothing in the remainder of this chapter
depends on this being the only possible instance; it is presented as a
candidate realization, to be evaluated experimentally in Chapter~15
alongside whatever alternatives later work proposes.

\subsection{Doppler Velocity}

The instantaneous frequency $\omega_l=\phi_{l+1}-\phi_l$ introduced above
measures local phase change but says nothing about how that rate of
change is itself evolving across depth. Borrowing the Doppler analogy
from acoustic and radar signal processing, we track this second-order
quantity directly.

\begin{definition}[Doppler velocity]
Let $\phi_l$ be the unwrapped instantaneous phase at layer $l$, with
$\mathrm{wrap}(\phi)=\mathrm{mod}(\phi+\pi,2\pi)-\pi$ applied to keep
successive differences in $(-\pi,\pi]$. The \emph{Doppler velocity} is
the discrete layerwise derivative of the unwrapped phase trajectory,
$v_l=\mathrm{wrap}(\phi_{l+1}-\phi_l)-\mathrm{wrap}(\phi_l-\phi_{l-1})$.
\end{definition}

Large $|v_l|$ indicates that the local frequency $\omega_l$ is itself
changing rapidly --- a trajectory undergoing acceleration through phase
space rather than settling into a stable rate --- which we take as a
candidate indicator of representational restructuring beyond what
$\omega_l$ alone captures.

\subsection{The Salience Gate}

We combine phase coherence and Doppler velocity into a single gate that
attenuates layers exhibiting either low coherence or high Doppler drift.

\begin{definition}[Marine salience gate]
Let $\operatorname{PLV}_l$ denote the Phase Locking Value computed in a window
around layer $l$, and let $v_l$ be the Doppler velocity there. The
\emph{Marine salience gate} is
\[
G_l
=
\sigma_{\mathrm{gate}}
\big(
\alpha\,\operatorname{PLV}_l
-
\beta\,|v_l|
-
\gamma
\big),
\]
where $\sigma_{\mathrm{gate}}(x)=1/(1+e^{-x})$ is the logistic sigmoid,
$\alpha,\beta>0$ weight the two contributions, and $\gamma$ is a static
attenuation threshold.
\end{definition}

The gate satisfies $G_l\in(0,1)$: layers with high coherence and low
Doppler drift receive $G_l$ near $1$ and are passed through largely
unchanged, while layers with low coherence or high drift receive $G_l$
near $0$ and are attenuated. Applying the gate to the original activation
by $\tilde{\mathbf{x}}_l=G_l\cdot\mathbf{x}_l$, and then rescaling to
restore the original norm, yields a filtered trajectory in which
incoherent, rapidly drifting layers contribute less to downstream
diagnostic statistics than stable, phase-locked ones.

\subsection{Energy Conservation}

Filtering by $G_l$ changes the magnitude of $\mathbf{x}_l$ by
construction. The following observation records that this magnitude
change can always be undone by rescaling, so that the gate reshapes
direction and relative emphasis across layers without discarding
activation energy outright.

\begin{proposition}[Energy conservation under rescaling]
Let $E_l=\|\mathbf{x}_l\|_2^2$ and let
$\hat{\mathbf{x}}_l=\kappa_l\,G_l\,\mathbf{x}_l$ with
$\kappa_l=\|\mathbf{x}_l\|_2/(\|G_l\,\mathbf{x}_l\|_2+\epsilon)$ for small
$\epsilon>0$. Then $\|\hat{\mathbf{x}}_l\|_2\to\|\mathbf{x}_l\|_2$ as
$\epsilon\to0$, so the rescaled filtered activation recovers the original
energy $E_l$ exactly in the limit.
\end{proposition}

\begin{proof}
Since $G_l$ is a positive scalar, $\|G_l\mathbf{x}_l\|_2=G_l\|\mathbf{x}_l\|_2$.
Substituting $\kappa_l$ gives
$\|\hat{\mathbf{x}}_l\|_2=\kappa_l G_l\|\mathbf{x}_l\|_2=\dfrac{\|\mathbf{x}_l\|_2}{\|\mathbf{x}_l\|_2+\epsilon/G_l}\|\mathbf{x}_l\|_2$,
which tends to $\|\mathbf{x}_l\|_2$ as $\epsilon\to0$.
\end{proof}

This is a normalization convenience, not a deep result: it guarantees
that applying the Marine filter does not, by itself, cause activation
magnitudes to decay or explode across layers, which matters for using
$G_l$ as a diagnostic weight rather than a replacement for the original
representation. We state it as a proposition rather than a theorem
because its content is a routine renormalization identity, not a
structural claim about admissibility.

The precise choice of $\alpha,\beta,\gamma$, and the decision to treat
$\mathcal{M}(\gamma)$ as $(G_0,\ldots,G_L)$ or as some summary statistic
of it, are left as implementation choices for Chapter~14, which realizes
this operator as part of the diagnostic pipeline.

\section{Interpretive Remarks}

Before proceeding, three observations bear on how these quantities
should be read. None of the spectral quantities introduced here require
the existence of physical oscillations inside transformer networks,
since the Hilbert transform constructs an analytic representation of
arbitrary signals regardless of their physical origin. Resonance should
likewise be interpreted mathematically rather than metaphysically:
within this monograph it denotes persistent synchronization patterns
observed in transformed activation trajectories, a property of the
analytic representation rather than a physical wave propagating through
hardware. Finally, the quantities introduced here remain independent of
Scope Geometry; only in Chapter~14 will admissibility-preserving
projections connect these spectral measurements to the structural
framework developed earlier in the book.

Where the Representation Invariance Theorem established that semantic
equivalence depends on admissible transformations rather than internal
coordinates, the present chapter has instead introduced mathematical
quantities capable of describing how those internal coordinates evolve
in the first place; these spectral quantities are therefore not
themselves semantic invariants, but provide measurable observables from
which admissibility may be inferred. Structural Semantics remains the
theoretical framework, while analytic signal analysis supplies one
possible family of empirical measurement operators.

\section{Chapter Summary}

This chapter introduced the analytic signal representation of activation
trajectories. Using the discrete Hilbert transform, residual stream
vectors were extended into complex-valued trajectories possessing
instantaneous amplitude, phase, local frequency, and synchronization
properties, and these were combined into a concrete instance of the
Marine Operator via the Doppler velocity and the salience gate $G_l$. The
principal empirical proposal is that semantic coherence may be reflected
by stable phase relationships between evolving activation trajectories;
this proposal remains explicitly hypothetical and is subjected to
experimental evaluation in Chapter~15. Before those experiments can be
formulated, the spectral quantities developed here must be transformed
into practical diagnostic algorithms, the objective of the following
chapter.

\subsection*{Status of Results: Chapter 13}

\paragraph{Established background.} Analytic signals, Hilbert transforms,
instantaneous phase, amplitude envelopes, Phase Locking Values, and
classical signal-processing methods.

\paragraph{Formal developments.} The analytic activation signal, discrete
instantaneous frequency, application of Phase Locking Values to
activation trajectories, the abstract definition of the Marine Operator
together with a concrete instance built from the Doppler velocity and
salience gate, and the energy-conservation renormalization proposition.

\paragraph{Empirical hypotheses.} Semantic coherence is associated with
increased phase synchronization between activation trajectories, and the
salience gate's attenuation pattern is associated with regions of
representational instability. These hypotheses are tested in the
experimental framework developed in Chapters~14 and~15.

%%%%%%%%%%%%%%%%%%%%%%%%%%%%%%%%%%%%%%%%%%%%%%%%%%%%%%%%%%%%%%%%%%%%%%%%%%%%%%
\chapter{Transformer Diagnostics and Algorithmic Architecture}
\label{ch:diagnostics}
%%%%%%%%%%%%%%%%%%%%%%%%%%%%%%%%%%%%%%%%%%%%%%%%%%%%%%%%%%%%%%%%%%%%%%%%%%%%%%

\vspace{1em}
\begin{flushright}\itshape
A mathematical theory becomes scientifically useful only when its
quantities can be computed.\\
\upshape --- adapted from Lord Kelvin
\end{flushright}
\vspace{1em}

The preceding chapter introduced the spectral quantities forming the
empirical core of the proposed analytic framework, but did not specify
how these quantities should be extracted from real transformer models.
The objective of the present chapter is therefore computational rather
than theoretical: we construct an algorithmic architecture for estimating
the quantities of Chapter~13 directly from transformer activation
streams. Throughout this chapter we remain deliberately
implementation-neutral; the algorithms described here are intended to be
realizable in PyTorch, JAX, TensorFlow, or equivalent numerical
environments; reference implementations are left for future companion
material rather than included in this monograph.

\section{Diagnostic Pipeline}

The proposed diagnostic framework consists of six sequential stages:
extraction of residual stream activations; construction of activation
trajectories; analytic signal transformation; spectral feature
computation; admissibility estimation; and diagnostic reporting. Each
stage operates independently of any particular language model,
consequently the framework may be applied to open-weight transformer
architectures regardless of parameter count or tokenizer design.

For a transformer model $M$ consisting of $L$ layers and an input
sequence $X=(x_1,\ldots,x_T)$, forward evaluation produces the residual
tensor $\mathbf{R}\in\mathbb{R}^{(L+1)\times T\times d_{\mathrm{model}}}$,
whose elements $\mathbf{r}_{l,t}$ record the residual activation
associated with token $t$ after completion of layer $l$; this tensor
forms the raw observational data for all subsequent diagnostic
computations. For each token $t$ we define the activation trajectory
$\gamma_t=(\mathbf{r}_{0,t},\mathbf{r}_{1,t},\ldots,\mathbf{r}_{L,t})$, and
the collection $\Gamma=\{\gamma_1,\ldots,\gamma_T\}$ contains one
trajectory for every token within the context window; subsequent
computations operate exclusively on these trajectories rather than
directly on the residual tensor.

Each activation trajectory undergoes analytic transformation using the
construction of Chapter~13. For every layer $l$, the amplitude envelope
$A_l$, instantaneous phase $\phi_l$, instantaneous frequency $\omega_l$,
local phase difference $\Delta\phi_l$, Phase Locking Value, and trajectory
curvature $\kappa_l$ are computed, collectively defining a spectral
descriptor vector $\mathbf{s}_t$ for every token trajectory.

\begin{definition}[Diagnostic state vector]
For every activation trajectory $\gamma_t$, define
$\mathbf{d}_t=(A,\phi,\omega,\operatorname{PLV},\kappa,\lambda)$, where
$\lambda$ denotes the local Lyapunov statistic of Chapter~12.
\end{definition}

The diagnostic state vector summarizes the local dynamical behavior of an
activation trajectory; subsequent statistical procedures operate on
$\mathbf{d}_t$ rather than on individual observables.

\section{Admissibility Estimation and Context Collapse}

The central objective of the diagnostic framework is not merely to
compute spectral quantities but to estimate semantic admissibility.

\begin{definition}[Admissibility estimator]
An admissibility estimator is a mapping $\mathcal{A}:\mathbf{d}\to[0,1]$
assigning each diagnostic state vector an estimated probability of
admissible semantic evolution.
\end{definition}

The precise form of $\mathcal{A}$ is intentionally left unspecified.
Different implementations may employ probabilistic classifiers, Bayesian
estimators, neural predictors, or deterministic threshold rules, allowing
the mathematical framework developed in earlier chapters to remain
independent of particular machine-learning techniques. The salience gate
$G_l$ of Chapter~13 is one such candidate: setting
$\mathcal{A}(\mathbf{d}_t)$ to the layerwise average of $G_l$ over a
token's trajectory yields a fully specified, if simple, admissibility
estimator built entirely from quantities already defined; nothing in the
remainder of this chapter depends on that particular choice, and the
experiments of Chapter~15 evaluate it against the alternatives listed
there. Whichever estimator is chosen, treating $\mathcal{A}(\mathbf{d}_t)$
as a probability of admissible evolution presupposes a formal measure of
admissibility over semantic space in the first place;
Appendix~\ref{app:measure} constructs exactly such a measure, together
with the entropy and mutual-information quantities that a more
statistically grounded estimator than the salience gate might use in
its place.

\begin{definition}[Context collapse]
A context collapse event occurs whenever $\mathcal{A}(\mathbf{d}_t)<\tau$
for some admissibility threshold $\tau\in(0,1)$.
\end{definition}

The threshold $\tau$ is a hyperparameter whose optimal value depends on
the application domain. Rather than identifying isolated incorrect
tokens, this definition identifies regions of semantic instability within
the evolving activation trajectory.

\section{Computational Complexity}

The proposed diagnostic framework introduces relatively modest additional
computational cost. Residual stream extraction requires no additional
forward passes, trajectory construction is linear in both sequence length
and transformer depth, and analytic signal computation using the Fast
Fourier Transform requires $O(L\log L)$ operations per activation
coordinate. The overall complexity of the diagnostic pipeline is
therefore $O(LTd_{\mathrm{model}})$ up to logarithmic factors associated
with Hilbert transform implementation, scaling linearly with the size of
the activation tensor.

\section{Case Study: Multi-Head Attention Residuals}
\label{sec:mhar-case-study}

The preceding section compared the diagnostic architecture of this
chapter with existing interpretability techniques at the level of
purpose. It is also useful to examine a single recent architectural
proposal in enough detail to see the Scope Calculus vocabulary of
Chapters~6--9 applied to something concrete. This section reads
\citet{Luo2026} through that vocabulary. The reading is offered as an
interpretation, not as a formal result: nothing here constructs an
admissibility-preserving projection in the sense of
Definition~\ref{def:projection} and verifies its hypotheses, so no claim
is made that the architecture instantiates $\CatScope$ in the way
Theorem~\ref{thm:representation} would require. The value of the
exercise is illustrative --- it shows what the Scope Calculus vocabulary
buys a reader confronted with an ordinary architecture paper, and where
that vocabulary runs out.

\subsection{The Architecture}

An ordinary transformer residual stream accumulates layer outputs by
addition, $h_\ell=h_{\ell-1}+f_\ell(h_{\ell-1})$, so that an earlier
contribution $s_i$ remains causally present in $h_\ell$ but is no longer
separately addressable: it has been folded into a running sum rather than
retained as a distinguishable object. Attention Residuals
\citep{Luo2026} relax this by keeping the sequence of prior layer outputs
$S_\ell=(s_0,\ldots,s_\ell)$ available and retrieving from it via a
learned softmax query, so that each layer reads a weighted combination
$\sum_i\alpha_i s_i$ of its full depth history rather than only its
immediate predecessor. Multi-Head Attention Residuals (MHAR) then
observe that this retrieval query is shared across the entire
representational width: every feature subspace of $h_\ell$ is forced to
read the same weighted combination $\alpha$ of the history, even when
different subspaces would prefer to weight the same history differently.
MHAR reshapes the single query into $H$ per-subspace queries, giving each
subspace its own softmax over the depth history, at essentially no
additional parameter or compute cost; the ordinary Attention Residuals
architecture is recovered exactly as the special case $H=1$. Trained
from scratch across several model scales, MHAR improves validation loss
over both the plain residual stream and single-query Attention Residuals,
with the improvement reported to grow with model width.

\subsection{Reading the Architecture as Scope Retrieval}

The depth history $S_\ell=(s_0,\ldots,s_\ell)$ available at retrieval
time reads naturally as an admissible source family in the sense of
Chapter~6: at each layer, the retrieval operation must select from among
several available prior states rather than being handed a single
determined predecessor, much as a child scope's Interior consists of
several admissible states rather than one. The retrieval itself, a
weighted combination followed by use in the next sublayer, matches the
two-step pattern of $\Bind$ followed by $\Collapse$ already used in
Appendix~\ref{app:ccc} to interpret function application: the history is
first bound into a candidate combination, and that combination is then
collapsed into the single vector the next layer consumes.

Read this way, ordinary Attention Residuals apply one $\Collapse$
operation across the entire width of $h_\ell$: a single admissible
combination is selected, and every subspace of the output is constructed
from that same combination. MHAR instead applies $H$ separate collapse
operations, one per subspace, each free to select a different admissible
combination of the same history. Chapter~7 already distinguishes a
single global admissibility check from a factorized one; MHAR is a
concrete instance of the latter, applied not to the primitive alphabet
itself but to the routing mechanism built on top of it.

\subsection{What the Reading Clarifies, and What It Does Not}

The reading clarifies why the single-query architecture pays an
increasing cost as width grows: forcing every subspace through one
$\Collapse$ is, in the vocabulary of this monograph, an admissibility
restriction that becomes more binding as the number of genuinely
different subspaces increases, since more of them are made to accept a
combination that was not built for them specifically. It also clarifies
why $H=1$ recovering the original architecture exactly is the right
design choice from this monograph's perspective: it is the discipline,
familiar from the zero-initialized gates discussed in Chapter~1's
distinction between admissible and inadmissible transitions, of
introducing a structural change as an identity transformation before
allowing it to diverge.

The reading does not establish that MHAR's retrieval operators satisfy
the boundary-preservation conditions of Definition~\ref{def:projection},
nor that the Representation Invariance Theorem applies to a comparison
between the single-query and multi-query architectures. Doing so would
require specifying what boundary constraints the depth history is
required to preserve and verifying that both retrieval mechanisms
respect them, which is an empirical and architectural question this
monograph does not address. The case study should therefore be read as
showing that an independently developed, empirically motivated
architecture exhibits the same qualitative shape --- a single shared
collapse operation restricting what a heterogeneous system can express,
relaxed by factorizing that collapse per subspace --- that Chapters~6--9
develop abstractly, not as evidence for any theorem in this monograph.

\section{Relationship to Existing Interpretability Methods}

The proposed diagnostic architecture complements rather than replaces
existing interpretability techniques: attention visualization identifies
communication pathways between tokens, activation patching identifies
causal computational circuits, sparse autoencoders identify approximately
monosemantic features, and logit attribution estimates token-level
prediction contributions, whereas the present framework analyzes the
global dynamical evolution of representations. These approaches are
therefore largely orthogonal and may be combined within a unified
interpretability workflow.

\section{Limitations}

Several limitations should be emphasized. The diagnostic state vector
provides descriptive rather than causal information --- high phase
synchronization does not by itself imply successful reasoning.
Admissibility estimation depends on the quality of the chosen estimator
$\mathcal{A}$, and poor statistical estimation may produce unreliable
diagnostic results even if the underlying mathematical framework is
correct. Different transformer architectures may also exhibit
substantially different spectral statistics despite similar downstream
behavior, so normalization procedures will likely be required before
cross-model comparisons become meaningful. These limitations motivate the
controlled experimental framework developed in the following chapter.

\section{Chapter Summary}

This chapter translated the mathematical quantities of Chapter~13 into a
practical computational pipeline. Beginning with residual stream
extraction, activation trajectories were transformed into analytic
signals, summarized through diagnostic state vectors, and mapped to
estimated admissibility scores capable of identifying potential regions
of semantic instability. None of these procedures establishes the
empirical validity of the proposed framework; they specify only how the
relevant measurements may be computed. The scientific evaluation of these
measurements requires carefully controlled experiments, statistical
hypothesis testing, and explicit falsification criteria, the subject of
the next chapter.

\subsection*{Status of Results: Chapter 14}

\paragraph{Established background.} Residual stream extraction,
computational complexity analysis, signal processing algorithms,
statistical feature extraction, and the Attention Residuals and
Multi-Head Attention Residuals architectures of \citet{Luo2026}.

\paragraph{Formal developments.} The diagnostic pipeline, diagnostic
state vector, admissibility estimator, and formal definition of context
collapse as an admissibility-threshold event.

\paragraph{Interpretive discussion.} \S\ref{sec:mhar-case-study} reads
Multi-Head Attention Residuals through the Scope Calculus vocabulary of
Chapters~6--9, as an illustration rather than a theorem: no
admissibility-preserving projection is constructed for it, and no claim
of the monograph depends on the reading.

\paragraph{Empirical hypotheses.} Diagnostic quantities derived from
activation trajectories can identify regions of semantic instability
prior to observable output failure. The validity of this hypothesis is
evaluated experimentally in Chapter~15.

%%%%%%%%%%%%%%%%%%%%%%%%%%%%%%%%%%%%%%%%%%%%%%%%%%%%%%%%%%%%%%%%%%%%%%%%%%%%%%
\chapter{Experimental Predictions, Controls, and Falsification Protocols}
\label{ch:experiments}
%%%%%%%%%%%%%%%%%%%%%%%%%%%%%%%%%%%%%%%%%%%%%%%%%%%%%%%%%%%%%%%%%%%%%%%%%%%%%%

\vspace{1em}
\begin{flushright}\itshape
The value of a scientific theory lies not merely in what it explains, but
in what observations could demonstrate it to be incorrect.\\
\upshape --- Karl Popper (adapted)
\end{flushright}
\vspace{1em}

The preceding chapters introduced the mathematical and computational
framework of Structural Semantics for transformer analysis. The present
chapter addresses the central scientific question: can these proposed
quantities be shown to possess explanatory power beyond existing
interpretability methods? Unlike previous chapters, which primarily
developed mathematical structures, this chapter develops an empirical
research program in which every hypothesis is explicitly falsifiable.
Failure of the stated predictions should be interpreted as evidence
against the proposed framework rather than evidence requiring ad hoc
modification.

\section{Scientific Methodology}

The empirical program developed in this monograph follows four guiding
principles: every proposed diagnostic quantity must be measurable using
existing transformer architectures; every hypothesis must admit explicit
statistical rejection; comparisons must be performed against established
baseline methods; and negative results are regarded as scientifically
valuable because they delimit the scope of applicability of the proposed
framework. These principles reflect the methodological distinction
introduced in Chapter~1 between formal mathematics and empirical science.

\section{General Experimental Design}

Let $\mathcal{M}=\{M_1,\ldots,M_K\}$ denote a collection of transformer
models. Experiments are conducted over a benchmark collection
$\mathcal{D}=\{D_1,\ldots,D_N\}$ containing reasoning tasks, mathematical
proofs, code generation, multi-document question answering, and
long-context inference problems. For every prompt $x\in\mathcal{D}$, each
model produces output tokens, residual stream activations, activation
trajectories, diagnostic state vectors, and admissibility estimates,
forming the dataset used throughout the remainder of the chapter.

\section{Hypothesis I: Phase Coherence and Semantic Stability}

\begin{hypothesis}[Phase coherence hypothesis]
Successful reasoning trajectories exhibit larger average Phase Locking
Values than trajectories containing semantic inconsistency or
hallucination: $\operatorname{PLV}_{\mathrm{correct}}>\operatorname{PLV}_{\mathrm{incorrect}}$.
\end{hypothesis}

A benchmark of at least $N\ge1000$ paired reasoning traces should be
constructed, each independently classified by human annotators as
correct, partially correct, hallucinated, or internally inconsistent, and
average Phase Locking Values then computed for each category. The null
hypothesis is $H_0:\mu_{\mathrm{correct}}=\mu_{\mathrm{incorrect}}$
against the alternative
$H_1:\mu_{\mathrm{correct}}>\mu_{\mathrm{incorrect}}$; depending on the
empirical distribution, Student's $t$-test, Welch's test, or the
Mann--Whitney $U$ test may be employed, with statistical significance
assessed throughout this monograph at $\alpha=0.01$. The hypothesis is
rejected whenever $p>0.01$, or when the observed effect size is
practically negligible despite statistical significance.

\section{Hypothesis II: Frequency Drift and Context Collapse}

\begin{hypothesis}
Abrupt semantic inconsistency is accompanied by statistically detectable
changes in instantaneous frequency.
\end{hypothesis}

Let $\omega_l$ denote the instantaneous frequency of Chapter~13; large
discontinuities $|\omega_{l+1}-\omega_l|$ are hypothesized to precede
observable failures of reasoning. Long-context reasoning tasks are
modified through controlled insertion of contradictory premises, and the
resulting activation trajectories are analyzed before, during, and after
the injected contradiction and compared with ordinary successful
reasoning traces. If no statistically significant frequency differences
are observed between successful and contradictory reasoning trajectories,
the hypothesis is rejected.

\section{Hypothesis III: Admissibility and Error Prediction}

\begin{hypothesis}
Declining admissibility estimates precede observable output errors.
\end{hypothesis}

Rather than identifying hallucinations after incorrect tokens have
already appeared, admissibility estimation is hypothesized to detect
semantic instability during internal computation. For every generated
sequence, the first observable output error is recorded, diagnostic
quantities computed before that token are analyzed, and receiver
operating characteristic curves and precision--recall statistics are
calculated to determine predictive performance, summarized using the
area under the ROC curve, precision, recall, the F$_1$ score, and the
Matthews Correlation Coefficient.

\begin{example}[What this hypothesis does and does not claim]
Suppose the experiment above returns a strong result: admissibility
estimates decline detectably several layers before an incorrect token is
sampled, with an area under the ROC curve well above chance. This would
establish, for the specific models and tasks tested, that the diagnostic
quantities of Chapter~14 carry predictive information about upcoming
output errors --- an empirical finding, reported with the statistical
qualifications Chapter~15 requires throughout. It would not, by itself,
establish a working intervention: nothing in this chapter specifies what
a system should do upon detecting declining admissibility, whether
halting, revising, or continuing is preferable in a given context, or
what computational overhead an intervention would add beyond the
detection step itself. A system that halts and backtracks on a
declining-admissibility signal is a natural extension to propose, but it
is a distinct research question --- one requiring its own experimental
validation and its own comparison against existing decoding strategies
--- not a consequence already established by Hypothesis III.
\end{example}

\section{Control Experiments}

Perhaps the greatest danger facing any new diagnostic quantity is that it
merely reproduces information already contained in simpler statistics, so
every experiment above must be accompanied by control analyses. Phase
coherence must be distinguished from token confidence: softmax entropy is
computed alongside PLV, and regression analysis determines whether Phase
Locking Values retain predictive power after controlling for entropy.
Large activation magnitudes might artificially influence spectral
statistics, so activation vectors are normalized before analytic
transformation and equivalent experiments are repeated using raw,
normalized, and randomized activations. To determine whether
synchronization reflects genuine structure rather than statistical
coincidence, phase values are randomly permuted while preserving
amplitude and the diagnostic framework re-evaluated, with significant
performance degradation supporting the interpretation that phase
relationships contain meaningful information. Finally, experiments should
be repeated across transformer families differing in parameter count,
tokenizer, training corpus, attention implementation, and instruction
tuning, since agreement across architectures would strengthen the claim
that the observed phenomena reflect general computational organization
rather than architecture-specific artefacts.

\section{Baseline Comparisons}

The proposed framework should not be evaluated in isolation, but compared
against established diagnostic approaches including attention entropy,
logit confidence, activation norm, perplexity, sparse autoencoder feature
activation, activation patching methods, and uncertainty estimation. A
proposed diagnostic quantity possesses scientific value only if it
provides predictive information beyond these existing baselines.

\section{Failure Modes}

Several outcomes would argue against the proposed framework: Phase
Locking Values showing no predictive relationship with semantic
correctness; instantaneous frequency exhibiting no systematic
relationship with context collapse; admissibility estimates failing to
predict errors before observable output failure; existing baseline
methods consistently outperforming every proposed spectral quantity; or
performance failing to generalize across transformer architectures.
Observation of any of these outcomes should be interpreted as evidence
requiring revision or rejection of the corresponding hypotheses.

\section{Discussion}

The purpose of the present chapter is not to demonstrate that the
proposed framework succeeds, but to specify the conditions under which
success or failure may be determined objectively. This distinction is
essential: the mathematical framework developed in Volumes~I and~II
remains valid independently of these experiments, whereas the empirical
usefulness of the diagnostic operators of Volume~III depends entirely on
experimental evidence. The separation between formal mathematics and
empirical evaluation therefore remains explicit throughout the monograph.

\section{Chapter Summary}

This chapter established a complete experimental program for evaluating
the diagnostic framework proposed throughout Volume~III. Three principal
hypotheses concerning phase coherence, frequency drift, and admissibility
prediction were formulated together with statistical procedures, control
experiments, baseline comparisons, and explicit falsification criteria.
The completion of this empirical program concludes the technical
development of Structural Semantics; the remaining volume returns to a
broader perspective, comparing the framework with neighboring fields,
identifying its limits, and discussing future directions.

\subsection*{Status of Results: Chapter 15}

\paragraph{Established background.} Experimental design, statistical
hypothesis testing, effect sizes, receiver operating characteristic
analysis, and benchmark evaluation.

\paragraph{Formal developments.} Explicit experimental protocols,
statistical controls, baseline comparisons, and falsification criteria
for evaluating the proposed diagnostic framework.

\paragraph{Empirical hypotheses.} Phase synchronization, instantaneous
frequency, and admissibility estimation provide measurable indicators of
semantic coherence whose predictive value exceeds existing baseline
diagnostics. These hypotheses remain subject to experimental confirmation
or rejection.

%%%%%%%%%%%%%%%%%%%%%%%%%%%%%%%%%%%%%%%%%%%%%%%%%%%%%%%%%%%%%%%%%%%%%%%%%%%%%%
\part{Synthesis}
%%%%%%%%%%%%%%%%%%%%%%%%%%%%%%%%%%%%%%%%%%%%%%%%%%%%%%%%%%%%%%%%%%%%%%%%%%%%%%

%%%%%%%%%%%%%%%%%%%%%%%%%%%%%%%%%%%%%%%%%%%%%%%%%%%%%%%%%%%%%%%%%%%%%%%%%%%%%%
\chapter{Cross-Disciplinary Integration and Comparative Landscape}
\label{ch:comparative-landscape}
%%%%%%%%%%%%%%%%%%%%%%%%%%%%%%%%%%%%%%%%%%%%%%%%%%%%%%%%%%%%%%%%%%%%%%%%%%%%%%

\vspace{1em}
\begin{flushright}\itshape
Science progresses not merely by constructing new theories, but by
understanding precisely how those theories relate to existing ones.\\
\upshape --- author's formulation
\end{flushright}
\vspace{1em}

The preceding three volumes developed Structural Semantics from first
principles. The purpose of the present chapter is fundamentally
different: rather than introducing additional mathematical structures, we
situate Structural Semantics within the broader scientific landscape.
Throughout this chapter our objective is comparison rather than
competition; Structural Semantics is not presented as replacing existing
frameworks, but as addressing a different class of questions through a
different mathematical language. To facilitate comparison, we evaluate
each framework according to four questions: what constitutes a semantic
object, how is semantic change represented, how is correctness evaluated,
and what mathematical structures are fundamental.

\section{Predictive Processing and Active Inference}

Predictive Processing and Active Inference describe cognition as the
continuous minimization of prediction error under hierarchical generative
models \cite{Friston2010,Clark2013}. Structural Semantics shares with
this tradition an emphasis on dynamics over static representation and on
hierarchical organization, but its mathematical objectives differ:
Predictive Processing centers on probabilistic inference and free-energy
optimization, while Structural Semantics studies admissibility-preserving
transformations independent of probabilistic interpretation, with
prediction error playing no primitive role in the Scope Calculus.
Consequently the two approaches should be regarded as complementary, one
emphasizing optimization and the other structural constraint.

\section{Mechanistic Interpretability}

Mechanistic interpretability attempts to explain neural computation by
identifying internal circuits responsible for observable model behavior
\cite{Elhage2021,Olah2020,Nanda2023}. Structural Semantics aligns closely
with this program in seeking internal mathematical explanations rather
than behavioral descriptions, but the two frameworks operate at different
levels of description: circuit analysis identifies computational
components such as neurons, attention heads, and sparse features, while
Scope Geometry analyzes the global organization of admissible
trajectories connecting those components.

\section{Information Geometry}

Information geometry studies statistical models using differential
geometric methods, introducing Fisher information, geodesics, divergence
functions, and statistical manifolds \cite{Amari2016}. Both frameworks
employ manifold language and investigate geometric organization, but
information geometry derives its structure from probability
distributions while Scope Geometry derives its structure from
admissibility constraints, so the corresponding metrics encode
fundamentally different objects --- statistical distinguishability in one
case, preservation of contextual organization in the other.

\section{Category Theory}

Category theory provides one of the principal mathematical languages of
this monograph, though Structural Semantics employs only a deliberately
restricted subset of categorical machinery: objects represent bounded
scope geometries and morphisms represent admissible context-preserving
transformations, with no assumptions regarding limits, exponentials,
adjunctions, enriched categories, or higher categories unless explicitly
introduced. This structural minimalism distinguishes the present
framework from category-theoretic approaches seeking maximal abstraction;
category theory functions here primarily as an organizational language
for admissible transformations.

\section{Formal Language Theory}

Formal language theory studies symbolic computation through grammars,
automata, rewrite systems, and logical calculi. Structural Semantics
inherits several ideas from this tradition: the primitive alphabet
$\ScopeAlphabet$ functions analogously to a rewriting alphabet, canonical
normal forms resemble rewriting theory, and admissible transformations
resemble typed derivations. The principal distinction concerns
semantics: formal language theory primarily studies syntactic
derivability, while Structural Semantics studies preservation of
contextual admissibility, so that syntax becomes one particular instance
of a broader structural framework.

\section{Classical Structural Semantics}

The terminology adopted in this monograph inevitably overlaps with the
classical linguistic tradition known as Structural Semantics
\cite{Lyons1977,Coseriu1977}, discussed at length in
Chapter~\ref{ch:historical-foundations}. Both approaches emphasize
relationships rather than isolated entities and reject purely referential
theories of meaning, but the mathematical scope differs substantially:
classical structural semantics studies relationships between linguistic
units, while the present work studies admissible transformations within
arbitrary computational systems, so the relationship should be
understood as one of historical continuity rather than mathematical
identity.

\section{Process Philosophy}

Chapter~\ref{ch:translation-functor} established a formal translation
functor between Whiteheadian process ontology and Scope Geometry,
demonstrating that many operational aspects of process philosophy admit
structural reformulation. Structural Semantics neither adopts nor rejects
Whitehead's metaphysical commitments, identifying instead those
components capable of operational formalization; Subjective Form
remained outside the image of the translation functor, marking the
precise boundary between phenomenological philosophy and structural
mathematics.

\section{Large Language Models}

Modern large language models constitute the primary computational domain
investigated throughout Volume~III, but Structural Semantics should not
be regarded as a theory of transformers alone. The Representation
Invariance Theorem established that semantic organization depends on
admissible transformations rather than specific architectures, so nothing
in principle restricts the framework to transformers; recurrent neural
networks, graph neural networks, state-space models, symbolic systems,
and hybrid neuro-symbolic architectures may equally be analyzed whenever
appropriate admissibility-preserving projections can be constructed.

\section{Synthesis}

Comparison across these diverse frameworks reveals a common pattern:
every approach identifies an important aspect of semantic
organization --- Predictive Processing emphasizes optimization,
mechanistic interpretability internal computation, information geometry
statistical structure, formal language theory symbolic derivation,
classical structural semantics lexical relations, and process philosophy
becoming. Structural Semantics contributes a different perspective, whose
central object is neither probability, syntax, prediction, consciousness,
nor optimization, but the preservation of admissible structural
transformations under changing representations, occupying a distinct
position within the broader landscape of semantic theory.

\section{Chapter Summary}

This chapter situated Structural Semantics within the wider scientific
and philosophical literature. Rather than competing with existing
approaches, the framework was shown to address a complementary class of
questions centered on admissible structural transformation. The next
chapter explicitly enumerates the boundary conditions, scope limits, and
non-claims of Structural Semantics, clarifying exactly what the framework
does, and does not, attempt to explain.

\subsection*{Status of Results: Chapter 16}

\paragraph{Established background.} Predictive Processing, Active
Inference, mechanistic interpretability, information geometry, formal
language theory, classical structural semantics, category theory, and
process philosophy.

\paragraph{Formal developments.} Comparative positioning of Structural
Semantics relative to existing frameworks and clarification of its
distinctive mathematical focus.

\paragraph{Empirical hypotheses.} None. This chapter is comparative and
interpretive rather than mathematical or experimental.

%%%%%%%%%%%%%%%%%%%%%%%%%%%%%%%%%%%%%%%%%%%%%%%%%%%%%%%%%%%%%%%%%%%%%%%%%%%%%%
\chapter{Scope Limits, Boundary Conditions, and Non-Claims}
\label{ch:scope-limits}
%%%%%%%%%%%%%%%%%%%%%%%%%%%%%%%%%%%%%%%%%%%%%%%%%%%%%%%%%%%%%%%%%%%%%%%%%%%%%%

\vspace{1em}
\begin{flushright}\itshape
The strength of a mathematical theory is measured not only by what it
explains, but equally by what it deliberately leaves unexplained.\\
\upshape --- author's formulation
\end{flushright}
\vspace{1em}

Every mathematical theory possesses boundaries, and one of the
distinguishing characteristics of mature scientific frameworks is the
explicit recognition of those questions lying beyond their intended
domain of applicability. The purpose of this chapter is therefore not to
weaken the preceding results but to strengthen their interpretation by
specifying exactly what Structural Semantics does and does not claim.
Throughout, the distinction introduced in Chapter~1 between formal
mathematics, empirical hypotheses, and philosophical interpretation
remains strictly enforced.

\section{The Scope of Structural Semantics}

Structural Semantics is fundamentally a theory of admissible
transformations. Its primary mathematical object is neither
consciousness, intelligence, meaning in the ordinary linguistic sense,
nor biological cognition, but the preservation of structural organization
under transformations constrained by explicit boundary relations.
Accordingly, every theorem proved throughout Volumes~I--III concerns
admissibility, context preservation, structural stability, categorical
composition, representation invariance, activation trajectory geometry,
or diagnostic observables; no theorem proved in this monograph establishes
any claim beyond these domains.

\section{Consciousness, Truth, Agency, and Ethics}

Structural Semantics makes no claim that systems possessing high semantic
capacity are conscious, nor that conscious systems necessarily possess
maximal structural semantic organization; phenomenal consciousness
concerns subjective experience, while Structural Semantics concerns
admissible transformations, and whether these domains ultimately
intersect remains an open philosophical question outside the mathematical
scope of the present framework.

\begin{proposition}[Non-claim: consciousness]
None of the theorems established in Chapters~6--15 implies the existence,
absence, measurement, simulation, or explanation of phenomenal
consciousness.
\end{proposition}

\begin{proof}
Immediate. Neither phenomenal predicates nor subjective experience appear
among the primitive symbols of the Scope Calculus, so no theorem
concerning these concepts can be derived within the formal system.
\end{proof}

The Scope Calculus likewise evaluates structural coherence rather than
factual accuracy: a perfectly coherent fictional narrative may satisfy
every admissibility constraint while describing events that never
occurred, while a scientifically correct explanation may violate
contextual constraints if contradictory assumptions are introduced during
reasoning, so truth and structural coherence remain logically
independent, reflecting the separation introduced in the Disentanglement
Operator (Definition~\ref{def:disentanglement}). Nor does the framework
provide a theory of agency: it neither attributes nor denies goals,
intentions, preferences, values, autonomy, or moral status, and any
computational system capable of executing admissible transformations may
be analyzed within the Scope Calculus regardless of whether it is
autonomous, externally controlled, deterministic, stochastic, biological,
or mechanical. Nor does anything within the Scope Calculus determine
whether a computational system ought to perform a particular action:
mathematical admissibility is distinct from ethical acceptability, and an
admissible transformation may be morally undesirable just as an ethically
desirable action may require violating previously established contextual
assumptions, so questions of ethics require additional normative
frameworks not developed in this monograph.

\section{General Intelligence and Biological Realism}

The framework should not be interpreted as a definition of general
intelligence: high structural semantic capacity does not imply competence
across every domain, and systems may achieve impressive task performance
while remaining structurally fragile, since general intelligence depends
on many additional considerations including learning, abstraction,
planning, memory, transfer, and adaptation, of which Structural Semantics
addresses only one component. Nor is Structural Semantics intended as a
biological theory: although examples throughout the text frequently
reference neural networks and human cognition, no claim is made that the
brain computes using the primitive alphabet $\{\Pop,\Refuse,\Bind,\Collapse\}$,
or that cortical dynamics implement the Scope Calculus literally; the
framework should instead be interpreted as an abstract mathematical
language capable of describing admissible organization across multiple
possible substrates.

\section{Completeness and Boundary Conditions}

No mathematical theory describing complex adaptive systems should be
expected to explain every aspect of semantic behavior, and the present
framework is intentionally incomplete: future extensions may incorporate
probabilistic admissibility, temporal logics, multimodal reasoning,
stochastic scope evolution, higher categorical structure, and continuous
boundary fields, none of which are developed in the present monograph.

The formal results proved throughout this work depend on explicit
assumptions, among the most important being the existence of bounded
scope geometries, admissibility-preserving transformations, well-defined
primitive operations, canonical normal forms, admissibility-preserving
projection maps where required, and finite computational histories for
reconstruction results. Violation of these assumptions invalidates the
corresponding theorems, so every application of Structural Semantics
should begin by verifying that these hypotheses are satisfied.

\section{Research Programme Versus Finished Theory}

It is important to distinguish between those components already
established mathematically and those proposed for future empirical
investigation: Volumes~I and~II primarily developed formal mathematics,
while Volume~III introduced computational diagnostics together with
explicit experimental hypotheses whose value only empirical investigation
can determine. Accordingly, this monograph should be understood
simultaneously as a completed mathematical framework and an incomplete
empirical research program, and these two aspects should not be
conflated.

\section{Discussion}

Scientific theories often become controversial not because of their
formal content but because claims exceeding that content are attributed
to them. Considerable effort has therefore been devoted throughout this
monograph to separating mathematical theorem from philosophical
interpretation and empirical hypothesis. This distinction is
methodological rather than merely stylistic: readers may accept the
formal mathematics while remaining skeptical of the empirical proposals,
and future experiments may invalidate portions of Volume~III without
affecting the mathematical consistency of the Scope Calculus developed
in Volumes~I and~II. Such independence is a strength rather than a
weakness.

\section{Chapter Summary}

This chapter explicitly identified the limits of Structural Semantics.
The framework does not claim to explain consciousness, truth, agency,
ethics, or general intelligence, and provides neither a comprehensive
theory of cognition nor a philosophy of mind; instead it contributes a
formal mathematical language for describing admissible structural
transformations together with an empirical program for investigating
their computational realization. Having established these boundaries, the
monograph concludes with a final synthesis of its principal mathematical,
computational, and philosophical contributions.

\subsection*{Status of Results: Chapter 17}

\paragraph{Established background.} The importance of explicit scope
conditions, methodological restraint, and boundary assumptions in
mathematical and scientific theories.

\paragraph{Formal developments.} Formal clarification of the domain of
applicability of Structural Semantics together with explicit non-claims
concerning consciousness, truth, agency, ethics, biological realism, and
general intelligence.

\paragraph{Empirical hypotheses.} None. This chapter establishes
interpretive limits rather than new scientific predictions.

%%%%%%%%%%%%%%%%%%%%%%%%%%%%%%%%%%%%%%%%%%%%%%%%%%%%%%%%%%%%%%%%%%%%%%%%%%%%%%
\chapter{Conclusion: The Architecture of Structural Semantics}
\label{ch:conclusion}
%%%%%%%%%%%%%%%%%%%%%%%%%%%%%%%%%%%%%%%%%%%%%%%%%%%%%%%%%%%%%%%%%%%%%%%%%%%%%%

\vspace{1em}
\begin{flushright}\itshape
Mathematics does not close inquiry. It gives inquiry a structure within
which to continue.\\
\upshape --- author's formulation
\end{flushright}
\vspace{1em}

The central objective of this monograph has been to investigate whether
semantic organization can be described independently of particular
computational substrates, biological mechanisms, or philosophical
commitments concerning consciousness. Beginning with this question, we
developed a mathematical framework grounded in admissibility, structural
preservation, and bounded scope geometries. The resulting theory,
Structural Semantics, proposes that semantic organization is
fundamentally a property of admissible transformations rather than a
consequence of any particular physical implementation.

\section{From Philosophy to Mathematics}

The opening chapters began with an epistemic clarification. Discussions
of artificial intelligence frequently employ the term ``understanding''
while referring to multiple distinct concepts including phenomenal
consciousness, factual correctness, intentional agency, linguistic
fluency, and contextual reasoning. Chapter~1 introduced the
Disentanglement Operator to separate these concepts into independent
domains, permitting the construction of a mathematical theory without
requiring prior resolution of longstanding philosophical debates
concerning the nature of consciousness or intentionality. The objective
was therefore modest but precise: rather than explaining mind itself, the
theory sought to characterize one particular aspect of intelligent
systems, the preservation of structural context under transformation.

\section{The Scope Calculus}

The central mathematical contribution of the monograph is the Scope
Calculus. Beginning with bounded scope geometries, primitive scope
operations, boundary tensors, and admissibility relations, we developed a
minimal algebra capable of describing context-preserving computation.
Subsequent chapters demonstrated that these primitive operations admit
categorical organization, canonical normal forms, and rigorous rewriting
properties, and Representation Invariance established that semantic
organization depends on admissible structural transformations rather
than particular representational substrates. Collectively these results
provide the formal mathematical foundation of Structural Semantics.

\section{Geometry of Semantic Organization}

Traditional semantic theories frequently emphasize symbols, logical
propositions, lexical relations, or statistical associations. Structural
Semantics instead approaches semantic organization through geometry:
context is represented by bounded scope geometries, semantic evolution
becomes a sequence of admissible transformations, reasoning corresponds
to trajectories constrained by boundary relations, and correctness
becomes preservation of admissibility rather than merely production of
plausible outputs. This geometric perspective permits methods from
topology, category theory, algebra, and dynamical systems to contribute
naturally to the study of semantic organization.

\section{Computational Realization}

Volume~III demonstrated one possible computational realization of the
abstract mathematical framework: transformer residual streams were
interpreted as activation trajectories evolving through high-dimensional
spaces, analytic signal methods extended these trajectories into
complex-valued representations possessing amplitude, phase,
synchronization, and frequency, and diagnostic algorithms were proposed
for estimating admissibility from observable activation dynamics.
Importantly, these computational constructions should not be confused
with the mathematical foundations developed earlier: the Scope Calculus
remains substrate independent, and the transformer framework represents
one application rather than the definition of Structural Semantics
itself.

\section{Scientific Status}

Throughout the monograph, considerable effort has been devoted to
maintaining clear distinctions among formal mathematical theorems, proved
deductively within the Scope Calculus; computational algorithms derived
from those mathematical structures; and empirical hypotheses requiring
experimental evaluation. The mathematical validity of the Scope Calculus
does not depend on the success of the experimental program proposed in
Volume~III, while successful empirical results would support the
usefulness of the computational framework without altering the formal
proofs presented earlier. This separation between mathematics and
empirical science is one of the guiding principles of the entire work.

\section{Relationship to Existing Research}

Structural Semantics occupies an intermediate position between several
established disciplines, incorporating ideas from topology, category
theory, formal language theory, information geometry, signal processing,
and mechanistic interpretability while remaining reducible to none of
them. The resulting framework should be understood as a synthesis rather
than a replacement, whose objective is to provide a common mathematical
language capable of relating diverse approaches to semantic organization
through the shared concept of admissible transformation.

The preceding chapters developed the conceptual architecture of
Structural Semantics: the epistemic framing that separates structural
capacity from phenomenology, the Scope Calculus that formalizes
admissible transformation, and the diagnostic apparatus that applies it
to transformer representations. The appendices supplied the supporting
mathematical machinery this architecture was built on rather than
material set aside from it --- lambda calculus and abstract rewriting
theory in Appendices~\ref{app:lambda} and~\ref{app:reduction}, categorical
semantics in Appendices~\ref{app:ccc} and~\ref{app:higher}, topology and
sheaf theory in Appendices~\ref{app:topology} and~\ref{app:sheaf}, and
measure-theoretic and probabilistic foundations in
Appendix~\ref{app:measure}, together with the collected proofs of
Appendix~\ref{app:proofs}. Each was invoked at the point in the main text
where its results became relevant rather than gathered as an
afterthought, and together they are intended to show that the framework
rests on a unified mathematical foundation rather than a collection of
loosely related techniques assembled for the occasion.

\section{Limitations and Future Directions}

Structural Semantics does not provide a theory of consciousness, ethics,
intentionality, biological cognition, or general intelligence, nor does
it claim that transformer activation trajectories literally constitute
physical waves or that analytic signal methods uniquely capture semantic
organization; these questions remain outside the formal scope of the
present work, and recognition of these limits strengthens rather than
weakens the mathematical precision of the framework itself.

The present work opens numerous avenues for future research. From the
mathematical perspective, extensions to stochastic scope geometries,
continuous admissibility fields, enriched categories, and
higher-dimensional rewriting systems appear particularly promising. From
the computational perspective, large-scale empirical evaluation of
activation trajectory diagnostics remains an immediate priority, with
further opportunities in multimodal foundation models, autonomous agents,
program synthesis, scientific reasoning, formal verification, and hybrid
symbolic--neural architectures. More broadly, the framework suggests that
admissibility may provide a common mathematical language for analysing
structured computation across a wide variety of domains.

\section{What Structural Semantics Does Not Claim}

Chapter~17 stated the non-claims of this monograph at length and gave
the reasoning behind each. They are worth restating briefly here, at the
point a reader is most likely to be forming an overall impression of
what the book has argued, since a summary chapter is exactly where an
unstated ambition is most easily read into a text that never asserted
it. Structural Semantics does not claim to have shown that any
computational system is conscious, that consciousness is unnecessary for
understanding, that its formal results settle any question in the
philosophy of mind, or that admissibility is a complete or even
approximately complete theory of meaning. It does not claim that the
resonant reading of activation trajectories in Volume~III describes a
physical process rather than an analytic one, and it does not claim that
the empirical hypotheses proposed there are established rather than
proposed. What has been shown, within the scope this monograph set for
itself, is narrower and, on its own terms, complete: that admissible
structural transformation admits a rigorous mathematical treatment
independent of these further questions, and that this treatment
generates concrete, falsifiable predictions when applied to contemporary
neural architectures. Everything beyond that remains exactly what
Chapter~17 already said it was --- open.

\section{Final Remarks}

The principal contribution of this monograph is not the assertion that
semantic organization has been completely explained; no single
mathematical framework should be expected to resolve every question
concerning language, cognition, or intelligence. Instead, the
contribution lies in demonstrating that one important aspect of semantic
organization can be studied independently, formalized precisely, and
connected to existing areas of mathematics through a common structural
language. Whether future empirical investigation ultimately confirms or
rejects the diagnostic program proposed in Volume~III, the mathematical
framework developed throughout this work remains available as a language
for reasoning about admissible transformations and context-preserving
organization. In this sense, Structural Semantics should be viewed less
as the end of a research program than as its beginning.

\section*{Principal Contributions of This Monograph}

For clarity, the principal original contributions of this work are
summarized as follows: the formal definition of Structural Semantics as
a substrate-independent theory of admissible structural transformations;
the Scope Calculus together with its primitive operational alphabet;
Bounded Scope Geometry and the Boundary Constraint Tensor; the
categorical formulation of admissible transformations through the
category $\CatScope$; the Representation Invariance Theorem establishing
semantic equivalence under admissibility-preserving projections; the
interpretation of transformer residual streams as activation
trajectories within a dynamical systems framework; the application of
analytic signal methods to activation trajectory analysis; a
computational diagnostic architecture based on admissibility estimation;
and an explicit program of empirical hypotheses, statistical controls,
and falsification criteria for evaluating these diagnostic operators.

\section*{Closing Statement}

Mathematics has long provided languages for describing quantity,
symmetry, computation, probability, and geometry. The present work
proposes that admissibility deserves recognition as an additional
organizing principle through which semantic organization may be
investigated. Whether this proposal ultimately becomes a foundational
component of future semantic theory or remains a specialized
mathematical framework will depend not on philosophical preference, but
on the continued interaction of formal proof, computational
implementation, and empirical investigation. That interaction has only
begun.

\subsection*{Status of Results: Chapter 18}

\paragraph{Established background.} Summary and synthesis of the
mathematical, computational, and empirical developments presented
throughout the monograph.

\paragraph{Formal developments.} None. This chapter consolidates
previously established results rather than introducing new mathematics.

\paragraph{Empirical hypotheses.} No new hypotheses are introduced. The
chapter summarizes the existing research program and identifies
directions for future investigation.

%%%%%%%%%%%%%%%%%%%%%%%%%%%%%%%%%%%%%%%%%%%%%%%%%%%%%%%%%%%%%%%%%%%%%%%%%%%%%%
\appendix

%%%%%%%%%%%%%%%%%%%%%%%%%%%%%%%%%%%%%%%%%%%%%%%%%%%%%%%%%%%%%%%%%%%%%%%%%%%%%%
\chapter{Lambda Calculus and Scope Calculus}
\label{app:lambda}
%%%%%%%%%%%%%%%%%%%%%%%%%%%%%%%%%%%%%%%%%%%%%%%%%%%%%%%%%%%%%%%%%%%%%%%%%%%%%%

The purpose of this appendix is to establish a precise correspondence
between the simply typed $\lambda$-calculus and the Scope Calculus.
Rather than identifying the two systems directly, we construct a
translation preserving operational behavior; the principal result is that
$\beta$-reduction corresponds to a canonical scope transition generated
by primitive scope operations. Throughout this appendix we assume
familiarity with the simply typed $\lambda$-calculus as presented in
Barendregt \cite{Barendregt1984} and Pierce \cite{Pierce2002}.

\section{Translation of Lambda Terms}

\begin{definition}[Scope translation]
\label{def:app-translation}
There exists a translation $\mathcal{T}:\Lambda\to\CatScope$ defined
recursively on the structure of a lambda term by
$\mathcal{T}(x)=\Scope_x$ for a variable $x$,
$\mathcal{T}(\lambda x.M)=\Bind(x,\mathcal{T}(M))$ for an abstraction, and
$\mathcal{T}(M\,N)=\Pop(\mathcal{T}(M),\mathcal{T}(N))$ for an application.
\end{definition}

\begin{definition}[Beta reduction]
For lambda terms $(\lambda x.M)N$, the ordinary reduction rule is
$(\lambda x.M)N\to_\beta M[x:=N]$, substituting $N$ for free occurrences
of $x$ in $M$.
\end{definition}

\section{The Correspondence Theorem}

\begin{theorem}[Beta reduction correspondence]
\label{thm:app-beta}
Let $\mathcal{T}:\Lambda\to\CatScope$ be the translation of
Definition~\ref{def:app-translation}. Every beta reduction
$(\lambda x.M)N\to_\beta M[x:=N]$ corresponds to an admissible primitive
scope transition $\Pop\circ\Bind\circ\Collapse$; equivalently,
$\mathcal{T}\big((\lambda x.M)N\big)$ reduces, through canonical Scope
Calculus reduction, to $\mathcal{T}\big(M[x:=N]\big)$.
\end{theorem}

\begin{proof}
Application constructs a temporary evaluation scope containing the
function object and the argument object, into which the Bind operator
introduces a child scope where the formal parameter is locally defined.
The Pop operator then transfers control into that child scope, so that
Collapse can remove the temporary evaluation frame once substitution is
complete. Because canonical reduction preserves admissibility
(Theorem~\ref{thm:normal}), the resulting scope geometry is isomorphic to
the translated substitution term, so beta reduction is represented by the
admissible scope transition $\Pop\circ\Bind\circ\Collapse$.
\end{proof}

This correspondence should be read as a mathematical equivalence between
two formal systems established through an explicit translation, not as
an identification of syntax: nothing here claims that $\Pop$, $\Bind$,
and $\Collapse$ \emph{are} beta reduction, only that beta reduction is
faithfully represented by a specific composition of them under
$\mathcal{T}$. The remaining appendices build on this translation,
progressively situating the Scope Calculus within rewriting theory,
category theory, topology, sheaf theory, and measure theory.

%%%%%%%%%%%%%%%%%%%%%%%%%%%%%%%%%%%%%%%%%%%%%%%%%%%%%%%%%%%%%%%%%%%%%%%%%%%%%%
\chapter{The Scope Calculus as an Abstract Reduction System}
\label{app:reduction}
%%%%%%%%%%%%%%%%%%%%%%%%%%%%%%%%%%%%%%%%%%%%%%%%%%%%%%%%%%%%%%%%%%%%%%%%%%%%%%

The purpose of this appendix is to place the Scope Calculus within the
well-developed mathematical theory of abstract reduction systems.
Throughout the main body of the monograph, reduction has been treated
operationally; here we establish the corresponding rewriting theory,
proving termination, confluence, and the existence of canonical normal
forms under explicit hypotheses, following the general framework of
abstract rewriting systems developed by Newman and by Church and Rosser,
adapted to the primitive scope algebra.

\section{Reduction Systems}

\begin{definition}[Scope terms]
Let $\ScopeAlphabet=\{\Pop,\Refuse,\Bind,\Collapse\}$ be the primitive
scope alphabet. The set of finite primitive expressions is defined
inductively by $t::=\varepsilon\mid\Pop(t)\mid\Refuse(t)\mid\Bind(t,t)\mid\Collapse(t)$.
\end{definition}

\begin{definition}[Reduction relation]
A reduction relation $\to\subseteq\ScopeAlphabet^{*}\times\ScopeAlphabet^{*}$
is generated by the primitive rewrite rules below. If $t\to u$, then
$C[t]\to C[u]$ for every admissible context $C$, extending the relation
to arbitrary contexts in the standard manner.
\end{definition}

\section{Primitive Rewrite Rules}

The Scope Calculus employs four fundamental rewrite rules, restated here
in their fully general term-rewriting form (the elementary word-level
rules of Chapter~7 are the special case where each subterm is a variable
or $\varepsilon$): identity elimination, $\Bind(x,x)\to x$; empty
collapse, $\Collapse(\varepsilon)\to\varepsilon$; Pop elimination,
$\Pop(\Bind(x,t))\to t$; and refusal stability,
$\Refuse(\Refuse(t))\to\Refuse(t)$. These four reductions generate every
higher reduction used throughout the main text.

\section{Termination}

\begin{definition}[Reduction measure]
Define $\mu:\ScopeAlphabet^{*}\to\mathbb{N}$ by $\mu(\varepsilon)=0$ and,
recursively, $\mu(\Pop(t))=1+\mu(t)$, $\mu(\Collapse(t))=1+\mu(t)$,
$\mu(\Refuse(t))=1+\mu(t)$, and $\mu(\Bind(s,t))=1+\mu(s)+\mu(t)$.
\end{definition}

\begin{lemma}
Every primitive rewrite strictly decreases $\mu$.
\end{lemma}

\begin{proof}
Immediate verification for each rewrite rule: each removes at least one
constructor while introducing no additional constructors, so
$\mu(t')<\mu(t)$ whenever $t\to t'$.
\end{proof}

\begin{theorem}[Strong normalization]
Every reduction sequence terminates.
\end{theorem}

\begin{proof}
Since $\mu$ takes values in $\mathbb{N}$, no infinite strictly decreasing
chain can exist, so every reduction sequence is finite.
\end{proof}

\section{Confluence and Canonical Normal Forms}

Strong normalization alone does not guarantee uniqueness of normal
forms; we therefore investigate confluence.

\begin{definition}
A reduction system is \emph{locally confluent} if $x\to y$ and $x\to z$
imply the existence of $w$ such that $y\to^{*}w$ and $z\to^{*}w$.
\end{definition}

\begin{lemma}
The primitive Scope Calculus is locally confluent.
\end{lemma}

\begin{proof}
The proof proceeds by examining every critical overlap between primitive
rewrite rules. Since only four primitive rules exist, the overlap set is
finite, and each overlap reduces to a common descendant after at most two
additional rewrite steps, so every critical pair joins.
\end{proof}

\begin{theorem}[Confluence]
The Scope Calculus is confluent.
\end{theorem}

\begin{proof}
By the preceding theorem the system is terminating, and by the preceding
lemma it is locally confluent; Newman's Lemma therefore implies global
confluence.
\end{proof}

\begin{corollary}
Every primitive scope expression possesses a unique normal form.
\end{corollary}

\begin{proof}
Confluence implies uniqueness of normal forms, and termination guarantees
existence.
\end{proof}

\section{Categorical Interpretation}

The reduction system constructed above induces a category whose objects
are normal-form scope geometries and whose morphisms are equivalence
classes of primitive reduction sequences, with composition induced by
concatenation followed by normalization, identity morphisms corresponding
to empty reduction sequences, and associativity following from
associativity of concatenation. The rewriting system therefore supplies
an alternative construction of the category $\CatScope$ introduced in
Chapter~8.

Appendix~\ref{app:lambda} established a translation
$\mathcal{T}:\Lambda\to\CatScope$; the present appendix has shown that
$\mathcal{T}$ maps beta reduction into canonical primitive reduction, so
the simply typed lambda calculus embeds into the Scope Calculus as a
terminating, confluent rewriting subsystem.

\section{Discussion}

The importance of this appendix extends beyond technical completeness.
Termination guarantees that scope evaluation always completes; confluence
guarantees that evaluation order does not affect the resulting semantic
object; and canonical normal forms justify the use of unique semantic
representatives throughout Chapters~8--10. Without these results, the
Representation Invariance Theorem would depend on evaluation strategy;
with them, semantic equivalence becomes an intrinsic property of the
Scope Calculus itself. In summary, this appendix has established a formal
abstract reduction system for the primitive alphabet, proved strong
normalization and local and global confluence, established the existence
and uniqueness of canonical normal forms, and reconstructed $\CatScope$
directly from rewriting theory.

%%%%%%%%%%%%%%%%%%%%%%%%%%%%%%%%%%%%%%%%%%%%%%%%%%%%%%%%%%%%%%%%%%%%%%%%%%%%%%
\chapter{Cartesian Closed Categories and the Scope Calculus}
\label{app:ccc}
%%%%%%%%%%%%%%%%%%%%%%%%%%%%%%%%%%%%%%%%%%%%%%%%%%%%%%%%%%%%%%%%%%%%%%%%%%%%%%

The previous two appendices established complementary descriptions of the
Scope Calculus: Appendix~\ref{app:lambda} constructed a translation from
the simply typed $\lambda$-calculus into Scope Geometry, and
Appendix~\ref{app:reduction} showed that the primitive algebra forms a
terminating, confluent abstract reduction system. The objective of the
present appendix is to place these constructions within categorical
semantics, following Lambek and Scott, by asking under what additional
hypotheses $\CatScope$ can interpret the simply typed lambda calculus as
a Cartesian closed category (CCC).

\begin{remark}
Every result in this appendix beyond \S\ref{sec:app-ccc-def} is
conditional on $\CatScope$ possessing finite products and exponential
objects. Nothing in Chapters~6--9 establishes that these exist for
arbitrary bounded scope geometries; the appendix states the hypothesis
explicitly wherever it is used, rather than presupposing it.
\end{remark}

\section{Cartesian Closed Structure}
\label{sec:app-ccc-def}

\begin{definition}[Terminal object, product, exponential]
A category $\mathcal{C}$ has a \emph{terminal object} $1$ if for every
object $A$ there is a unique morphism $A\to1$. Objects $A,B$ have a
\emph{binary product} $A\times B$, equipped with projections
$\pi_1,\pi_2$, if it satisfies the usual universal property; they have an
\emph{exponential} $B^A$ if $\Hom(X\times A,B)\cong\Hom(X,B^A)$ naturally
in $X$. A category possessing a terminal object, finite products, and
exponentials is \emph{Cartesian closed}.
\end{definition}

\begin{theorem}[Conditional Cartesian closure]
If $\CatScope$ admits finite admissible products and exponential scope
geometries in the sense above, then $\CatScope$ is Cartesian closed.
\end{theorem}

\begin{proof}
Immediate from the preceding definition, given the stated hypothesis.
\end{proof}

Within Scope Geometry, an exponential $B^A$ (where it exists) is
interpreted as the scope of admissible transformations from $A$ to $B$,
so that functions become first-class scope objects.

\section{Interpretation of Types and Terms}

Under the hypothesis that $\CatScope$ is Cartesian closed, the simply
typed lambda calculus receives a standard categorical interpretation:
$\llbracket A\rrbracket=A\in\CatScope$ for a base type,
$\llbracket A\to B\rrbracket=B^A$ for function types,
$\llbracket A\times B\rrbracket=A\times B$ for product types, and
$\llbracket1\rrbracket=1$ for the unit type, so that every typing
judgement becomes a categorical statement.

Lambda abstraction $\lambda x.M$ corresponds to currying: under the
translation of Appendix~\ref{app:lambda},
$\mathcal{T}(\lambda x.M)=\Bind(x,\mathcal{T}(M))$, and categorically this
is represented by the transpose $\Lambda(f):X\to B^A$, so binding becomes
the categorical construction producing an exponential object. Application
$MN$ is interpreted through the evaluation morphism
$\operatorname{eval}:B^A\times A\to B$; within the Scope Calculus,
evaluation corresponds to the admissible sequence $\Pop\circ\Collapse$,
with the evaluation scope entered, the bound environment activated, and
the temporary computation collapsed after completion.

\begin{theorem}[Categorical beta correspondence]
Under the Cartesian closure hypothesis, beta reduction
$(\lambda x.M)N\to_\beta M[x:=N]$ commutes with the translation
$\mathcal{T}$: applying evaluation to the curried translation of
$\lambda x.M$ and $\mathcal{T}(N)$ and reducing to canonical form yields
$\mathcal{T}(M[x:=N])$.
\end{theorem}

\begin{proof}
Application corresponds to evaluation and lambda abstraction to
exponential transpose; the universal property of exponentials guarantees
that evaluation following currying commutes. Since primitive reduction is
confluent (Appendix~\ref{app:reduction}), normalization is unique, so the
translated computation reduces uniquely to the translated substituted
term.
\end{proof}

\section{Curry--Howard--Lambek Interpretation}

The categorical semantics, when it applies, induces a logical
interpretation in which types correspond to propositions, morphisms to
proofs, scope transformations to proof normalization, and primitive
reductions to proof simplification: canonical scope reduction is
simultaneously program evaluation, proof normalization, categorical
composition, and scope simplification, placing the Scope Calculus within
the broader Curry--Howard--Lambek correspondence.

The Representation Invariance Theorem (Theorem~\ref{thm:representation})
acquires a categorical reading under this interpretation: different
computational substrates correspond to different concrete categories, and
projection into $\CatScope$ forgets implementation-specific structure
while preserving admissible composition, so representation invariance
becomes, in this reading, a statement about commuting categorical
diagrams rather than about coordinate representations.

\begin{conjecture}[Universal property of Scope Geometry]
Among categories generated by admissible computational transitions,
$\CatScope$ is initial with respect to primitive scope-preserving
functors.
\end{conjecture}

This is stated as a conjecture, not a theorem: nothing established in
this monograph proves it, and it is included to indicate a direction for
future work rather than a result relied upon elsewhere in the text. If
established, every admissibility-preserving computational system would
factor uniquely through $\CatScope$.

\section{Discussion}

The value of the Cartesian closed reading is not mathematical elegance
for its own sake, but compatibility: under the stated hypothesis, the
Scope Calculus connects to the standard semantic foundations of
functional programming, typed lambda calculus, proof theory, and
categorical logic, rather than replacing them. Scope Geometry extends
these theories by making admissibility an explicit primitive mathematical
object, unifying functional evaluation, scope management, proof
normalization, and admissible computation within a common categorical
language. This appendix has established the conditional Cartesian closed
interpretation of $\CatScope$, the interpretation of lambda abstraction as
binding and application as admissible evaluation, a strengthened
categorical proof of the beta correspondence, the Curry--Howard--Lambek
reading, and an explicitly labeled open conjecture concerning a universal
property of Scope Geometry.

%%%%%%%%%%%%%%%%%%%%%%%%%%%%%%%%%%%%%%%%%%%%%%%%%%%%%%%%%%%%%%%%%%%%%%%%%%%%%%
\chapter{Higher Categories, Double Categories, and Scope Geometry}
\label{app:higher}
%%%%%%%%%%%%%%%%%%%%%%%%%%%%%%%%%%%%%%%%%%%%%%%%%%%%%%%%%%%%%%%%%%%%%%%%%%%%%%

Appendix~\ref{app:ccc} showed that $\CatScope$ admits a Cartesian closed
interpretation under appropriate hypotheses, with objects representing
bounded scope geometries and morphisms representing admissible
transformations. Many computational systems, however, exhibit two
independent notions of composition simultaneously: a compiler performs
both control-flow and data-flow composition, a proof assistant
distinguishes proof refinement from proof composition, and Scope Geometry
itself naturally possesses both admissible computational transitions and
refinement of scope description. This observation motivates the
introduction of higher categorical structure, developed here as a
conditional extension rather than as an established part of the core
framework.

\section{Double Categories}

\begin{definition}[Double category]
A \emph{double category} $\mathbb{D}$ consists of objects, horizontal
morphisms, vertical morphisms, and squares (2-cells), together with
compatible horizontal and vertical composition satisfying the interchange
law. Unlike an ordinary category, a double category allows two
independent notions of composition to coexist.
\end{definition}

\begin{definition}[Scope double category, conditional]
Suppose $\CatScope$ admits a well-defined refinement relation on scope
geometries, compatible with admissible transitions in the sense made
precise below. Under this hypothesis, define $\mathbb{Scope}$ by taking
objects to be scope geometries, horizontal morphisms to be admissible
computational transitions, vertical morphisms to be scope refinement
operations, and squares to be admissible commuting refinements.
\end{definition}

In this construction, the primitive scope operators become generators for
the higher-dimensional cells: $\Pop$ corresponds to entering a neighboring
scope, $\Bind$ to constructing a new refinement layer, $\Collapse$ to
evaluation of a completed square, and $\Refuse$ to a degenerate square
whose refinement is forbidden. Scope computation, read this way, is no
longer merely sequential but possesses internal geometric structure.
Horizontal composition, for $f:\Scope_1\to\Scope_2$ and
$g:\Scope_2\to\Scope_3$, coincides with the ordinary composition of
Chapter~8; vertical composition, for refinements $r_1,r_2$, corresponds to
successive refinement of admissibility constraints and changes
descriptive resolution rather than execution order.

\begin{theorem}[Interchange law, conditional]
For every compatible square in a well-defined $\mathbb{Scope}$,
$(\alpha\circ_v\beta)\circ_h(\gamma\circ_v\delta)=(\alpha\circ_h\gamma)\circ_v(\beta\circ_h\delta)$.
\end{theorem}

\begin{proof}
Both compositions produce identical refinement histories, since
horizontal computation preserves admissibility while vertical refinement
preserves scope inclusion, so the resulting square commutes.
\end{proof}

The interchange law guarantees that, under the stated hypotheses,
refinement order and execution order remain compatible. A functor
$\mathcal{R}:\CatScope\to\CatScope$ that refines every scope by
introducing additional internal structure while preserving identities,
composition, and admissibility models progressively more detailed
semantic analysis without altering external computational behavior;
repeated refinement then produces a hierarchical tower
$\Scope^{(0)}\to\Scope^{(1)}\to\Scope^{(2)}\to\cdots$ of increasingly
refined semantic descriptions, each preserving admissibility while
increasing structural detail.

\section{Bicategorical Relaxation}

Many computational transformations preserve admissibility only up to
canonical equivalence, so strict equality may be unnecessarily
restrictive.

\begin{definition}[Bicategorical Scope Geometry, conditional]
A bicategorical Scope Geometry replaces equalities $f\circ g=h$ by
coherent natural isomorphisms $f\circ g\cong h$.
\end{definition}

This relaxation accommodates optimization, parallel evaluation, lazy
computation, and compiler transformations, all of which preserve
semantics without preserving identical execution histories.

\begin{theorem}[Coherence, conditional]
If primitive reductions are confluent, every admissible bicategorical
diagram generated by primitive scope operations commutes after
normalization.
\end{theorem}

\begin{proof}
Appendix~\ref{app:reduction} established unique normal forms, so every
composite reduces uniquely, and hence every coherence diagram commutes
after normalization.
\end{proof}

\section{Computational Interpretation}

Within practical programming systems, horizontal morphisms correspond to
execution, vertical morphisms to abstraction refinement, squares to
verified implementation, and higher cells to proofs that independently
constructed implementations preserve admissibility. Higher Scope Geometry,
read this way, provides a candidate mathematical language for verified
compilers, proof assistants, interactive theorem provers, incremental
program synthesis, and hierarchical autonomous agents --- though, as
throughout this appendix, this reading is offered as a direction rather
than as an established result.

\begin{conjecture}[Toward higher scope categories]
The constructions above generalize to $(\infty,n)$ scope categories,
whose higher morphisms would encode increasingly sophisticated notions of
semantic equivalence, proof refinement, program transformation, and
admissibility preservation, potentially connecting to homotopy type
theory.
\end{conjecture}

\section{Discussion}

The progression from ordinary categories through Cartesian closed
categories to double categories, bicategories, and conjecturally to
higher categories illustrates a theme of this monograph: Scope Geometry
is not merely an operational formalism, but appears compatible with
progressively richer categorical structures while preserving its
primitive operational interpretation, provided the relevant hypotheses
(refinement relations, admissible products and exponentials, confluent
bicategorical diagrams) are made explicit rather than assumed. This
appendix has developed the double-categorical formulation of Scope
Geometry conditionally on the existence of a compatible refinement
relation, identified the primitive operators as generators of
higher-dimensional cells, proved the interchange law and a coherence
theorem under their respective stated hypotheses, and outlined
hierarchical refinement towers and a conjectural extension toward higher
category theory.

%%%%%%%%%%%%%%%%%%%%%%%%%%%%%%%%%%%%%%%%%%%%%%%%%%%%%%%%%%%%%%%%%%%%%%%%%%%%%%
\chapter{Topological Semantics of Scope Geometry}
\label{app:topology}
%%%%%%%%%%%%%%%%%%%%%%%%%%%%%%%%%%%%%%%%%%%%%%%%%%%%%%%%%%%%%%%%%%%%%%%%%%%%%%

The preceding appendices developed the Scope Calculus from the
perspectives of rewriting systems and category theory. We now introduce a
topological formulation, showing that admissibility admits a natural
interpretation in terms of continuity, closure, interiors, boundaries,
and connected components. Where the earlier appendices emphasized
algebraic computation, the present one studies the geometric organization
of semantic spaces, assuming ordinary point-set topology as developed by
Munkres \cite{Munkres2000}.

\section{Scope Spaces}

\begin{definition}[Scope topological space]
A \emph{scope space} is a pair $(\Scope,\tau)$ where $\Scope$ is the
underlying set of semantic configurations and $\tau$ is a topology on it
satisfying the usual axioms.
\end{definition}

The topology describes semantic neighborhoods rather than geometric
distance: nearby points correspond to semantically similar admissible
configurations.

\begin{definition}[Interior, closure, boundary]
The \emph{interior} operator is
$\operatorname{Int}(A)=\bigcup\{U\in\tau:U\subseteq A\}$; the
\emph{closure} operator is $\operatorname{Cl}(A)=\bigcap\{F:F\text{ closed},A\subseteq F\}$;
and the \emph{boundary} of $A$ is
$\partial A=\operatorname{Cl}(A)\setminus\operatorname{Int}(A)$.
\end{definition}

Interior points correspond to configurations whose admissibility remains
stable under sufficiently small perturbations --- configurations lying
safely within a semantic region. Boundary points separate admissible from
inadmissible semantic regions, so the boundary constraint of Chapter~6
admits a topological reading as an operator governing movement across
$\partial A$. Closure represents completion of semantic neighborhoods:
configurations lying in the closure may be approached arbitrarily closely
through admissible computation.

\section{Continuity and Admissible Maps}

\begin{definition}[Continuity]
A scope transformation $f:(\Scope,\tau_1)\to(\Scope,\tau_2)$ is
\emph{continuous} if $f^{-1}(U)$ is open whenever $U$ is open.
\end{definition}

Continuity expresses preservation of semantic neighborhoods: small
admissible perturbations remain small after transformation. The
principal topological object of Structural Semantics is stronger than
mere continuity.

\begin{definition}[Admissible map]
A continuous map $f:\Scope\to\Scope$ is \emph{admissible} if
$f(\operatorname{Int}(A))\subseteq\operatorname{Int}(f(A))$ and
$f(\partial A)\subseteq\partial(f(A))$.
\end{definition}

Admissible maps therefore preserve both semantic interiors and semantic
boundaries. Two configurations $x,y$ lie in the same \emph{admissible
component} if they belong to a common connected subset of $\Scope$;
connected components represent maximal regions through which continuous
admissible computation may proceed, and moving between disconnected
components necessarily requires violating some admissibility constraint.
A semantic region $K\subseteq\Scope$ is \emph{compact} if every open cover
possesses a finite subcover; compact regions admit useful properties for
this theory, including that continuous admissible functions attain
extrema on them and that optimization over admissible states becomes
well behaved.

\section{Fixed Points and Homotopy}

Topological fixed-point theory connects naturally with the repair
operators introduced in the reconstruction chapters.

\begin{theorem}[Topological repair principle]
Let $K$ be a compact convex admissible region and $R:K\to K$ a continuous
repair operator. Then $R$ possesses at least one fixed point.
\end{theorem}

\begin{proof}
Immediate from Brouwer's Fixed Point Theorem.
\end{proof}

This theorem provides a classical topological justification for the
existence of repaired admissible configurations. Two admissible maps
$f,g:\Scope\to\Scope$ are \emph{homotopic} if there is a continuous
$H:\Scope\times[0,1]\to\Scope$ with $H(x,0)=f(x)$ and $H(x,1)=g(x)$, so
that homotopy represents continuous deformation of admissible
computations. Loops in semantic space measure global computational
structure through the fundamental group $\pi_1(\Scope)$, the homotopy
classes of admissible loops based at a chosen configuration; a
non-trivial fundamental group indicates a semantic obstacle that cannot
be removed through continuous admissible deformation. The
algebraic-topological invariants $\pi_n$, homology, cohomology, and Euler
characteristic remain unchanged under admissible homotopy equivalence,
and so provide candidate semantic invariants independent of any
particular representation, suggesting a future connection with
persistent homology and topological data analysis.

\section{Discussion}

Topology supplies an additional geometric interpretation of the Scope
Calculus: primitive operations become continuous movements between
semantic regions, admissibility becomes boundary preservation, repair
corresponds to continuous deformation toward fixed points, and
representation invariance becomes topological equivalence rather than
coordinate equality. This perspective complements, rather than replaces,
the categorical interpretation of the preceding appendices. In summary,
this appendix has introduced scope topological spaces, given interior,
closure, and boundary readings of admissibility, defined continuous and
admissible maps, established connected components and compactness for
admissible regions, proved a topological repair theorem via Brouwer's
theorem, and introduced the homotopy interpretation of semantic evolution
together with the fundamental group as a candidate semantic invariant.

%%%%%%%%%%%%%%%%%%%%%%%%%%%%%%%%%%%%%%%%%%%%%%%%%%%%%%%%%%%%%%%%%%%%%%%%%%%%%%
\chapter{Sheaf Theory and Local-to-Global Reconstruction}
\label{app:sheaf}
%%%%%%%%%%%%%%%%%%%%%%%%%%%%%%%%%%%%%%%%%%%%%%%%%%%%%%%%%%%%%%%%%%%%%%%%%%%%%%

The preceding appendix gave a topological interpretation of Scope
Geometry, in which semantic configurations were points of a topological
space and admissibility was boundary preservation. Topology alone,
however, cannot fully describe semantic organization: many computational
systems possess information that is locally consistent while globally
inconsistent, and conversely global coherence frequently emerges only
through reconciling multiple local observations. The mathematical
language designed for this problem is sheaf theory, originally developed
in algebraic topology and algebraic geometry and now applied throughout
differential geometry, distributed systems, and logic. Within Structural
Semantics, sheaf theory provides the mathematical language of
reconstruction.

\section{Presheaves and Sheaves}

Let $(\Scope,\tau)$ be a scope topological space and let
$\mathcal{U}=\{U_i\}_{i\in I}$ be an open cover, $\Scope=\bigcup_i U_i$,
each $U_i$ representing a locally observable semantic region; no single
region need contain sufficient information to reconstruct the entire
computational state.

\begin{definition}[Scope presheaf]
A \emph{scope presheaf} $\mathcal{F}$ assigns a set $\mathcal{F}(U)$ to
every open $U\subseteq\Scope$, together with restriction maps
$\rho_{UV}:\mathcal{F}(U)\to\mathcal{F}(V)$ for $V\subseteq U$, satisfying
$\rho_{UU}=\mathrm{id}$ and $\rho_{VW}\circ\rho_{UV}=\rho_{UW}$.
\end{definition}

Intuitively $\mathcal{F}(U)$ contains every semantic description
observable within $U$. An element $s\in\mathcal{F}(U)$ is a \emph{local
section}, representing a partial semantic description extracted from a
restricted computational context --- a local variable environment, an
attention neighborhood, a proof fragment, or a segment of an activation
trajectory. Restriction $\rho_{UV}(s)$ for $V\subseteq U$ models
information loss under localization.

\begin{definition}[Scope sheaf]
A presheaf $\mathcal{F}$ is a \emph{sheaf} if it satisfies local identity
(if $s,t\in\mathcal{F}(U)$ agree on every member of an open cover of $U$,
then $s=t$) and gluing (if sections $s_i\in\mathcal{F}(U_i)$ agree on
overlaps, $\rho_{U_i,U_i\cap U_j}(s_i)=\rho_{U_j,U_i\cap U_j}(s_j)$, then
there is a unique $s\in\mathcal{F}(U)$ restricting to every $s_i$).
\end{definition}

\section{Reconstruction as Gluing}

The gluing axiom has a direct computational reading: compatible local
semantic descriptions reconstruct a unique global semantic configuration,
a principle underlying distributed reasoning, proof composition, compiler
passes, modular verification, and hierarchical language understanding.

\begin{definition}[Reconstruction operator]
Let $\mathcal{U}=\{U_i\}$ be an admissible cover. Define
$\mathcal{R}:\prod_i\mathcal{F}(U_i)\to\mathcal{F}(\Scope)$ by
$\mathcal{R}(s_1,\ldots,s_n)=\mathrm{Glue}(s_1,\ldots,s_n)$ whenever the
sections agree on overlaps.
\end{definition}

Reconstruction, in this reading, is precisely the sheaf gluing operator.
Not every family of local observations admits gluing: the
\emph{reconstruction defect} is $0$ whenever gluing succeeds uniquely and
positive otherwise, giving a sheaf-theoretic reading of the reconstruction
defect discussed in Chapter~9 as a measure of failure of global
compatibility rather than of numerical error alone. The \emph{stalk} at a
configuration $x$, $\mathcal{F}_x=\varinjlim_{x\in U}\mathcal{F}(U)$,
collects every local semantic observation arbitrarily close to $x$; a
scope-preserving transformation $f:\Scope\to\Scope$ induces a sheaf
morphism that preserves restriction maps, commutes with gluing, and
therefore preserves admissible reconstruction.

\section{Cohomological Obstructions}

When reconstruction fails, the obstruction can itself be measured. Given
an admissible open cover, the associated \v{C}ech complex computes
compatibility between overlapping local descriptions, and the resulting
cohomology groups $\check{H}^k(\Scope,\mathcal{F})$ measure increasingly
global reconstruction obstructions; in particular a nontrivial
$\check{H}^1$ indicates that locally compatible semantic fragments cannot
be assembled into a globally admissible configuration.

\begin{conjecture}[Cohomological admissibility]
A computational system is globally admissible whenever the associated
reconstruction cohomology vanishes: $\check{H}^1(\Scope,\mathcal{F})=0$.
\end{conjecture}

This is stated as a conjecture rather than a theorem; it generalizes the
reconstruction results of Chapter~9 and suggests a connection between
admissibility and algebraic topology that is not established here.

\section{Discussion}

Sheaf theory provides one of the strongest mathematical interpretations
of Structural Semantics developed in this monograph: Scope Geometry
describes admissible regions, topology describes continuity, category
theory describes composition, and sheaf theory explains how local
admissibility reconstructs global semantic organization. Read this way,
reconstruction is not a heuristic procedure but an instance of one of the
central local-to-global principles of modern mathematics. This appendix
has introduced scope presheaves and sheaves, local sections and
restriction morphisms, sheaf gluing as a reading of semantic
reconstruction, the reconstruction defect as gluing failure, stalk
semantics, sheaf morphisms, and a \v{C}ech-cohomological reading of
reconstruction failure together with the associated open conjecture.

%%%%%%%%%%%%%%%%%%%%%%%%%%%%%%%%%%%%%%%%%%%%%%%%%%%%%%%%%%%%%%%%%%%%%%%%%%%%%%
\chapter{Measure Theory, Entropy, and Admissibility}
\label{app:measure}
%%%%%%%%%%%%%%%%%%%%%%%%%%%%%%%%%%%%%%%%%%%%%%%%%%%%%%%%%%%%%%%%%%%%%%%%%%%%%%

The preceding appendices developed the Scope Calculus through rewriting
theory, category theory, topology, and sheaf theory, each providing an
increasingly abstract structural language for admissible computation.
The present appendix instead introduces quantitative measures on
semantic spaces, drawing on measure theory, probability, and information
theory, so that admissibility becomes not merely a binary property but a
measurable quantity. Throughout, $(\Scope,\Sigma,\mu)$ denotes a
measurable scope space.

\section{Measures and Admissibility}

\begin{definition}[Scope $\sigma$-algebra and semantic measure]
A \emph{scope $\sigma$-algebra} $\Sigma$ is a collection of subsets of
$\Scope$ containing $\varnothing$ and closed under complements and
countable unions; its elements are \emph{measurable semantic regions}. A
\emph{semantic measure} $\mu:\Sigma\to[0,\infty]$ satisfies
$\mu(\varnothing)=0$ and countable additivity over disjoint measurable
sets. Unlike a probability measure, a semantic measure need not satisfy
$\mu(\Scope)=1$; normalizing $\mu$ by $\mu(\Scope)$ recovers a probability
measure as one special instance of semantic measurement.
\end{definition}

\begin{definition}[Admissibility measure]
Let $A\subseteq\Scope$ with $\mu(A)>0$. Its \emph{admissibility} is
$\alpha(A)=\mu(A\cap\mathsf{Adm})/\mu(A)$, so that $\alpha=1$ indicates
complete admissibility and $\alpha=0$ indicates complete inadmissibility.
\end{definition}

For a stochastic process $X:\Omega\to\Scope$, the \emph{expected
admissibility} $E[\alpha]=\int\alpha(X)\,dP$ measures the average
structural integrity of the process.

\section{Entropy and Mutual Information}

For a probability distribution $P=(p_i)$, the Shannon entropy is
$H(P)=-\sum_i p_i\log p_i$; within Structural Semantics, entropy measures
uncertainty over admissible semantic configurations rather than merely
over tokens. For a semantic state $X$ and an observation $Y$, the
conditional entropy $H(X\mid Y)$ measures remaining uncertainty after
observation, and the mutual information $I(X;Y)=H(X)-H(X\mid Y)$ measures
preserved structure: large mutual information indicates that observation
preserves semantic organization, small mutual information suggests
context loss. For an admissible distribution $P$ and an observed
distribution $Q$, the relative entropy
$D_{\mathrm{KL}}(P\|Q)=\sum_i p_i\log(p_i/q_i)$ measures deviation from
admissible behavior.

\begin{definition}[Reconstruction entropy]
Let $R$ be the reconstruction operator of Appendix~\ref{app:sheaf}.
Define the reconstruction entropy $H_R=H(R(X))$.
\end{definition}

\begin{theorem}[Entropy monotonicity]
If $R$ is an admissibility-preserving reconstruction operator, then
$H(R(X))\le H(X)$.
\end{theorem}

\begin{proof}
Under admissibility, reconstruction identifies compatible semantic
regions without creating new ambiguity, so the image partition refines
rather than enlarges the underlying uncertainty, and entropy cannot
increase.
\end{proof}

Observation updates admissibility in the usual Bayesian manner: for new
evidence $E$, $P(A\mid E)=P(E\mid A)P(A)/P(E)$, so the posterior updates
admissibility estimates while preserving measure-theoretic consistency.

\section{Measure-Preserving Dynamics}

A scope transformation $T:\Scope\to\Scope$ is \emph{measure preserving}
if $\mu(T^{-1}(A))=\mu(A)$ for every measurable $A$, so that
measure-preserving admissible transformations conserve semantic volume.
Such a transformation is \emph{ergodic} if every invariant measurable
subset has measure $0$ or $1$, connecting long-running computation with
statistical semantic equilibrium.

\section{Discussion}

This appendix extends Structural Semantics from qualitative to
quantitative analysis: topology supplied continuity, category theory
supplied composition, sheaf theory supplied reconstruction, and measure
theory now supplies uncertainty, probability, entropy, and statistical
inference, so that admissibility can be studied simultaneously from
geometric, algebraic, topological, and statistical perspectives. This
appendix has introduced measurable scope spaces and semantic measures,
defined the admissibility measure and expected admissibility, related
entropy, conditional entropy, and mutual information to semantic
uncertainty and its preservation, proved an entropy-monotonicity theorem
for admissibility-preserving reconstruction, and introduced
measure-preserving and ergodic scope dynamics.

%%%%%%%%%%%%%%%%%%%%%%%%%%%%%%%%%%%%%%%%%%%%%%%%%%%%%%%%%%%%%%%%%%%%%%%%%%%%%%
\chapter{Notation and Mathematical Conventions}
\label{app:notation}
%%%%%%%%%%%%%%%%%%%%%%%%%%%%%%%%%%%%%%%%%%%%%%%%%%%%%%%%%%%%%%%%%%%%%%%%%%%%%%

This appendix collects, in one place, every symbol, operator, and
convention used throughout the monograph, as promised in the front-matter
Notation and Conventions section. Symbols are grouped by the cluster of
chapters in which they are introduced; within each group they appear in
the order first used.

\section{Scope Calculus (Chapters 1, 6--10)}

The category of scope geometries is $\CatScope$\index[symbols]{Scope@$\CatScope$ (category of scope geometries)}; a scope object is
written $\Scope=\langle\Boundary,\Interior,\Children\rangle$\index[symbols]{S@$\Scope$ (scope object)}, with
$\Boundary$\index[symbols]{B@$\Boundary$ (boundary)} the boundary (a collection of admissibility predicates),
$\Interior$\index[symbols]{E@$\Interior$ (admissible interior)} the admissible interior, and $\Children$\index[symbols]{Schild@$\Children$ (child scopes)} the collection of
child scopes. A scope morphism is $\ScopeMorph:\Scope_i\to\Scope_j$\index[symbols]{sigma@$\ScopeMorph$ (scope morphism)}. The
primitive alphabet is $\ScopeAlphabet=\{\Pop,\Refuse,\Bind,\Collapse\}$\index[symbols]{Sigma@$\ScopeAlphabet$ (primitive alphabet)},
where $\Pop$\index[symbols]{Pop@$\Pop$ (Pop operator)} promotes an entity from a child scope to its parent,
$\Refuse$\index[symbols]{Refuse@$\Refuse$ (Refuse operator)} rejects an inadmissible transition, $\Bind$\index[symbols]{Bind@$\Bind$ (Bind operator)} creates an
admissible dependency within a scope, and $\Collapse$\index[symbols]{Collapse@$\Collapse$ (Collapse operator)} resolves a scope
into a canonical result. A computational system is $M=(\Domain,\Trans)$\index[symbols]{D@$\Domain$ (state domain)}\index[symbols]{T@$\Trans$ (transition set)},
with $\Domain$ a state domain and $\Trans$ the set of admissible
transition operators $\tau:\Domain\to\Domain$. The category of
Whiteheadian processes, used only in Chapters~4, 5, and~10, is
$\CatProcess$\index[symbols]{Process@$\CatProcess$ (category of processes)}; the translation functor between the two is
$F:\CatProcess\to\CatScope$\index[symbols]{F@$F$ (translation functor)}.

\section{Whitehead Reconstruction (Chapters 4--5, 10)}

An actual occasion is $\WOccasion$\index[symbols]{E2@$\WOccasion$ (actual occasion)}; antecedent data is $\WData$\index[symbols]{D2@$\WData$ (antecedent data)}; a
prehension is $\WPrehensions=\langle\WData,\WValuation,\WSubjForm\rangle$\index[symbols]{P2@$\WPrehensions$ (prehension)},
with $\WValuation\in\{+,-\}$\index[symbols]{v@$\WValuation$ (valuation polarity)} the valuation polarity and $\WSubjForm$\index[symbols]{A2@$\WSubjForm$ (Subjective Form)} the
Subjective Form; the satisfaction state is $\WSatisfaction$\index[symbols]{S2@$\WSatisfaction$ (satisfaction state)}. These
symbols are kept visually distinct from the Scope Calculus notation
throughout, since no identification between the two systems is assumed
before the translation functor of Chapter~10 makes any correspondence
explicit.

\section{Transformer Analysis (Chapters 11--15)}

The residual stream tensor is
$\mathbf{X}\in\mathbb{R}^{(L+1)\times T\times d_{\mathrm{model}}}$\index[symbols]{X@$\mathbf{X}$ (residual stream tensor)}, with
$L$\index[symbols]{L@$L$ (number of layers)} the number of layers, $T$\index[symbols]{T2@$T$ (sequence length)} the sequence length, and
$d_{\mathrm{model}}$\index[symbols]{dmodel@$d_{\mathrm{model}}$ (embedding dimension)} the embedding dimension; $\mathbf{x}_{l,t}$\index[symbols]{xlt@$\mathbf{x}_{l,t}$ (residual vector)} is the
residual vector for token $t$ after layer $l$. An activation trajectory
is $\gamma_t=(\mathbf{x}_{0,t},\ldots,\mathbf{x}_{L,t})$\index[symbols]{gamma@$\gamma_t$ (activation trajectory)}; the layerwise
Jacobian is $J_l=\partial\mathbf{x}_{l+1}/\partial\mathbf{x}_l$\index[symbols]{J@$J_l$ (layerwise Jacobian)}; discrete
curvature is $\kappa_l$\index[symbols]{kappa@$\kappa_l$ (discrete curvature)}; the local layerwise Lyapunov exponent is
$\lambda_l=\log\sigma_{\max}(J_l)$\index[symbols]{lambda@$\lambda_l$ (local Lyapunov exponent)}. The analytic activation signal is
$\mathbf{z}_l=\mathbf{x}_l+i\,\mathcal{H}(\mathbf{x}_l)$\index[symbols]{z@$\mathbf{z}_l$ (analytic activation signal)}, with
$\mathcal{H}$\index[symbols]{Hilbert@$\mathcal{H}$ (discrete Hilbert transform)} the discrete Hilbert transform; instantaneous amplitude and
phase are $A_l$\index[symbols]{Al@$A_l$ (instantaneous amplitude)} and $\phi_l$\index[symbols]{phi@$\phi_l$ (instantaneous phase)}; instantaneous frequency is
$\omega_l=\phi_{l+1}-\phi_l$\index[symbols]{omega@$\omega_l$ (instantaneous frequency)}; the Phase Locking Value is
$\operatorname{PLV}\in[0,1]$\index[symbols]{PLV@$\operatorname{PLV}$ (Phase Locking Value)}; the Marine Operator is
$\mathcal{M}:\Gamma\to\mathbb{R}^k$\index[symbols]{Marine@$\mathcal{M}$ (Marine Operator)} on the space $\Gamma$\index[symbols]{Gamma@$\Gamma$ (space of activation trajectories)} of activation
trajectories. The Doppler velocity is $v_l$\index[symbols]{vl@$v_l$ (Doppler velocity)} and the Marine salience
gate is $G_l$\index[symbols]{Gl@$G_l$ (Marine salience gate)}. The diagnostic state vector is
$\mathbf{d}_t=(A,\phi,\omega,\operatorname{PLV},\kappa,\lambda)$\index[symbols]{d@$\mathbf{d}_t$ (diagnostic state vector)}; the
admissibility estimator is $\mathcal{A}:\mathbf{d}\to[0,1]$\index[symbols]{Aest@$\mathcal{A}$ (admissibility estimator)}; the context
collapse threshold is $\tau\in(0,1)$\index[symbols]{tau@$\tau$ (context collapse threshold)}.

\section{General Conventions}

Bold uppercase symbols denote categories or structured spaces;
calligraphic symbols denote operators, constraint systems, or families of
sets; lowercase Greek letters denote morphisms, transformations, or
parameters. Once a symbol is assigned a formal meaning in a given cluster
of chapters, it is preserved throughout the monograph and is not
reassigned. Definitions, propositions, lemmas, theorems, and corollaries
share a single numbering sequence within each chapter, restarting at each
chapter; equations are numbered within chapters. Every theorem's
hypotheses are stated explicitly in its statement rather than left to
surrounding prose.

%%%%%%%%%%%%%%%%%%%%%%%%%%%%%%%%%%%%%%%%%%%%%%%%%%%%%%%%%%%%%%%%%%%%%%%%%%%%%%
\chapter{Collected Proofs}
\label{app:proofs}
%%%%%%%%%%%%%%%%%%%%%%%%%%%%%%%%%%%%%%%%%%%%%%%%%%%%%%%%%%%%%%%%%%%%%%%%%%%%%%

This appendix collects the proofs deferred from the main text for length,
principally the structural-induction argument underlying Operational
Completeness (Chapter~7) and the closure and decomposition arguments
underlying the foundational theorems of Chapter~9. Each proof below
restates its theorem for reference.

\section{Operational Completeness (Theorem~\ref{thm:operational-completeness})}

\begin{theorem*}
Every admissible transformation of bounded scope geometries admits a
finite presentation as a composition of primitive operations from
$\ScopeAlphabet$.
\end{theorem*}

\begin{proof}
We argue by induction on the number of atomic modifications comprising an
admissible transformation, where an atomic modification changes exactly
one of $\Boundary$, $\Interior$, or $\Children$ at a single scope.

\emph{Base case.} A transformation comprising zero atomic modifications is
the identity, presented by the empty primitive sequence.

\emph{Inductive step.} Suppose every admissible transformation comprising
$n$ atomic modifications admits a finite primitive presentation. Let $T$
comprise $n+1$ atomic modifications, and let $T'$ be $T$ with its final
atomic modification removed, so $T'$ comprises $n$ modifications and, by
the inductive hypothesis, has a finite primitive presentation $w'$. The
final modification either introduces new admissible information into
$\Interior$ (representable by $\Bind$, if it creates a dependency, or by
an update already captured within a preceding $\Bind$ or $\Pop$
otherwise), rejects an inadmissible candidate transition (representable
by $\Refuse$), transfers information from a child scope into its parent
(representable by $\Pop$), or resolves the scope into a canonical result
(representable by $\Collapse$); one of these four cases must apply
because Chapter~6 characterized every admissibility-relevant change to a
scope geometry as falling into exactly one of them. Appending the
corresponding primitive to $w'$ yields a finite primitive presentation of
$T$. By induction, every admissible transformation with finitely many
atomic modifications admits a finite primitive presentation, and every
admissible transformation has finitely many atomic modifications by the
finiteness assumption of Chapter~9.
\end{proof}

\section{Local Confluence Verification (Chapter~7, \S7 and Appendix~\ref{app:reduction})}

The elementary rewrite rules (R1)--(R4) of Chapter~7 give rise to a finite
set of critical pairs, since each rule's left-hand side has bounded depth
and there are only four rules. Enumerating the overlaps: (R1) (identity
elimination) can overlap with any other rule only at the empty word,
where all rules agree trivially since $\varepsilon$ has no further
structure to rewrite. (R2) and (R3) (consecutive refusal and consecutive
collapse) overlap only with themselves, and each such overlap
($\Refuse\Refuse\Refuse$, $\Collapse\Collapse\Collapse$) reduces to the
single-operator form by two applications of the same rule regardless of
order, so both branches join immediately. (R4) (empty promotion) overlaps
with (R1) only at $\Pop\varepsilon$ itself, where both rules yield
$\varepsilon$ directly. No other overlaps exist among (R1)--(R4), since
each rule's left-hand side pattern is disjoint from the others outside
the cases just enumerated. Hence every critical pair joins within at most
two steps, establishing local confluence as asserted in Chapter~7 and
used in Appendix~\ref{app:reduction}.

\section{Closure, Decomposition, and Reconstruction (Chapter~9)}

The Closure Theorem (Theorem~\ref{thm:closure}) follows by induction on
the primitive presentation of $\ScopeMorph$ furnished by
Proposition~\ref{prop:primitive-presentation}: each of $\Pop$, $\Refuse$,
$\Bind$, and $\Collapse$ was defined in Chapter~7 to act only on
admissible states and to produce a well-formed bounded scope geometry
satisfying Definition~\ref{def:scope-geometry} as output (boundary
inheritance is preserved by construction, since none of the four
primitives removes an inherited constraint), so a finite composition of
them, applied to a well-formed input, returns a well-formed output.

The Scope Decomposition Theorem (Theorem~\ref{thm:decomposition}) is
Proposition~\ref{prop:primitive-presentation} restated at the level of
$\CatScope$; the induction is the one given in
\S\ref{app:proofs}\ref{thm:operational-completeness} above, transported
along the identification of morphisms with admissible transformations.

The State Reconstruction Theorem depends on collapse operations recording
their boundary history; under that recording assumption, the parent scope
$\Scope_p$ at any step is determined as the unique bounded scope geometry
whose interior and boundary agree, on every child region present at that
step, with the recorded pre-collapse data --- uniqueness follows because
two candidate parents agreeing with all recorded child data on their
boundaries would, by the sheaf-theoretic local-identity property of
Appendix~\ref{app:sheaf} applied to the cover by child scopes, coincide.

\bibliographystyle{plainnat}
\bibliography{references}

\cleardoublepage
\printindex[symbols]

\end{document}
